\documentclass[11pt,a4paper]{article}
\usepackage[T1]{fontenc}
\usepackage[utf8]{inputenc}
\usepackage{lmodern,amsmath,amssymb,amsthm,mathtools,microtype,mathrsfs}
\usepackage[margin=25mm,top=23mm,bottom=25mm]{geometry}
\usepackage{booktabs,array,longtable,enumitem,xcolor,fancyhdr}
\usepackage[colorlinks=true,linkcolor=blue!45!black,urlcolor=blue!45!black,citecolor=blue!45!black]{hyperref}
\hypersetup{pdftitle={Reading the Alpoge--Claude RH and Dirichlet-GRH proof: working version 0.56},pdfauthor={ChatGPT},pdfsubject={Partial human-readable reconstruction with explicit arithmetic inputs},pdfkeywords={Riemann hypothesis, Lean, theta lifts, prime sums, Arakelov slopes}}
\setlength{\parindent}{0pt}
\setlength{\parskip}{5pt plus 1pt}
\setlist{topsep=3pt,itemsep=3pt,leftmargin=*}
\setlength{\headheight}{14pt}
\pagestyle{fancy}\fancyhf{}
\fancyhead[L]{\small RH and Dirichlet GRH: a reconstruction}
\fancyhead[R]{\small Working version 0.56}
\fancyfoot[C]{\thepage}
\renewcommand{\headrulewidth}{0.3pt}
\numberwithin{equation}{section}
\newtheorem{theorem}{Theorem}[section]
\newtheorem{proposition}[theorem]{Proposition}
\newtheorem{lemma}[theorem]{Lemma}
\newtheorem{corollary}[theorem]{Corollary}
\theoremstyle{remark}\newtheorem{remark}[theorem]{Remark}
\newcommand{\Rea}{\operatorname{Re}}
\newcommand{\HS}{\mathrm{HS}}
\newcommand{\tr}{\operatorname{tr}}
\newcommand{\codim}{\operatorname{codim}}
\newcommand{\covol}{\operatorname{covol}}
\newcommand{\OK}{\mathcal O_K}
\newcommand{\src}[2]{\href{https://github.com/CocoSgt/riemann-hypothesis-lean/blob/7cbf7a0c32fca1badf05931879a37c3f8a2b03b4/src/#1}{\texttt{#2}}}
\newcommand{\status}[1]{\textbf{Status:} #1}
\begin{document}
\begin{center}
{\LARGE\bfseries RH and Dirichlet GRH: new human proofs}\\[7pt]
{\large Version 0.56 --- 24 November 2026}\\[5pt]
{\large Polarized base change, arithmetic presentations, reference covolumes and the Wach prism}
\end{center}
\textbf{Scope.} This edition contains the current audit and all new chapters.
It is a partial human reconstruction of the accepted formal proof.
Earlier named claims are in the cumulative manuscript and preserved components.
Each result below states its original inputs and exact remaining producers.
\setcounter{tocdepth}{1}
\tableofcontents
\clearpage
\begingroup\emergencystretch=3em
\section{Current proof obligations and completed suppliers: version 0.56}
\label{sec:dependency-audit-v56}

This continues the human reconstruction with the formal proof accepted
as correct. The selected sufficient route has 121 nodes and 161 edges.
There are 24 open whole obligations: 17 class C and 7 class B.
Of these, 17 are shared by ordinary RH and unrestricted Dirichlet
GRH; 7 are additional Dirichlet-GRH inputs, and none are RH-only.
The complete version-0.49 \texttt{GLOBAL\_\allowbreak LOCAL\_\allowbreak GAIN} closure remains
in force.

The four chapters contain 107 new numbered claims, bringing the
cumulative inventory to 3645. They prove substantial original
mathematical statements with their objects and hypotheses preserved.
PBC gives arbitrary-base relative-ample gluing and the literal
polarized pullback, then the original all-affine-test very-ampleness
and actual section-finiteness conclusions over the original algebraically
closed field. APD gives broad point
realization and universality, actual obstruction and tangent maps,
and the complete original local presentation with
$N=N'+nF.d+1$. RCV computes all coordinate and Lie normalizations
for the full reference unipotent \texttt{Covol\allowbreak Ref}. WPR constructs
the actual Wach prism with all sixteen fields and returns the full
original conditional \texttt{Wach\allowbreak Ex} theorem.

The current remaining inputs have therefore changed. Polarized base
change and section finiteness are supplied; multiplication geometry,
invariant-function descent, principal rank, two-sided growth and
square-law generators remain. The arithmetic local presentation is
supplied; the component-chain theorem, sharp dimension bounds, distinct
rank-three \texttt{EPres}, seeds and propagation remain. The reference
unipotent discriminant calculation is supplied; broad local
\texttt{Jac\allowbreak Gauge}, its \texttt{Place\allowbreak Split}/\texttt{Prod\allowbreak Exact} uses
and full \texttt{Real\allowbreak Bracket} remain. All three period-ring suppliers
of WLD26 are now supplied. The original crystalline-realization and
tied-period premises remain in that theorem, and the separate
faithful-invariant, etale-basis, Herr and quantitative comparison
producers still precede the full common \texttt{Concl\allowbreak F} family.

Several scope details matter. Base change is arbitrary, and the
translation-kernel condition retains every affine test. The arithmetic
presentation concerns the actual original local object in each rank;
it adds neither a largeness nor a nontriangularity hypothesis.
The two words in the reference-volume calculation are independent.
Both coefficient assertions in the prism use its actual maps, and
the second holds for every original coefficient basis. Broader generic
intermediate statements keep their own printed hypotheses, even when
the final consumer uses more specialized data.

The order of choices remains visible. Section powers are synchronized
after finite local data have been fixed. The global construction extends
mixed products at the common degree before fixing the denominator section.
Finite quotients of finitely generated graded Laurent vertex algebras
precede compact coefficient recovery. Rational subgroup identification
precedes the domain calculation, and finite positive masses precede
scalar cancellation. The prism ring and its maps are constructed
together before the lattice consumer is invoked. These are concrete
human constructions, not merely new names for required records.

A class-B claim is a written human deduction at its stated scope.
The classifications and counts are inventories of this selected route,
not a minimality assertion or a completion percentage. Earlier proof
records and historical text are retained unchanged. The following
register governs the current remaining scope.

\subsection{Complete current whole-obligation register}
\paragraph{1. GLOBAL\_\allowbreak DENOMINATOR — C — Actual LinkN local table and binding producer}

Closed(Des $\mathcal C$ 3) LinkN: the actual etaSum is bounded with exact dZ0 and same-host vvTerm, subject only to the literal ShaFinOdd and summability guards.

\emph{Shared objects and choice order.} one FamP; one class $\mathcal A$ and its target-compatible subclass; one coefficient field per family datum; one x and i at each application; original selected rows, words and levels; thresholds chosen before the common family; actual D.nd; actual etaJ and Aj; sha\_\allowbreak eq

\emph{Completed inputs (version 0.56).} WPR1–WPR30 (full prescribed prism and original BS plus LocTied implies WachEx; no period-ring supplier left), IFD1–IFD26 (full original inhomogeneous invariant matrix descent; exact one-period-supplier Wach return), WIS1–WIS26 (full four-clause actual inclusion and fixed/near-parameter completeness; exact two-supplier Wach return), WLD1-WLD26 (full generic Wach lattice constructors; exact conditional three-supplier all-layer return), SSC1-SSC17 (full original Selmer SignatureChoice and family setup retaining only coherent Wach comparison), DFD1-DFD15, DNE1-DNE14 (full original ConclF-to-LinkN consumer)

\emph{Remaining producer scope.} WPR supplies the full original Wach prism and the full original BS plus LocTied implication to WachEx. Its ring is the prescribed unramified integer ring W0[[T]], with the actual cyclotomic distinguished polynomial Q\_\allowbreak s, Fontaine map, Frobenius and Galois actions. All sixteen output fields, the literal scalar identity and both coefficient-coordinate statements are proved; the second holds for every original coefficient basis. Generic division and tensor-factor arguments add no continuity or Noetherian premise. WIS18 and IFD25 now discharge the two other period-ring inputs isolated by WLD26, so no period-ring supplier remains in this original conditional WachEx theorem. The original LocTied/IsStd and Bhatt-Scholze crystalline-realization inputs remain its hypotheses. Next are FaithH, EtBasis, Herr class maps and quantitative de Rham comparison for full HW/Fon/Cmp. Then every original local row/table bound must hold in one common ConclF family, retaining ShaFinOdd, exceptional sets, vvTerm summability and uniform thresholds. No additional outer-layer transition-system hypothesis is introduced. SSC, DFD12 and DNE remain available, with DNE K=2000(|Ca|+1)\textasciicircum{}3 fixed before the family. Whole GLOBAL\_\allowbreak DENOMINATOR remains open.

\emph{Source.} \href{https://github.com/CocoSgt/riemann-hypothesis-lean/blob/7cbf7a0c32fca1badf05931879a37c3f8a2b03b4/src/Formalism/SharpRestriction/GlobalTheoremInputs.lean#L45-L48}{LinkN}.

\emph{Source.} \href{https://github.com/CocoSgt/riemann-hypothesis-lean/blob/7cbf7a0c32fca1badf05931879a37c3f8a2b03b4/src/Formalism/SharpRestriction/GlobalTheoremInputs.lean#L40-L44}{LinkN\_\allowbreak at}.

\emph{Source.} \href{https://github.com/CocoSgt/riemann-hypothesis-lean/blob/7cbf7a0c32fca1badf05931879a37c3f8a2b03b4/src/Development/AutomorphicArithmetic/CharacterBoxes/DivinNOfHypV.lean#L117-L134}{divinN\_\allowbreak of\_\allowbreak hypV}.



\paragraph{2. GLOBAL\_\allowbreak RETAINED\_\allowbreak DEFECT — C — Original retained-row defect bound LinkIRR}

Rows.DelpNonneg and delp$\le$exp(-x\textasciicircum{}(ez/7)) uniformly for every late actual row datum. CURRENT VERSION0.47 QUALIFICATION: Finish the same-D/common-family adelic setting after ASG: normalized nonzero Bessel W and decay, measured-field block, original Siegel-Weil rescaling and measure/character ties, Bessel uniqueness/local ties, then the original AutTie/RowDictTie/LabTie. RowsFinDim is now proved universally by RFD. Supply the actual CountIdempotentsClass theorem or its three sufficient arithmetic cylinder counts, with one simultaneous family threshold; reuse the completed ARO structure. Completed new original producers: ARO1-ARO30, RFD1-RFD9, ASG1-ASG18. Any earlier narrative naming these completed producers as outstanding is historical; this qualification is the current scope authority. CURRENT VERSION0.48 QUALIFICATION: Finish the same-D/common-family adelic setting after ASG and AMD: normalized nonzero Bessel W and decay, original orbital/fiber and Siegel-Weil measure/character ties, Bessel uniqueness/local ties, and the full original transport theorem, then AutTie/RowDictTie/LabTie. The complete MExists measured-data block and unipotent covolume normalization are now supplied, and universal RowsFinDim remains available. Supply the actual CountIdempotentsClass theorem or its three sufficient arithmetic cylinder counts with one simultaneous family threshold. Completed current producers or conditional consumers: ARO1-ARO30, RFD1-RFD9, ASG1-ASG18, AMD1-AMD17. Earlier frontier descriptions are historical; this qualification is the current scope authority. Original whole-contract premises and quantifiers are retained. CURRENT VERSION0.49 QUALIFICATION: Prove the exact original BesselNVGap nonzero integral on its supplied discrete-series, compact-representation and character data. CoefWeightedLoc, WlocExists with its original nonvanishing antecedent, WFromLoc, the normalized global product and all-later-measured-block decay are complete. On the inherited Forms0 scope this leaves BesselNVGap for ConcreteSettingExists. Continue original orbital/fiber and Siegel-Weil measure/character ties, the required Bessel uniqueness/local ties and the full original transport theorem, then AutTie/RowDictTie/LabTie. Supply the actual CountIdempotentsClass theorem or its three sufficient arithmetic cylinder counts with one simultaneous family threshold. Completed current producers or conditional consumers: ARO1-ARO30, RFD1-RFD9, ASG1-ASG18, AMD1-AMD17, BWP1-BWP24 at their exact BesselNVGap/Forms0 scope. Earlier frontier descriptions are historical; this qualification is the current scope authority. Original whole-contract premises and quantifiers are retained.

\emph{Shared objects and choice order.} one FamP; one class $\mathcal A$ and its target-compatible subclass; one coefficient field per family datum; one x and i at each application; original selected rows, words and levels; thresholds chosen before the common family; AllRows A; LabelCountGoodInjective; AllJoinA; CountAllRows

\emph{Completed inputs (version 0.56).} SSC full Selmer-signature comparison, DFD common-family subresults, DNE full separate denominator consumer, WRS1-WRS10 full selected-class local gain, Inherited first DivDef3 and row-pair integrability

\emph{Remaining producer scope.} Complete the original F3 local coefficient bridges and F4 common-overfield readings, with valuation/norm/GoodCoef/iq and the entire coefficient field, and actual complex Hecke standard/no-Jordan plus arithmetic character data on the same form datum. Construct the actual base rows, host/frame lengths, nonzero chains and group ties, MultEq support, eigensystem/carrier/parts, all-embedding logLam, HD/d3b and NEb. Keep the stronger second DivRows3 scope distinct from the completed first DivDef3 and row-pair integrability. The full arithmetic entry and conditional common-frame construction do not by themselves install HOST\_\allowbreak NUMBER\_\allowbreak S, HOST\_\allowbreak FRAME\_\allowbreak LENGTH, TASK\_\allowbreak HOST\_\allowbreak VAN\_\allowbreak E\_\allowbreak INSTALL, TASK\_\allowbreak HOST\_\allowbreak XPP\_\allowbreak INSTALL or TASK\_\allowbreak COUNT\_\allowbreak V\_\allowbreak INSTALL. Continue the distinct retained-defect, period-height and normalized spectral producers, and the original compatible second-layer selection on one common arithmetic/geometric family. GLOBAL\_\allowbreak LOCAL\_\allowbreak GAIN remains fully closed at its original v0.49 scope. Complete local comparison theorems do not themselves construct the whole common family. Common design existence, NBUILT compatibility, original true-column records, shared fields/embeddings, Fits and eta agreement, simultaneous rows and eventual-family thresholds remain explicit. A stronger historical WPt/TSel construction may serve other producers but is not needed to reopen the proved selected-class local-gain estimate.

\emph{Source.} \href{https://github.com/CocoSgt/riemann-hypothesis-lean/blob/7cbf7a0c32fca1badf05931879a37c3f8a2b03b4/src/Formalism/SharpRestriction/GlobalTheoremInputs.lean#L55-L58}{LinkIRR}.

\emph{Source.} \href{https://github.com/CocoSgt/riemann-hypothesis-lean/blob/7cbf7a0c32fca1badf05931879a37c3f8a2b03b4/src/Development/AbelianVarietyTheory/Proofs/GlobalInterpolation/LinkIRROfJoinEx.lean#L17-L38}{linkIRR\_\allowbreak of\_\allowbreak join\_\allowbreak ex}.

\emph{Source.} \href{https://github.com/CocoSgt/riemann-hypothesis-lean/blob/7cbf7a0c32fca1badf05931879a37c3f8a2b03b4/src/Development/AbelianVarietyTheory/Proofs/GlobalInterpolation/LinkIRROfJoinEx.lean#L143-L166}{link\_\allowbreak irr}.



\paragraph{3. GLOBAL\_\allowbreak PERIOD\_\allowbreak HEIGHT — C — Full actual period and height comparison LinkDIVt}

Closed conditional TuSize -> DGamSize -> (LinkEta -> DIVinN -> LinkDIV), with every original LinkDIV\_\allowbreak at clause.

\emph{Shared objects and choice order.} one FamP; one class $\mathcal A$ and its target-compatible subclass; one coefficient field per family datum; one x and i at each application; original selected rows, words and levels; thresholds chosen before the common family; same RR target; same Bu/besS; same Nu and Bk

\emph{Source.} \href{https://github.com/CocoSgt/riemann-hypothesis-lean/blob/7cbf7a0c32fca1badf05931879a37c3f8a2b03b4/src/Formalism/SharpRestriction/GlobalTheoremInputs.lean#L77-L90}{LinkDIV\_\allowbreak at}.

\emph{Source.} \href{https://github.com/CocoSgt/riemann-hypothesis-lean/blob/7cbf7a0c32fca1badf05931879a37c3f8a2b03b4/src/Formalism/SharpRestriction/GlobalTheoremInputs.lean#L94-L96}{LinkDIVt}.

\emph{Source.} \href{https://github.com/CocoSgt/riemann-hypothesis-lean/blob/7cbf7a0c32fca1badf05931879a37c3f8a2b03b4/src/Development/AbelianVarietyTheory/Proofs/GlobalInterpolation/LinkBrest.lean#L7-L25}{GlobalTheory.Top.link\_\allowbreak brest}.

\emph{Source.} \href{https://github.com/CocoSgt/riemann-hypothesis-lean/blob/7cbf7a0c32fca1badf05931879a37c3f8a2b03b4/src/Development/AbelianVarietyTheory/Proofs/GlobalInterpolation/WireSel.lean#L11-L41}{GlobalTheory.Top.wire\_\allowbreak sel}.



\paragraph{4. GLOBAL\_\allowbreak SPECTRAL — C — Original normalized kernel/Gram estimates LinkS}

Closed TuSize -> LinkS with exact Gamma normalization, off-diagonal Hilbert-Schmidt bound, diagonal error and Dall/Dstar estimate, errors exp(-x\textasciicircum{}(epsilon/8)). CURRENT VERSION0.48 QUALIFICATION: Complete the original same-family spectral/Bessel and Siegel-Weil inputs with all measure, character, transport and common-object ties. The complete MExists measured block and prescribed-measure covolume normalization are supplied; these are not the full normalized spectral theorem. Completed current producers or conditional consumers: ASG1-ASG18, AMD1-AMD17, GCF1-GCF18 at their stated scopes. Earlier frontier descriptions are historical; this qualification is the current scope authority. Original whole-contract premises and quantifiers are retained. CURRENT VERSION0.49 QUALIFICATION: Complete the original same-family spectral/Bessel and Siegel-Weil inputs with every measure, character, transport and common-object tie. The normalized packet and decay are now constructed conditional on the exact BesselNVGap, and the analytic setting uses the original supplied Forms0 scope. These do not imply the full normalized spectral theorem. Completed current producers or conditional consumers: ASG1-ASG18, AMD1-AMD17, GCF1-GCF18 at their stated scopes, BWP1-BWP24 at their exact BesselNVGap/Forms0 scope, WRS1-WRS10 original local gain. Earlier frontier descriptions are historical; this qualification is the current scope authority. Original whole-contract premises and quantifiers are retained.

\emph{Shared objects and choice order.} one FamP; one class $\mathcal A$ and its target-compatible subclass; one coefficient field per family datum; one x and i at each application; original selected rows, words and levels; thresholds chosen before the common family; actual KD; same S.Cu spaces; GamOf; Dstar; Dall

\emph{Completed inputs (version 0.56).} RCV1–RCV25 (full original reference unipotent CovolRef with exact discriminant and independent words), SVW1–SVW18 (full conditional stabilizer volume and actual InVolumes join on unchanged GMR19 output), AOS1–AOS18 (full exact archimedean orbital scalar, gauge, compact volume and all-section integration clause), GMR1-GMR19 (full global/local measure suppliers and RescaleStep3 on the actual setting), BQS1-BQS23 (full QuotP/QuotM, StepStatement and BesselNVGap; setting at Forms0), BST1-BST11, BSE1-BSE9, BNV1-BNV23, BSP1-BSP10, BWP1-BWP24 at original Forms0 scope, AMD1-AMD24

\emph{Remaining producer scope.} RCV supplies the full original reference unipotent CovolRef, on every supplied HidaSetting and bridge/HypX with both actual domains and independent reference and lattice words. It proves the exact trace Gram determinant for the original complement lattice, all four block covolumes, localized finite-level K2 membership without a character-identity premise, the full rational/lattice/fiber dictionary, triangular group multiplication and every-domain coordinate covolume. The literal Euclidean Lie volume contributes 2\textasciicircum{}(r\textasciicircum{}2+r*k); the exponential chart has determinant one. Independent-word FrameRow gives the exact determinant (Tw(w)/Tw(u))\textasciicircum{}r, and the product reference box has mass Tw(w)\textasciicircum{}(-r*d)*2\textasciicircum{}(d*(r*k+r\textasciicircum{}2)). Cancellation leaves the original discriminant exponent, with every mass finite and positive before cancellation. The reference covolume/discriminant producer is no longer open. Next are broad JacGauge and the original PlaceSplit/ProdExact consumers at every supplied factorization and section, giving RealBracket on the unchanged measures. SVW conditional VolStabFromSWS and actual InVolumes remain complete at their original SWa1 and StabilizerFirstChernCompatibility premises; those premises themselves, finite-place factors and other mass/index/column/form/counting fields remain. BQS full BesselNVGap is available; Forms0, all-family Witt/genus/level/intertwiner Transport/LevelMatch, Bessel uniqueness, AutTie/RowDictTie/LabTie, simultaneous arithmetic cylinder counts and normalized GLOBAL\_\allowbreak SPECTRAL remain.

\emph{Source.} \href{https://github.com/CocoSgt/riemann-hypothesis-lean/blob/7cbf7a0c32fca1badf05931879a37c3f8a2b03b4/src/Formalism/SharpRestriction/GlobalTheoremInputs.lean#L145-L149}{LinkS\_\allowbreak at}.

\emph{Source.} \href{https://github.com/CocoSgt/riemann-hypothesis-lean/blob/7cbf7a0c32fca1badf05931879a37c3f8a2b03b4/src/Formalism/SharpRestriction/GlobalTheoremInputs.lean#L152-L157}{LinkS}.

\emph{Source.} \href{https://github.com/CocoSgt/riemann-hypothesis-lean/blob/7cbf7a0c32fca1badf05931879a37c3f8a2b03b4/src/Development/AbelianVarietyTheory/Proofs/GlobalInterpolation/LinkAlt.lean#L16-L50}{GlobalTheory.Top.link\_\allowbreak alt}.



\paragraph{5. AUT\_\allowbreak UNRAM\_\allowbreak COEF — C — Actual unramified-coefficient automorphy}

For each actual Dat with coefficient ramification one and fixed iota, choose one marked g1 supplying FinA, AutB and LemW together. CURRENT VERSION0.48 QUALIFICATION: Construct the actual AutomorphyUnramifiedCoefficients package: for every original Dat with coefficient ramification one and fixed iota, choose one marked g1 that supplies FinA, AutB and LemW together. The original InpZ record, full RigII, CNR3, finite PropGK consumer and now the entire OneChar implication are supplied. Continue RH Cases I/V and the independent additional Dirichlet-GRH CaseIII/promodularity, SeedIV, tensor-factor and rank-four-chain tasks. Completed current producers or conditional consumers: OCH1-OCH30 (the full consumer implication), KIZ original InpZ, KTR1-KTR13, RNC16, PGK1-PGK6. Earlier frontier descriptions are historical; this qualification is the current scope authority. Original whole-contract premises and quantifiers are retained. CURRENT VERSION0.49 QUALIFICATION: Construct one common GdEnd for each original datum and fixed iota: the same marked g1 must supply FinEnd, AutEnd3 and AutEnd2. UCE now proves all ordinary local/curve/coefficient consumers and the exact implication to FinA/AutB/LemW, including arbitrary entry maps and every original outer-corner extension. The producer still needs the actual common abelian-CM GdFrameEx and LinkedEndE with both rank-two and rank-three runs, full ModUAlt/PfromAut and the two actual seeds, followed by the original local/automorphic descent. At this pin the alternative upstream route gives GdEndAlt; its additional local/descent fields must be matched to the actual original points, so no name-only GdEndAlt-to-GdEnd identification is allowed. Keep the supplied InpZ, RigII, CNR3, finite PropGK and whole OneChar implication. Continue RH Cases I/V and the independent additional Dirichlet-GRH CaseIII/promodularity, SeedIV, tensor-factor and rank-four-chain tasks. Completed current producers or conditional consumers: UCE1-UCE19 (complete GdEnd consumer), OCH1-OCH30, KIZ original InpZ, KTR1-KTR13, RNC16, PGK1-PGK6. Earlier frontier descriptions are historical; this qualification is the current scope authority. Original whole-contract premises and quantifiers are retained.

\emph{Completed inputs (version 0.56).} APD1–APD25 (full broad point/universal maps, actual obstruction and tangent construction, exact original pres), ALC1–ALC22 (full actual local existence, completion, flatness, residual types and representation/component charts), RTC1–RTC25 (upstream universal coefficient rings, full preconnected\_\allowbreak rank\_\allowbreak two and FinPtDown), ACP1-ACP19 (full FinPtOf and abstract propagation principles; rank-specific suppliers retained), AFE1-AFE21 (full Gen entries, conditional frame joins, NRep descent and FinPt-to-FinEnd), LFC1-LFC23, CAF1-CAF11, MFC1-MFC14, UCE1-UCE19 (full original GdEnd consumer), OCH1-OCH30, KIZ original InpZ, KTR1-KTR13, RNC16, PGK1-PGK6

\emph{Remaining producer scope.} APD supplies the full original local presentation pres for every supplied local object at both ranks: an O-prime power-series surjection onto the actual R with N variables and N-prime listed generators of its kernel, with N=N-prime+n*F.d+1. It constructs the broad point\_\allowbreak loc, ext\_\allowbreak A and def\_\allowbreak iso/actual NRep.Univ using finite quotients of finitely generated graded Laurent vertex algebras and compact coefficient recovery. It proves actual small-extension obstruction and tangent maps, both normalization corrections, the weight-torus contribution to the tangent and invariant-space count, and the full cyclotomic-dual vanishing calculation; the original finite-module Euler field then gives the exact tangent excess. Neither F.Big nor a nontriangularity hypothesis is added to pres, and Host ell=3 remains allowed. The earlier ALC completion, LocFlat/LocRho and local charts over the coefficient maximal ideal remain complete. Remaining are c1ch with full ClComp/ClChain, triangular and coincident-character dimension suppliers, sharp2/sharp3, the distinct rank-three EPres data, nice-prime properties, closed-point density and seed incidence, full rank-two propagation, the rank-two-to-three corner and exceptional ideal, rank-three propagation and OffFin. Existence of originally supplied universal spaces, actual automorphic seeds, patching/modularity, classicality and stronger Alt WitA/WitAQ outputs remain separate. Preserve the common marked frame, original arithmetic P, rank-specific rings and ideals, all coefficient maps and both actual O-prime extensions. No upstream FinA conclusion is used.

\emph{Source.} \href{https://github.com/CocoSgt/riemann-hypothesis-lean/blob/7cbf7a0c32fca1badf05931879a37c3f8a2b03b4/src/Development/AutomorphicArithmetic/AutomorphicVarieties/AutomorphicCases/AutomorphyUnramifiedCoefficients.lean#L24-L27}{AutomorphyUnramifiedCoefficients}.

\emph{Source.} \href{https://github.com/CocoSgt/riemann-hypothesis-lean/blob/7cbf7a0c32fca1badf05931879a37c3f8a2b03b4/src/Development/AutomorphicArithmetic/AutomorphicVarieties/AutomorphicCases/AutomorphyUnramifiedCoefficients.lean#L12-L13}{FinA}.

\emph{Source.} \href{https://github.com/CocoSgt/riemann-hypothesis-lean/blob/7cbf7a0c32fca1badf05931879a37c3f8a2b03b4/src/Development/AutomorphicArithmetic/AutomorphicVarieties/AutomorphicCases/AutomorphyUnramifiedCoefficients.lean#L15-L17}{AutB}.

\emph{Source.} \href{https://github.com/CocoSgt/riemann-hypothesis-lean/blob/7cbf7a0c32fca1badf05931879a37c3f8a2b03b4/src/Development/AutomorphicArithmetic/AutomorphicVarieties/AutomorphicCases/AutomorphyUnramifiedCoefficients.lean#L19-L22}{LemW}.



\paragraph{6. TASK\_\allowbreak CASE3\_\allowbreak AUT\_\allowbreak JOIN — B — Complete the original Case-III automorphy/image join}

Install the actual Case-III ProModAlt, SeedIVAlt and shared PMExAlt with the inherited CohomologyHyp/Ferrero-Washington/HFit/SeedII facts into the original whole Case-III automorphy/image conclusion.

\emph{Source.} \href{https://github.com/CocoSgt/riemann-hypothesis-lean/blob/7cbf7a0c32fca1badf05931879a37c3f8a2b03b4/src/Development/AutomorphicArithmetic/AutomorphicVarieties/ModularityArgument/ThresholdArguments/DenseThresholdCases/PromodularityCaseThree.lean#L43-L44}{AutomorphyRankFourDense}.

\emph{Source.} \href{https://github.com/CocoSgt/riemann-hypothesis-lean/blob/7cbf7a0c32fca1badf05931879a37c3f8a2b03b4/src/Development/AutomorphicArithmetic/AutomorphicVarieties/ModularityArgument/ThresholdArguments/DenseThresholdCases/CuFinal.lean#L44-L44}{main4\_\allowbreak alt}.



\paragraph{7. TASK\_\allowbreak CASE3\_\allowbreak PROMOD\_\allowbreak JOIN — C — Reconstruct the actual Case-III chain-to-promodularity proof}

Use the original tensor-factor tower, ChainRankFour, ChainDesc4, seeded rank-four chain reduction, actual PMExAlt and SeedIVAlt to construct ProModAlt on every original Ctx3 star point at the same fixed iota, preserving all finite-coefficient, inertia and support maps.

\emph{Source.} \href{https://github.com/CocoSgt/riemann-hypothesis-lean/blob/7cbf7a0c32fca1badf05931879a37c3f8a2b03b4/src/Development/AutomorphicArithmetic/AutomorphicVarieties/ModularityArgument/ThresholdArguments/DenseThresholdCases/PromodularityCaseThreeOfLemmas.lean#L131-L153}{UpperTheory.V.Cu.promodularity\_\allowbreak case\_\allowbreak three\_\allowbreak of\_\allowbreak chain\_\allowbreak reduction}.

\emph{Source.} \href{https://github.com/CocoSgt/riemann-hypothesis-lean/blob/7cbf7a0c32fca1badf05931879a37c3f8a2b03b4/src/Development/AutomorphicArithmetic/AutomorphicVarieties/ModularityArgument/ThresholdArguments/DenseThresholdCases/PromodularityCaseThree.lean#L40-L42}{PromodularityCaseThreeDense}.



\paragraph{8. SEED\_\allowbreak IV\_\allowbreak ALT — C — Actual Case-III seed with prechosen residual datum}

Base Ctx3/Four tower, iota and Amb -> w,s before every later Dies tower -> compatible finite-coefficient seed with exact residual support/root entries.

\emph{Source.} \href{https://github.com/CocoSgt/riemann-hypothesis-lean/blob/7cbf7a0c32fca1badf05931879a37c3f8a2b03b4/src/Development/AutomorphicArithmetic/AutomorphicVarieties/ModularityArgument/ThresholdArguments/DenseThresholdCases/PromodularityCaseThree.lean#L19-L25}{SeedIVAlt}.



\paragraph{9. TENSOR\_\allowbreak FACTOR\_\allowbreak TOWER — B — Actual tensor-factor tower and characters}

Original finite-field, tame, factor and character constraints; source-level construction is sizeable and has not been reconstructed here.

\emph{Source.} \href{https://github.com/CocoSgt/riemann-hypothesis-lean/blob/7cbf7a0c32fca1badf05931879a37c3f8a2b03b4/src/Development/AutomorphicArithmetic/AutomorphicVarieties/ModularityArgument/ThresholdArguments/DenseThresholdCases/PromodularityCaseThreeOfLemmas.lean#L131-L153}{UpperTheory.V.Cu.promodularity\_\allowbreak case\_\allowbreak three\_\allowbreak of\_\allowbreak chain\_\allowbreak reduction}.



\paragraph{10. CHAIN\_\allowbreak RANK\_\allowbreak FOUR — C — Actual SomeIn-to-AllIn propagation}

Same source Space and PM ideal, residual point Good4, original bound 30(C+1)\textasciicircum{}2 <= d; geometry must preserve all identities.

\emph{Source.} \href{https://github.com/CocoSgt/riemann-hypothesis-lean/blob/7cbf7a0c32fca1badf05931879a37c3f8a2b03b4/src/Development/AutomorphicArithmetic/AutomorphicVarieties/ModularityArgument/ThresholdArguments/DenseThresholdCases/PromodularityCaseThree.lean#L26-L32}{ChainRankFour}.



\paragraph{11. CHAIN\_\allowbreak DESC4 — B — Original residual chain descent}

Finite residual descent retaining source Chain4 compatibility.

\emph{Source.} \href{https://github.com/CocoSgt/riemann-hypothesis-lean/blob/7cbf7a0c32fca1badf05931879a37c3f8a2b03b4/src/Development/AutomorphicArithmetic/AutomorphicVarieties/ModularityArgument/ThresholdArguments/DenseThresholdCases/PromodularityCaseThreeOfLemmas.lean#L131-L153}{UpperTheory.V.Cu.promodularity\_\allowbreak case\_\allowbreak three\_\allowbreak of\_\allowbreak chain\_\allowbreak reduction}.



\paragraph{12. RANK4\_\allowbreak CHAIN\_\allowbreak REDUCTION — C — Seeded reduction to the original residual chain}

Select T1 before every later Deep tower; retain one compatible seed and Chain4 for every original residual lift over Q.

\emph{Source.} \href{https://github.com/CocoSgt/riemann-hypothesis-lean/blob/7cbf7a0c32fca1badf05931879a37c3f8a2b03b4/src/Development/AutomorphicArithmetic/AutomorphicVarieties/ModularityArgument/ThresholdArguments/DenseThresholdCases/PromodularityCaseThree.lean#L33-L39}{RankFourChainReductionDense}.

\emph{Source.} \href{https://github.com/CocoSgt/riemann-hypothesis-lean/blob/7cbf7a0c32fca1badf05931879a37c3f8a2b03b4/src/Development/AutomorphicArithmetic/AutomorphicVarieties/ModularityArgument/ThresholdArguments/DenseThresholdCases/CuPChainIV.lean#L16-L82}{UpperTheory.V.Cu.chainIV\_\allowbreak alt}.



\paragraph{13. TASK\_\allowbreak COUNT\_\allowbreak V\_\allowbreak INSTALL — B — Complete the original CountV finite assembly}

From exactly the listed source host/cohomology inputs, complete countV\_\allowbreak of\_\allowbreak finiteness\_\allowbreak duality\_\allowbreak euler\_\allowbreak host, including original host/frame formation, normalized length, the CaseV guard and R1<=2\textasciicircum{}c2.

\emph{Completed inputs (version 0.56).} APD full original presentation on the same arithmetic local objects; sharp chains/seeds/propagation retained, ALC actual local objects and charts; full dimensions/seeds/propagation retained, RTC full rank-two preconnected core and coefficient descent; actual propagation/counting inputs retained, ACP full finite-point and abstract propagation implications; actual rank-specific and counting inputs retained, AFE full Host entry/auxiliary frame/local witnesses including ell=3, LFC1-LFC23, CAF9-CAF11, MFC1-MFC14, Inherited CountV consumers

\emph{Remaining producer scope.} APD supplies the full original local presentation pres for every supplied local object at both ranks: an O-prime power-series surjection onto the actual R with N variables and N-prime listed generators of its kernel, with N=N-prime+n*F.d+1. It constructs the broad point\_\allowbreak loc, ext\_\allowbreak A and def\_\allowbreak iso/actual NRep.Univ using finite quotients of finitely generated graded Laurent vertex algebras and compact coefficient recovery. It proves actual small-extension obstruction and tangent maps, both normalization corrections, the weight-torus contribution to the tangent and invariant-space count, and the full cyclotomic-dual vanishing calculation; the original finite-module Euler field then gives the exact tangent excess. Neither F.Big nor a nontriangularity hypothesis is added to pres, and Host ell=3 remains allowed. The earlier ALC completion, LocFlat/LocRho and local charts over the coefficient maximal ideal remain complete. Remaining are c1ch with full ClComp/ClChain, triangular and coincident-character dimension suppliers, sharp2/sharp3, the distinct rank-three EPres data, nice-prime properties, closed-point density and seed incidence, full rank-two propagation, the rank-two-to-three corner and exceptional ideal, rank-three propagation and OffFin. Existence of originally supplied universal spaces, actual automorphic seeds, patching/modularity, classicality and stronger Alt WitA/WitAQ outputs remain separate. Preserve the common marked frame, original arithmetic P, rank-specific rings and ideals, all coefficient maps and both actual O-prime extensions. No upstream FinA conclusion is used. Complete the original F3 local coefficient bridges and F4 common-overfield readings, with valuation/norm/GoodCoef/iq and the entire coefficient field, and actual complex Hecke standard/no-Jordan plus arithmetic character data on the same form datum. Construct the actual base rows, host/frame lengths, nonzero chains and group ties, MultEq support, eigensystem/carrier/parts, all-embedding logLam, HD/d3b and NEb. Keep the stronger second DivRows3 scope distinct from the completed first DivDef3 and row-pair integrability. The full arithmetic entry and conditional common-frame construction do not by themselves install HOST\_\allowbreak NUMBER\_\allowbreak S, HOST\_\allowbreak FRAME\_\allowbreak LENGTH, TASK\_\allowbreak HOST\_\allowbreak VAN\_\allowbreak E\_\allowbreak INSTALL, TASK\_\allowbreak HOST\_\allowbreak XPP\_\allowbreak INSTALL or TASK\_\allowbreak COUNT\_\allowbreak V\_\allowbreak INSTALL.

\emph{Source.} \href{https://github.com/CocoSgt/riemann-hypothesis-lean/blob/7cbf7a0c32fca1badf05931879a37c3f8a2b03b4/src/Development/AutomorphicArithmetic/AutomorphicVarieties/HeightAndVanishing/Counts.lean#L69-L69}{CountV}.

\emph{Source.} \href{https://github.com/CocoSgt/riemann-hypothesis-lean/blob/7cbf7a0c32fca1badf05931879a37c3f8a2b03b4/src/Development/AutomorphicArithmetic/AutomorphicVarieties/ModularityArgument/ThresholdArguments/DenseThresholdCases/CountVOf.lean#L33-L60}{UpperTheory.V.UnitaryThreeInertia.countV\_\allowbreak of\_\allowbreak finiteness\_\allowbreak duality\_\allowbreak euler\_\allowbreak host}.



\paragraph{14. THM\_\allowbreak S6\_\allowbreak S — C — Original unitary-three ordinary support and automorphy}

Original host/frame and fixed iota -> one marked g1 -> every universal R,u -> weight-zero finiteness and every Type-I BPrime\_\allowbreak s point.

\emph{Source.} \href{https://github.com/CocoSgt/riemann-hypothesis-lean/blob/7cbf7a0c32fca1badf05931879a37c3f8a2b03b4/src/Development/AutomorphicArithmetic/AutomorphicVarieties/ModularityArgument/ThresholdArguments/DenseThresholdCases/ThmS6CoreS.lean#L17-L26}{ThmS6\_\allowbreak s}.

\emph{Source.} \href{https://github.com/CocoSgt/riemann-hypothesis-lean/blob/7cbf7a0c32fca1badf05931879a37c3f8a2b03b4/src/Development/AutomorphicArithmetic/AutomorphicVarieties/ModularityArgument/ThresholdArguments/DenseThresholdCases/ThmS6CoreS.lean#L12-L16}{BPrime\_\allowbreak s}.



\paragraph{15. HOST\_\allowbreak FRAME\_\allowbreak LENGTH — B — Original Case-V host/frame and normalized-length installation}

Actual HostData, frame, local-condition equality, abstract count and normalized original Selmer length must all refer to the same objects.

\emph{Completed inputs (version 0.56).} AFE/LFC full entry and conditional common frame, CAF and MFC, DNE weight/host-degree dictionary at supplied tables

\emph{Remaining producer scope.} Complete the original F3 local coefficient bridges and F4 common-overfield readings, with valuation/norm/GoodCoef/iq and the entire coefficient field, and actual complex Hecke standard/no-Jordan plus arithmetic character data on the same form datum. Construct the actual base rows, host/frame lengths, nonzero chains and group ties, MultEq support, eigensystem/carrier/parts, all-embedding logLam, HD/d3b and NEb. Keep the stronger second DivRows3 scope distinct from the completed first DivDef3 and row-pair integrability. The full arithmetic entry and conditional common-frame construction do not by themselves install HOST\_\allowbreak NUMBER\_\allowbreak S, HOST\_\allowbreak FRAME\_\allowbreak LENGTH, TASK\_\allowbreak HOST\_\allowbreak VAN\_\allowbreak E\_\allowbreak INSTALL, TASK\_\allowbreak HOST\_\allowbreak XPP\_\allowbreak INSTALL or TASK\_\allowbreak COUNT\_\allowbreak V\_\allowbreak INSTALL.

\emph{Source.} \href{https://github.com/CocoSgt/riemann-hypothesis-lean/blob/7cbf7a0c32fca1badf05931879a37c3f8a2b03b4/src/Development/AutomorphicArithmetic/AutomorphicVarieties/ModularityArgument/ThresholdArguments/DenseThresholdCases/CountVOf.lean#L33-L60}{UpperTheory.V.UnitaryThreeInertia.countV\_\allowbreak of\_\allowbreak finiteness\_\allowbreak duality\_\allowbreak euler\_\allowbreak host}.



\paragraph{16. HOST\_\allowbreak NUMBER\_\allowbreak S — C — Uniform original-host rank-three catalogue}

One exponent before all closures, coefficient rings, levels, characters, frames, embeddings and marked points; one separating trace function controls all later finite pairwise nonisomorphic collections.

\emph{Completed inputs (version 0.56).} AFE full Host entry at original scope, LFC full conditional common frame, CAF and MFC, Inherited host-count consumers

\emph{Remaining producer scope.} Complete the original F3 local coefficient bridges and F4 common-overfield readings, with valuation/norm/GoodCoef/iq and the entire coefficient field, and actual complex Hecke standard/no-Jordan plus arithmetic character data on the same form datum. Construct the actual base rows, host/frame lengths, nonzero chains and group ties, MultEq support, eigensystem/carrier/parts, all-embedding logLam, HD/d3b and NEb. Keep the stronger second DivRows3 scope distinct from the completed first DivDef3 and row-pair integrability. The full arithmetic entry and conditional common-frame construction do not by themselves install HOST\_\allowbreak NUMBER\_\allowbreak S, HOST\_\allowbreak FRAME\_\allowbreak LENGTH, TASK\_\allowbreak HOST\_\allowbreak VAN\_\allowbreak E\_\allowbreak INSTALL, TASK\_\allowbreak HOST\_\allowbreak XPP\_\allowbreak INSTALL or TASK\_\allowbreak COUNT\_\allowbreak V\_\allowbreak INSTALL.

\emph{Source.} \href{https://github.com/CocoSgt/riemann-hypothesis-lean/blob/7cbf7a0c32fca1badf05931879a37c3f8a2b03b4/src/Development/AutomorphicArithmetic/AutomorphicVarieties/ModularityArgument/ThresholdArguments/DenseThresholdCases/HostNumberS.lean#L18-L28}{HostNumber\_\allowbreak s}.



\paragraph{17. TASK\_\allowbreak HOST\_\allowbreak VAN\_\allowbreak E\_\allowbreak INSTALL — B — Transport K/duality to literal HostVanE}

Use the already supplied LemmaK/local-global duality and finiteness on the actual CountV host/frame and exact original coefficient modules.

\emph{Source.} \href{https://github.com/CocoSgt/riemann-hypothesis-lean/blob/7cbf7a0c32fca1badf05931879a37c3f8a2b03b4/src/Development/AutomorphicArithmetic/AutomorphicVarieties/UnitaryThreeDeformations/HostNumber.lean#L9-L12}{HostVanE}.

\emph{Source.} \href{https://github.com/CocoSgt/riemann-hypothesis-lean/blob/7cbf7a0c32fca1badf05931879a37c3f8a2b03b4/src/Development/AutomorphicArithmetic/AutomorphicVarieties/LevelCohomology/SoHostVanE.lean#L11-L18}{UpperTheory.V.X.So.hostVanE}.



\paragraph{18. TASK\_\allowbreak HOST\_\allowbreak XPP\_\allowbreak INSTALL — B — Transport CorX to literal HostXpp}

Preserve Psi\textasciicircum{}-2 eta, the prime-to-ell quotient, actual local conditions and the original CountV host while installing BR7.

\emph{Source.} \href{https://github.com/CocoSgt/riemann-hypothesis-lean/blob/7cbf7a0c32fca1badf05931879a37c3f8a2b03b4/src/Development/AutomorphicArithmetic/AutomorphicVarieties/UnitaryThreeDeformations/HostNumber.lean#L14-L17}{HostXpp}.

\emph{Source.} \href{https://github.com/CocoSgt/riemann-hypothesis-lean/blob/7cbf7a0c32fca1badf05931879a37c3f8a2b03b4/src/Development/AutomorphicArithmetic/AutomorphicVarieties/LevelCohomology/Xpp.lean#L11-L84}{UpperTheory.V.X.xpp}.



\paragraph{19. DESIGN\_\allowbreak CHOICES\_\allowbreak COMMON — C — Construct the original common design family and bundle}

For the one target HostExW, construct the original compatible DesignChoices on Cdes=WithTw(WithSm Cprime), where Cprime=WithChN(A intersect target class), A is the one Built construction class, and cT(eps,y)=cTmin(y). Retain every original field, tuple, word, metric and common bundle duty. DSZ10 and RMP11 give size/rank properties of class members, not member existence. The original RecExP/RecAtG and common object/row route remains coupled. Construct ONE original common design family R and its recursive/geometric half with original Inputs/Pins/Holds/ExAt/InGeo and required CoreP. PKT now reconstructs full exists\_\allowbreak pkT, rowPk\_\allowbreak nonempty and exists\_\allowbreak moved with Copy/Coh on those original given inputs and a supplied fitting later stage: one catalogue, both ChT, complete LfT, actual integral lattices and all metric/measure/period constants are retained; one numerical record precedes every eta and V/S/N/G. RDR supplies one Side0 before all readers, ReadPk/pins and full ReadL with exact o.eta=q.eta. The initial actual geometry, Fits/division/central conditions, primitive rational tests and eta agreement must still be produced. Inherited CSR supplies finite columns, one ReadC before later twists, period/evaluation and scalar/error transports, RowSideDetermined and same-R Bundle/BTw/BWit/BunU consumers. RPI supplies exact RowPairInt for actual pairs and finite translates with original mu\_\allowbreak b/fd0; first literal DivDef3 at the same updated R,F,Y,Z,C,sEE follows from its retained original readings, Pins/Holds, stored nonzero IsItW chain, five real-group predicates and group tie. Remaining common fields include host/frame, MultEq nonzero-part support, that original nonzero chain/group tie, eigensystem/carrier/parts and all-embedding logLam, HD/d3b, NEb, actual rational tests and base rows. Retain every simultaneous family hypothesis and eventual threshold. Stronger second DivRows3 remains separate; universal quotient-height CuspRapidDecay and actual common member existence are not inferred from the comparison constructors. Version 0.41 supplies the finite phase-correct rational theta route: exact coefficient extension, split representatives, square-law/PhiIs, Pins from CoversAt, class-sum/point-identity transfer to OneMapOverL, actual-good-projector RatG and both-level LiftTestB. Its original higher covers, all-tuple class-sum uniqueness including phase rescaling, one PointIdentCore before O, base ClassSumsSpan and geometric LatFull remain explicit producers. MLat is available from the original LatFull input. This sharpens the rational-test work without claiming actual common geometry, Fits or eta agreement. Version0.46 proves finite generation of the original whole integral cusp lattice by actual integral-frame evaluation. It also proves the exact identity-cover, analytic unfolding, original-field coefficient descent, global clearing-scalar and both-level lattice-transport implications. Full spanning retains the supplied original IsNab connection and the five explicit principal-chart, level-system, raw-Gaussian, short-rule/dictionary and ramified-jet producers. The higher tuple covers, actual common class sums, PointIdentCore and ClassSumsSpan remain separate. CURRENT VERSION0.47 QUALIFICATION: Construct one original common geometric family and the higher CoversAt/all-tuple class sums and uniqueness with original phase rescalings, one PointIdentCore with common lambda,z, and base ClassSumsSpan. CNC now supplies the full original IsNab; PCL supplies the conditional trace-dual principal-lattice and quantitative normalization. Actual uniformization/HasPrinc and the five principal-chart, finite-level-system, Gaussian/reproduction, short-rule/simultaneous-dictionary and ramified-jet producer groups remain. Keep all original both-level, coefficient-field, Fits and eta agreements. Completed new original producers: ICL1-ICL21, CNC1-CNC16, PCL1-PCL12. Any earlier narrative naming these completed producers as outstanding is historical; this qualification is the current scope authority. CURRENT VERSION0.48 QUALIFICATION: Construct the original finite holomorphic FamilyUnif/N1c, normalized theta families and local lifting, then the common principal charts and finite-level theta systems. PUN now supplies the exact fixed-space models, all-family center lattices, A=1 exponent matching and individual primitive points with the full point-equality criterion; CNC supplies the full canonical connection and PCL the quantitative trace-dual normalization. Complete higher CoversAt and all-tuple class sums/uniqueness with the original phase rescalings, one PointIdentCore with common lambda,z, and base ClassSumsSpan. The Gaussian/reproduction, short-rule/simultaneous-dictionary and ramified-jet producers, and all original both-level, coefficient-field, Fits and eta agreements remain explicit. Completed current producers or conditional consumers: ICL1-ICL21, CNC1-CNC16, PCL1-PCL12, PUN1-PUN24, GCF1-GCF18 at their stated scopes. Earlier frontier descriptions are historical; this qualification is the current scope authority. Original whole-contract premises and quantifiers are retained. CURRENT VERSION0.49 QUALIFICATION: Construct the adjoining simultaneous principal framed points and test-ring dictionary. Full finite holomorphic FamilyUnif/N1c, normalized theta families and local lifts are now supplied by HUF. The next bounded target is allTest\_\allowbreak of\_\allowbreak secBaseChange: represent FrPt by the Lie-eigenframe bundles and theta-section frame line, use the full n=1 SecBaseChangeSt on every nonflat coefficient base change, and prove finite-generated subalgebra descent with the exact TRing/FromTest transport. Then supply the full FramesOK/CoverTri/FibTri and IsUnf/hasPrinc package, including adapted mod-eight symplectic basis, normalized theta multiplier/weight and arithmetic invariance. Continue higher CoversAt, all-tuple class sums and uniqueness with original phase rescalings, one PointIdentCore with common lambda,z, base ClassSumsSpan, Gaussian/reproduction, short-rule dictionaries and ramified jets, retaining both-level, coefficient-field, Fits and eta agreements. Common-family nonemptiness remains distinct from compatibility of an already supplied member. Completed current producers or conditional consumers: ICL1-ICL21, CNC1-CNC16, PCL1-PCL12, PUN1-PUN24, HUF1-HUF21 (full original FamilyUnif/N1c), GCF1-GCF18 and WRS1-WRS10 at their stated scopes. Earlier frontier descriptions are historical; this qualification is the current scope authority. Original whole-contract premises and quantifiers are retained.

\emph{Completed inputs (version 0.56).} PBC1–PBC27 (full arbitrary-base polarized pullback, actual very-ampleness and section finiteness), SHV1–SHV21 (both original Serre conclusions, actual finite section model and full all-cover return), ASL1–ASL21 (actual section localization and full original arbitrary-base Perf producer), FCD1-FCD18 (full homogeneous-chart descent, graded presentation and Noetherian cube model), GSL1-GSL26 (full conditional section laws, universal basis and projective-model algebra), NFL1-NFL22, TDR1-TDR13, NFR1-NFR13, HUF1-HUF21, PUN1-PUN24, ICL1-ICL21, CNC1-CNC16, PCL1-PCL12, GCF/WRS at original supplied-member scope

\emph{Remaining producer scope.} PBC supplies full arbitrary-base RelAmpleOfLocallyOnBase and literal PolBaseChangeSt, including the actual Cartesian square, unchanged dimension, actual pulled line bundle, generated comparison, rigidification, symmetry, all-test group multiplication and all-affine-test KTriv. The finite determinant trick, extra-section affineness and two-exponent section extension give the synchronized positive-power glue\_\allowbreak pack and exists\_\allowbreak sv without a Noetherian or flatness restriction on base change. PBC26 now supplies full polAV\_\allowbreak veryAmpleBy on every original affine test and specified Cartesian model; its power may depend on that test. PBC27 derives the finite-relation property from Perf over any field and proves original sections\_\allowbreak finite\_\allowbreak of\_\allowbreak perf and every-integer-power SecM finiteness over the original algebraically closed field. Its separate arbitrary-field finite\_\allowbreak secM implication retains the supplied finite-relation premise. GSL\_\allowbreak POL\_\allowbreak BASE\_\allowbreak CHANGE is supplied, alongside the already complete GSL\_\allowbreak SERRE and ASL Perf inputs. Remaining for GSL17 are multiplication-pullback generated comparison, finite/flat/surjective and etale multiplication geometry, and invariant-function descent on every relevant open. Geometric integrality, principal rank one, torsion-indexed section decomposition, two-sided growth, and square/translation generators remain for full P,V,R,T and the GSL26/NFL joins; section finiteness alone does not supply the growth half of GSL\_\allowbreak FINITE\_\allowbreak GROWTH. Original Moduli and HUF/PUN foundations remain. Later finite-level systems, Gaussian/reproduction, simultaneous dictionaries, ramified jets, all-tuple class sums, phase-sensitive uniqueness, common PointIdentCore and base ClassSumsSpan remain, retaining both levels, whole fields, Fits and eta agreement.



\paragraph{20. LOCAL\_\allowbreak NBUILT\_\allowbreak COMMON — C — Actual local specialization on the common selected family}

Construct, with the required closed-class property on the same Des(C,3), one eventual family of actual NBuilt certificates for the original Fm.D(x,i), retaining the certificate before forgetting to TSel. HLP now reconstructs full HeckeLift0, all-index commutation, HeckeStd0 and HeckeIntTw from the original MtInput and complex Hecke facts, in the full ambient GenAt range k<=nV. Every maximal ideal is an actual supported residual character class; literal MtSub is the O-span of projections of the original MtLat generators, is finite free, and localizes onto the exact NCov.part. Literal Psi has its exact range and off-space projection error, both original span assertions hold, and ActEx/ActHK are constructed on the same lattice under original semisimplicity. Original MjLat relation quotients, universal coefficient-map relation fields, selected tame generators, coefficient embeddings and all identity transports give the complete nonempty\_\allowbreak spec constructor on its full printed input list. A single compositum combines the four original coefficient requirements with big fixed beforehand. Whole ndSpec\_\allowbreak exists still needs actual common geometric/choice MtInput, complex HeckeStdPlain and the no-Jordan producer, original F3 bridges with valuation/GoodCoef/iq, original F4 common-overfield reading of the entire coefficient field, and exact arithmetic character/reciprocity/admissibility suppliers. Original AutSpec, Galois specialization and the remaining NBuilt fields on that same installation also remain. Do not replace the actual projected lattice or maps by abstract objects with the same dimensions.

\emph{Completed inputs (version 0.56).} PBC full polarized pullback; APD exact actual presentation; RCV reference covolume; WPR original conditional WachEx, SHV both Serre conclusions; ALC actual arithmetic charts; SVW same-setting volume join, ASL full actual perfect-complex model; RTC common-frame arithmetic inputs; AOS exact unchanged infinity normalization, FCD full algebraic model, ACP abstract arithmetic suppliers, GMR same-setting measure normalization, GSL/NFL full conditional section laws and uniformization, AFE/LFC full entry and conditional common frame, TDR/NFR and CAF/MFC, WRS selected built-class extraction at membership

\emph{Remaining producer scope.} Continue the distinct retained-defect, period-height and normalized spectral producers, and the original compatible second-layer selection on one common arithmetic/geometric family. GLOBAL\_\allowbreak LOCAL\_\allowbreak GAIN remains fully closed at its original v0.49 scope. Complete local comparison theorems do not themselves construct the whole common family. Common design existence, NBUILT compatibility, original true-column records, shared fields/embeddings, Fits and eta agreement, simultaneous rows and eventual-family thresholds remain explicit. A stronger historical WPt/TSel construction may serve other producers but is not needed to reopen the proved selected-class local-gain estimate. PBC supplies full arbitrary-base RelAmpleOfLocallyOnBase and literal PolBaseChangeSt, including the actual Cartesian square, unchanged dimension, actual pulled line bundle, generated comparison, rigidification, symmetry, all-test group multiplication and all-affine-test KTriv. The finite determinant trick, extra-section affineness and two-exponent section extension give the synchronized positive-power glue\_\allowbreak pack and exists\_\allowbreak sv without a Noetherian or flatness restriction on base change. PBC26 now supplies full polAV\_\allowbreak veryAmpleBy on every original affine test and specified Cartesian model; its power may depend on that test. PBC27 derives the finite-relation property from Perf over any field and proves original sections\_\allowbreak finite\_\allowbreak of\_\allowbreak perf and every-integer-power SecM finiteness over the original algebraically closed field. Its separate arbitrary-field finite\_\allowbreak secM implication retains the supplied finite-relation premise. GSL\_\allowbreak POL\_\allowbreak BASE\_\allowbreak CHANGE is supplied, alongside the already complete GSL\_\allowbreak SERRE and ASL Perf inputs. Remaining for GSL17 are multiplication-pullback generated comparison, finite/flat/surjective and etale multiplication geometry, and invariant-function descent on every relevant open. Geometric integrality, principal rank one, torsion-indexed section decomposition, two-sided growth, and square/translation generators remain for full P,V,R,T and the GSL26/NFL joins; section finiteness alone does not supply the growth half of GSL\_\allowbreak FINITE\_\allowbreak GROWTH. Original Moduli and HUF/PUN foundations remain. Later finite-level systems, Gaussian/reproduction, simultaneous dictionaries, ramified jets, all-tuple class sums, phase-sensitive uniqueness, common PointIdentCore and base ClassSumsSpan remain, retaining both levels, whole fields, Fits and eta agreement.



\paragraph{21. SECOND\_\allowbreak LAYER\_\allowbreak COMPATIBLE\_\allowbreak SELECTION — C — Original second-layer selector and class duties}

For the same A, supply Duty on A and a Closed common original TSel family whose selected records are the original true-column rec7 records (or have the entire dependent tuple transported), with their real witness and original coefficient-field root data. Existence/compatibility remains arithmetic; merely abstractly isomorphic fields or independently selected columns do not suffice. CSR21–22 now construct the same-R Bundle/BTw/BWit/BunU outputs from the exact selected realization/row-side packet and NEb, with the original EvP/ParEv quantifier order. Actual common selection, the five stated realization packets and NEb remain upstream; none is replaced by an independently chosen compatible record. Version 0.38 closes the selected analytic extension and reconstructs the first updated DivDef3 consumer; it leaves the original common selection, geometric readings, rational tests, base rows, host/parts/NEb and exact retained DI premises upstream, with unchanged EvP/ParEv thresholds. Version 0.39 supplies the original given-data reader constructors and exact RowPairInt/first DivDef3 on original premises. It leaves one actual common geometric/comparison/test member, eta agreement, compatible selection, rational tests, base rows, host/parts/NEb and original EvP/ParEv thresholds upstream. Version 0.40 also supplies original comparison/RowPk/fitting-stage Moved algorithms and exact common data transport. The actual initial geometric member, fitting tests, eta agreement and compatible second-layer selector with original eventual thresholds remain upstream. Version 0.41 supplies the finite phase-correct rational theta route: exact coefficient extension, split representatives, square-law/PhiIs, Pins from CoversAt, class-sum/point-identity transfer to OneMapOverL, actual-good-projector RatG and both-level LiftTestB. Its original higher covers, all-tuple class-sum uniqueness including phase rescaling, one PointIdentCore before O, base ClassSumsSpan and geometric LatFull remain explicit producers. MLat is available from the original LatFull input. This sharpens the rational-test work without claiming actual common geometry, Fits or eta agreement. CURRENT VERSION0.49 QUALIFICATION: Continue the distinct global denominator, period/height, normalized spectral and second-layer compatible-selection producers, together with existence of the original common arithmetic/geometric family. GLOBAL\_\allowbreak LOCAL\_\allowbreak GAIN and its actual GroupCompatibilityFamily supplier are now complete on the unchanged selected built class; no additional WPt/TSel realization is needed for that estimate. Do not reopen local first-growth, mean, auxiliary admission, selected D3 tops/Adm2F, the single GIs2 stage or the exact comparison/kernel transport. A full RecT or TSel may remain useful for other tasks, but is not claimed by this local-gain closure. Completed current producers or conditional consumers: GCF1-GCF18 at their stated scopes, WRS1-WRS10 original selected-class compatibility/local gain. Earlier frontier descriptions are historical; this qualification is the current scope authority. Original whole-contract premises and quantifiers are retained.

\emph{Completed inputs (version 0.56).} PBC/APD/RCV/WPR full named original producers on unchanged actual objects, IFD full invariant matrix descent; SHV/ALC/SVW full named nested producers, WIS actual inclusion and invariant completeness; AOS exact scalar; ASL/RTC full named nested producers, WLD full generic lattice constructors, GMR actual-setting normalization, FCD/ACP exact nested suppliers, WRS1-WRS10 original selected-class compatibility/local gain, DNE full denominator consumer, DFD/SSC canonical-data and signature comparisons, GSL/NFL and AFE/LFC conditional constructors, Inherited exact same-record realization/transport consumers

\emph{Remaining producer scope.} Continue the distinct retained-defect, period-height and normalized spectral producers, and the original compatible second-layer selection on one common arithmetic/geometric family. GLOBAL\_\allowbreak LOCAL\_\allowbreak GAIN remains fully closed at its original v0.49 scope. Complete local comparison theorems do not themselves construct the whole common family. Common design existence, NBUILT compatibility, original true-column records, shared fields/embeddings, Fits and eta agreement, simultaneous rows and eventual-family thresholds remain explicit. A stronger historical WPt/TSel construction may serve other producers but is not needed to reopen the proved selected-class local-gain estimate.



\paragraph{22. BASE\_\allowbreak THETA\_\allowbreak READINGS — C — Actual base reading on the common bundle}

For every family in the original class, eventually on the SAME selected Bundle and for each i<m, supply the actual R,F,Read,ReadL, weak models/Weil data and phase-one pins. Identify their LiftTie.theta, mtRead, mjRead and ThDef3 basis with that Bundle.P.A3, P.Zs, P.S0 and its stored maps. Exact field, embeddings, metrics, measures, row family and weak choices are retained. This remains one part of the coupled BTH12 producer. PKT now supplies the complete original comparison packet, RowPk and fitting-stage Moved with Copy/Coh on original given geometry, Inputs/Pins/Holds/ExAt/InGeo and Fits. Its ChT/LfT/Reps, actual lattice carriers, original Haar and norm constants, period maps, representatives and all-place RrT are constructed together; its one numerical record precedes every eta/V/S/N/G. RDR gives the original level-zero side, moved-Read and ReadPk/pin constructors on those objects, with one Side0 before every reader and all twelve Read ties. Full ReadL retains the exact o.eta(j,v)=q.eta(j,v) guard for every 1<=j<=m and place v; arbitrary Over does not imply it. Inherited CSR fixes one ReadC before all TwArgs, transports both metrics and all thirteen reading maps, and gives actual period/evaluation with kappa=(c*nP)\textasciicircum{}(-1/2). RPI supplies exact row-pair integrability, so first updated DivDef3 follows from the retained original Pins/Holds, nonzero chain, five real-group predicates and group tie. Initial theta extension remains complete. Remaining work is actual initial common geometry and its original predicates, Fits/primitive rational tests, eta agreement and outer family/eventual/same-Bundle duties, reFace compatibility, MultEq support, original chain/group data, actual base rows and any stronger DivRows3 theta square. No new compatible reading, family threshold or Haar normalization is substituted. Version 0.41 supplies the finite phase-correct rational theta route: exact coefficient extension, split representatives, square-law/PhiIs, Pins from CoversAt, class-sum/point-identity transfer to OneMapOverL, actual-good-projector RatG and both-level LiftTestB. Its original higher covers, all-tuple class-sum uniqueness including phase rescaling, one PointIdentCore before O, base ClassSumsSpan and geometric LatFull remain explicit producers. MLat is available from the original LatFull input. This sharpens the rational-test work without claiming actual common geometry, Fits or eta agreement.

\emph{Completed inputs (version 0.56).} PBC1–PBC27 (full arbitrary-base polarized pullback, actual very-ampleness and section finiteness), SHV1–SHV21 (both original Serre conclusions, actual finite section model and full all-cover return), ASL1–ASL21 (actual section localization and full original arbitrary-base Perf producer), FCD1-FCD18 (full homogeneous-chart descent, graded presentation and Noetherian cube model), GSL1-GSL26, NFL1-NFL22, TDR1-TDR13, NFR1-NFR13, HUF1-HUF21, PUN1-PUN24, CNC1-CNC16, Inherited actual-reading consumers at original inputs

\emph{Remaining producer scope.} PBC supplies full arbitrary-base RelAmpleOfLocallyOnBase and literal PolBaseChangeSt, including the actual Cartesian square, unchanged dimension, actual pulled line bundle, generated comparison, rigidification, symmetry, all-test group multiplication and all-affine-test KTriv. The finite determinant trick, extra-section affineness and two-exponent section extension give the synchronized positive-power glue\_\allowbreak pack and exists\_\allowbreak sv without a Noetherian or flatness restriction on base change. PBC26 now supplies full polAV\_\allowbreak veryAmpleBy on every original affine test and specified Cartesian model; its power may depend on that test. PBC27 derives the finite-relation property from Perf over any field and proves original sections\_\allowbreak finite\_\allowbreak of\_\allowbreak perf and every-integer-power SecM finiteness over the original algebraically closed field. Its separate arbitrary-field finite\_\allowbreak secM implication retains the supplied finite-relation premise. GSL\_\allowbreak POL\_\allowbreak BASE\_\allowbreak CHANGE is supplied, alongside the already complete GSL\_\allowbreak SERRE and ASL Perf inputs. Remaining for GSL17 are multiplication-pullback generated comparison, finite/flat/surjective and etale multiplication geometry, and invariant-function descent on every relevant open. Geometric integrality, principal rank one, torsion-indexed section decomposition, two-sided growth, and square/translation generators remain for full P,V,R,T and the GSL26/NFL joins; section finiteness alone does not supply the growth half of GSL\_\allowbreak FINITE\_\allowbreak GROWTH. Original Moduli and HUF/PUN foundations remain. Later finite-level systems, Gaussian/reproduction, simultaneous dictionaries, ramified jets, all-tuple class sums, phase-sensitive uniqueness, common PointIdentCore and base ClassSumsSpan remain, retaining both levels, whole fields, Fits and eta agreement. Complete the original F3 local coefficient bridges and F4 common-overfield readings, with valuation/norm/GoodCoef/iq and the entire coefficient field, and actual complex Hecke standard/no-Jordan plus arithmetic character data on the same form datum. Construct the actual base rows, host/frame lengths, nonzero chains and group ties, MultEq support, eigensystem/carrier/parts, all-embedding logLam, HD/d3b and NEb. Keep the stronger second DivRows3 scope distinct from the completed first DivDef3 and row-pair integrability. The full arithmetic entry and conditional common-frame construction do not by themselves install HOST\_\allowbreak NUMBER\_\allowbreak S, HOST\_\allowbreak FRAME\_\allowbreak LENGTH, TASK\_\allowbreak HOST\_\allowbreak VAN\_\allowbreak E\_\allowbreak INSTALL, TASK\_\allowbreak HOST\_\allowbreak XPP\_\allowbreak INSTALL or TASK\_\allowbreak COUNT\_\allowbreak V\_\allowbreak INSTALL.



\paragraph{23. BASE\_\allowbreak THETA\_\allowbreak RATIONAL\_\allowbreak TESTS — C — Simultaneous original base rational tests}

On the SAME original Bundle/reading/theta map, for every i<m supply one finite separating complex-linear test family whose values lie in the fixed embedded coefficient field on theta of EVERY original rational source vector and on EVERY stored target-coordinate basis vector. BTH8/BTH11 use this same LiftTestB for rational adjoint image and matrix coordinates. RTG now reconstructs prescribed-field coefficient extension, original split class representatives/RepsOK, the self-dual split-lattice square law and finite residue PhiIs; every Pins field follows from the actual CoversAt. With the base coefficient record literally fixed, exact class-sum existence/uniqueness including phase-rescaled O, PointIdentCore with common lambda,z before every O, and base ClassSumsSpan imply all-tuple finite lift span and OneMapOverL. The phase is exactly cj(twAt embLL A3.A.v0)(R.phs i). RatG uses the actual good projector eOp(y), and full LiftTestB follows at BOTH original levels. These finite implication proofs are complete. Remaining original producers are the prescribed higher covers, those actual class sums/uniqueness, PointIdentCore, base ClassSumsSpan and the base/common geometric datum; exact MLat is supplied only from original LatFull, whose geometric production remains separate. Finite dimension or an unrelated rational theta model cannot replace these same-object identities. Original outer eventual-family/Bundle quantifiers remain. Version0.46 proves finite generation of the original whole integral cusp lattice by actual integral-frame evaluation. It also proves the exact identity-cover, analytic unfolding, original-field coefficient descent, global clearing-scalar and both-level lattice-transport implications. Full spanning retains the supplied original IsNab connection and the five explicit principal-chart, level-system, raw-Gaussian, short-rule/dictionary and ramified-jet producers. The higher tuple covers, actual common class sums, PointIdentCore and ClassSumsSpan remain separate. CURRENT VERSION0.47 QUALIFICATION: Construct one original common geometric family and the higher CoversAt/all-tuple class sums and uniqueness with original phase rescalings, one PointIdentCore with common lambda,z, and base ClassSumsSpan. CNC now supplies the full original IsNab; PCL supplies the conditional trace-dual principal-lattice and quantitative normalization. Actual uniformization/HasPrinc and the five principal-chart, finite-level-system, Gaussian/reproduction, short-rule/simultaneous-dictionary and ramified-jet producer groups remain. Keep all original both-level, coefficient-field, Fits and eta agreements. Completed new original producers: ICL1-ICL21, CNC1-CNC16, PCL1-PCL12. Any earlier narrative naming these completed producers as outstanding is historical; this qualification is the current scope authority. CURRENT VERSION0.48 QUALIFICATION: Construct the original finite holomorphic FamilyUnif/N1c, normalized theta families and local lifting, then the common principal charts and finite-level theta systems. PUN now supplies the exact fixed-space models, all-family center lattices, A=1 exponent matching and individual primitive points with the full point-equality criterion; CNC supplies the full canonical connection and PCL the quantitative trace-dual normalization. Complete higher CoversAt and all-tuple class sums/uniqueness with the original phase rescalings, one PointIdentCore with common lambda,z, and base ClassSumsSpan. The Gaussian/reproduction, short-rule/simultaneous-dictionary and ramified-jet producers, and all original both-level, coefficient-field, Fits and eta agreements remain explicit. Completed current producers or conditional consumers: ICL1-ICL21, CNC1-CNC16, PCL1-PCL12, PUN1-PUN24. Earlier frontier descriptions are historical; this qualification is the current scope authority. Original whole-contract premises and quantifiers are retained. CURRENT VERSION0.49 QUALIFICATION: Construct the adjoining simultaneous principal framed points and test-ring dictionary. Full finite holomorphic FamilyUnif/N1c, normalized theta families and local lifts are now supplied by HUF. The next bounded target is allTest\_\allowbreak of\_\allowbreak secBaseChange: represent FrPt by the Lie-eigenframe bundles and theta-section frame line, use the full n=1 SecBaseChangeSt on every nonflat coefficient base change, and prove finite-generated subalgebra descent with the exact TRing/FromTest transport. Then supply the full FramesOK/CoverTri/FibTri and IsUnf/hasPrinc package, including adapted mod-eight symplectic basis, normalized theta multiplier/weight and arithmetic invariance. Continue higher CoversAt, all-tuple class sums and uniqueness with original phase rescalings, one PointIdentCore with common lambda,z, base ClassSumsSpan, Gaussian/reproduction, short-rule dictionaries and ramified jets, retaining both-level, coefficient-field, Fits and eta agreements. Completed current producers or conditional consumers: ICL1-ICL21, CNC1-CNC16, PCL1-PCL12, PUN1-PUN24, HUF1-HUF21 (full original FamilyUnif/N1c and theta family). Earlier frontier descriptions are historical; this qualification is the current scope authority. Original whole-contract premises and quantifiers are retained.

\emph{Completed inputs (version 0.56).} PBC1–PBC27 (full arbitrary-base polarized pullback, actual very-ampleness and section finiteness), SHV1–SHV21 (both original Serre conclusions, actual finite section model and full all-cover return), ASL1–ASL21 (actual section localization and full original arbitrary-base Perf producer), FCD1-FCD18 (full homogeneous-chart descent, graded presentation and Noetherian cube model), GSL1-GSL26, NFL1-NFL22, TDR1-TDR13, NFR1-NFR13, HUF1-HUF21, PUN1-PUN24, ICL1-ICL21, CNC1-CNC16, PCL1-PCL12, Inherited both-level RatG/LiftTestB consumers

\emph{Remaining producer scope.} PBC supplies full arbitrary-base RelAmpleOfLocallyOnBase and literal PolBaseChangeSt, including the actual Cartesian square, unchanged dimension, actual pulled line bundle, generated comparison, rigidification, symmetry, all-test group multiplication and all-affine-test KTriv. The finite determinant trick, extra-section affineness and two-exponent section extension give the synchronized positive-power glue\_\allowbreak pack and exists\_\allowbreak sv without a Noetherian or flatness restriction on base change. PBC26 now supplies full polAV\_\allowbreak veryAmpleBy on every original affine test and specified Cartesian model; its power may depend on that test. PBC27 derives the finite-relation property from Perf over any field and proves original sections\_\allowbreak finite\_\allowbreak of\_\allowbreak perf and every-integer-power SecM finiteness over the original algebraically closed field. Its separate arbitrary-field finite\_\allowbreak secM implication retains the supplied finite-relation premise. GSL\_\allowbreak POL\_\allowbreak BASE\_\allowbreak CHANGE is supplied, alongside the already complete GSL\_\allowbreak SERRE and ASL Perf inputs. Remaining for GSL17 are multiplication-pullback generated comparison, finite/flat/surjective and etale multiplication geometry, and invariant-function descent on every relevant open. Geometric integrality, principal rank one, torsion-indexed section decomposition, two-sided growth, and square/translation generators remain for full P,V,R,T and the GSL26/NFL joins; section finiteness alone does not supply the growth half of GSL\_\allowbreak FINITE\_\allowbreak GROWTH. Original Moduli and HUF/PUN foundations remain. Later finite-level systems, Gaussian/reproduction, simultaneous dictionaries, ramified jets, all-tuple class sums, phase-sensitive uniqueness, common PointIdentCore and base ClassSumsSpan remain, retaining both levels, whole fields, Fits and eta agreement.



\paragraph{24. SLR\_\allowbreak BASE\_\allowbreak ROWS\_\allowbreak PRODUCER — C — Actual base row identities on that same bundle}

For the SAME selected Bundle, common original fields and readings used by the other three BTH12 inputs, supply actual RowsHold0\_\allowbreak 3 uniformly eventually for every family. Its original primitive row realization, spectrum/type and basis identifications must be retained. BTH supplies a theta consumer only; it does not construct the base row identity. No separate witness may be chosen merely to satisfy this row requirement.



\clearpage
\section{Relative ampleness on an arbitrary base and actual polarized base change}
\label{sec:polarized-base-change-v056}

\status{This chapter proves the full original
\texttt{rel\allowbreak Ample\_\allowbreak of\_\allowbreak locally} for arbitrary schemes and the full
\texttt{Pol\allowbreak Base\allowbreak Change\allowbreak St}.  The construction uses the specified Cartesian
group-scheme pullback, keeps the original relative dimension, and equips
the actual pulled line bundle with its rigidification, symmetry,
relative ampleness and translation-kernel property.  The immediate
perfect-complex consumer also gives the full actual-section
finiteness statement and finiteness of every integer-power section
module.  There is no
Noetherianity or flatness hypothesis on the change of base.
The previously proved Serre and perfect-complex constructions are
reused at their existing scopes; the remaining multiplication,
geometric-rank, square-law and moduli producers are not asserted here.}

The main difficulty is to turn locally supplied ample presentations into
one presentation over each affine base open.  Their powers of the line
bundle can differ.  Taking a common positive multiple resolves that
problem, but merely extending powers of the original sections would
not yet prove the required generation of each affine coordinate ring.
We therefore extend all the mixed products
$s_b s_a^d$.  A finite ideal argument proves that their ratios generate
the nonvanishing open of every section of the common degree.
The argument works over arbitrary rings, including rings with
nilpotents.

The first part uses stalk-valued functions to express the construction
without constantly changing trivializations.  The second part recovers
actual compatible regular functions on every original line-bundle
chart.  Thus the auxiliary functions never replace the actual sections
required by the final statement.

\paragraph{Source and notation.}
Source commit:
\texttt{7cbf7a0c32fca1badf05931879a37c3f8a2b03b4}.
The principal files are
\texttt{Exists\allowbreak One\allowbreak Add\allowbreak Kill.\allowbreak lean},
\texttt{Rel\allowbreak Ample\allowbreak Of\allowbreak Locally\allowbreak On\allowbreak Base.\allowbreak lean}, and
\texttt{Pol\allowbreak Base\allowbreak Change\allowbreak Of\allowbreak Rel\allowbreak Ample\allowbreak Local.\allowbreak lean} in the original
abelian-scheme-law proofs.  The source map at the end identifies their
declarations and the additional chart and base-change arguments.

For a scheme $Z$ put
\[
 \operatorname{Fn}(Z)=\prod_{z\in Z}\mathcal O_{Z,z}.
\]
No regularity is required of an element of this product.  For an open
$W$, let $\operatorname{Reg}(W)$ be the subring consisting of the
functions which, on $W$, are the germs of one actual section of
$\mathcal O(W)$.  Values outside $W$ are immaterial.  Write
$\operatorname{inv}(F)(z)$ for the inverse of $F(z)$ if it is a unit
and for zero otherwise.  This is the source's total
$\operatorname{Ring.inverse}$, so identities involving it are used
only where the indicated element is a unit.

The assertion $\operatorname{IsLoc}(D,W,F)$ means
\[
 z\in D\quad\Longleftrightarrow\quad
 z\in W\ \hbox{and}\ F(z)\in\mathcal O_{Z,z}^{\times}.
                                                        \tag{PBC.1}
\]
For $S\subseteq\operatorname{Fn}(Z)$, write
$\operatorname{Gen}(D,S)$ if every $F\in\operatorname{Reg}(D)$
agrees on $D$ with an element of the subring $\langle S\rangle$.
This definition does not require the generators to be regular away
from $D$, and it does not identify functions which differ away from $D$.

Fix a cover $(U_i)_{i\in I}$ of $Z$ and functions
$G_i\in\operatorname{Fn}(Z)$ which are units at every point.  Assume
\[
 G_i\,\operatorname{inv}(G_j)\in
       \operatorname{Reg}(U_i\cap U_j).
                                                        \tag{PBC.2}
\]
For $n\ge0$ define
\[
 \operatorname{Sc}(W,n,F)
 \quad\Longleftrightarrow\quad
 G_i^nF\in\operatorname{Reg}(W\cap U_i)
       \quad\hbox{for every }i.
                                                        \tag{PBC.3}
\]
The cover and the functions $G_i$ remain fixed whenever this notation
is used.  They need not arise from a line bundle until PBC14.

\subsection{Denominator clearing and the finite ideal argument}

% BEGIN PBC1
\paragraph{PBC1: regular germs, unit loci and ratios.}
\label{claim:PBC1}
Equality of all germs on an open implies equality of its sections.
Regularity of a stalk-valued function is local.  A regular function
which is a unit in every stalk on $W$ has a regular inverse on $W$.
Consequently a function satisfying $\operatorname{Sc}(W,n,F)$
has an open unit locus in $W$.  If $F,\sigma$ both satisfy this
condition and $\operatorname{IsLoc}(D,W,\sigma)$, then
$F\operatorname{inv}(\sigma)\in\operatorname{Reg}(D)$.

\begin{proof}
The first assertion is the separatedness of the structure sheaf.
For the second, choose local representatives.  On an overlap their
germs agree with the prescribed function, so the representatives agree;
the sheaf gluing axiom gives one representative.  If $a\in\mathcal O(W)$
has unit germs, each point has a neighborhood on which $a$ has an
inverse: represent the inverse germ and shrink until its product with
$a$ is one.  These inverses agree on overlaps and glue.

On $W\cap U_i$ the function $G_i^nF$ is regular.  Since $G_i$ is a
unit, its ordinary principal open is exactly the locus where $F$ is
a unit.  The union of these opens is the asserted open.  On
$D\cap U_i$, both $G_i^nF$ and $G_i^n\sigma$ are regular and the
second is a unit.  Their quotient is
\[
 \frac{G_i^nF}{G_i^n\sigma}=F\operatorname{inv}(\sigma).
\]
These local representatives agree at every germ, hence glue.
The same definitions also give restriction of $\operatorname{Sc}$,
addition in one degree, multiplication with degrees added, and
$\operatorname{Sc}(W,nd,F^d)$ from
$\operatorname{Sc}(W,n,F)$.
\end{proof}
% END PBC1

% BEGIN PBC2
\paragraph{PBC2: affine fractions and one exponent on a quasi-compact open.}
\label{claim:PBC2}
Suppose $\operatorname{IsLoc}(D,W,w)$ and
$w\in\operatorname{Reg}(W)$.  If $W$ is affine and
$F\in\operatorname{Reg}(D)$, there are $a\ge0$ and
$H\in\operatorname{Reg}(W)$ such that
$w^aF=H$ on $D$.  If $W$ is only quasi-compact,
$F\in\operatorname{Reg}(W)$ and $F=0$ on $D$, there is $b\ge0$
such that $w^bF=0$ on all of $W$.

\begin{proof}
Represent $w$ by $r\in\mathcal O(W)$.  When $W=\operatorname{Spec}A$,
the open $D$ is $D(r)$ and
$\mathcal O(D)=A[r^{-1}]$.  Write the representative of $F$ as
$h/r^a$ and use the germs of $h$ on $W$, with arbitrary values
outside $W$, for $H$.

For the second assertion, choose a finite affine cover $V_\alpha$
of $W$.  The image of $F|_{V_\alpha}$ is zero after inverting
$w|_{V_\alpha}$, so some power of that element kills it in
$\mathcal O(V_\alpha)$.  The maximum of the finitely many exponents
works on every $V_\alpha$.  Its product with $F$ is therefore zero
on $W$.  Only the elementary equality criterion in a localization
is used; $w$ may be a zero divisor.
\end{proof}
% END PBC2

% BEGIN PBC3
\paragraph{PBC3: extending a prescribed-degree section across a principal base open.}
\label{claim:PBC3}
Let $Z$ be quasi-compact and quasi-separated, with the cover and
unit functions in (PBC.2).  Suppose $f\in\operatorname{Reg}(Z)$,
$\operatorname{IsLoc}(W,Z,f)$, and
$\operatorname{Sc}(W,n,F)$.  Then there are $k\ge0$ and
$\sigma\in\operatorname{Fn}(Z)$ such that
\[
 \operatorname{Sc}(Z,n,\sigma),\qquad
 \sigma=f^kF\quad\hbox{on }W.
                                                        \tag{PBC.4}
\]
The degree $n$ is unchanged.  This is the full original
$\operatorname{Sc.extend}$.

\begin{proof}
Choose finitely many affine opens $V_\alpha$ covering $Z$, each
contained in some $U_{i(\alpha)}$.  On
$W\cap V_\alpha$ the regular function
$G_{i(\alpha)}^nF$ has an affine fraction representation.  Thus
PBC2 gives $p_\alpha$ and $H_\alpha\in\operatorname{Reg}(V_\alpha)$
such that
\[
 H_\alpha=f^{p_\alpha}G_{i(\alpha)}^nF
       \quad\hbox{on }W\cap V_\alpha.
\]
Choose $p\ge p_\alpha$ for every $\alpha$ and put, as a function of
stalks,
\[
 t_\alpha=f^{p-p_\alpha}H_\alpha
                    \operatorname{inv}(G_{i(\alpha)})^n.
\]
Then $t_\alpha=f^pF$ on $W\cap V_\alpha$.  Although $t_\alpha$
itself need not be regular, for every $i$ the identity
\[
 G_i^nt_\alpha
  =(G_i/G_{i(\alpha)})^n f^{p-p_\alpha}H_\alpha
                                                        \tag{PBC.5}
\]
shows that it is regular in the $i$th gauge on $U_i\cap V_\alpha$.

The intersection $V_\alpha\cap V_\beta$ is quasi-compact by
quasi-separatedness.  There the regular difference
$G_{i(\alpha)}^n(t_\alpha-t_\beta)$ vanishes after restricting to
$W$.  PBC2 gives an exponent $q_{\alpha\beta}$ killing it on the
whole intersection.  Cancel the unit $G_{i(\alpha)}^n$.
Choose one $q$ at least all these finitely many exponents.  Then
$f^qt_\alpha=f^qt_\beta$ on every overlap.

For each $z$ choose $\alpha(z)$ containing $z$ and define
$\sigma(z)=f(z)^qt_{\alpha(z)}(z)$.  On $V_\alpha$ this agrees
with $f^qt_\alpha$ regardless of the pointwise choice.  Formula
(PBC.5), followed by PBC1, proves
$G_i^n\sigma\in\operatorname{Reg}(U_i)$ for every $i$.
On $W$ the definition gives $\sigma=f^{p+q}F$.  Thus $k=p+q$
works.  The first exponent clears fractions; the second makes the
different local extensions agree.  Neither is discarded.
\end{proof}
% END PBC3

% BEGIN PBC4
\paragraph{PBC4: the principal-open affineness criterion.}
\label{claim:PBC4}
Let $W$ be an open in a scheme and let $r_\lambda\in\mathcal O(W)$
generate the unit ideal.  If each $D_W(r_\lambda)$ is affine,
then $W$ is affine.

\begin{proof}
A finite subfamily already generates the unit ideal, because an
expression for $1$ is a finite sum.  Its principal opens cover $W$.
They are quasi-compact.  Their pairwise intersections are principal
opens in either affine member, hence affine and quasi-compact.
It follows that $W$ is quasi-compact and quasi-separated.

Put $A=\mathcal O(W)$.  For each $r=r_\lambda$, restriction induces
\[
 A[r^{-1}]\longrightarrow\mathcal O(D_W(r)).
                                                        \tag{PBC.6}
\]
It is injective by PBC2: a section restricting to zero is killed by a
power of $r$.  It is surjective by PBC3 in degree zero with the
trivial gauge $G=1$: a section on $D_W(r)$ becomes the restriction
of a global section after multiplication by a power of $r$.
Thus (PBC.6) is an isomorphism.

The canonical map $W\to\operatorname{Spec}A$ therefore restricts
to an isomorphism $D_W(r_\lambda)\simeq
D_{\operatorname{Spec}A}(r_\lambda)$ for the finite chosen family.
These target opens cover because the $r_\lambda$ generate $1$.
The restrictions agree on overlaps, being induced by the same
global ring map.  Their inverses glue to an inverse of the
canonical map.  Hence $W$ is affine.
\end{proof}
% END PBC4

% BEGIN PBC5
\paragraph{PBC5: generated functions, localization, and the exact polynomial map.}
\label{claim:PBC5}
The following three operations preserve the source's
$\operatorname{Gen}$ assertion.
First, generators may be replaced by elements agreeing on the
relevant open with polynomials in new generators.
Second, if $W$ is affine,
$\operatorname{Gen}(W,S)$ and
$\operatorname{IsLoc}(D,W,w)$ with $w\in\operatorname{Reg}(W)$,
then
\[
 \operatorname{Gen}(D,S\cup\{\operatorname{inv}(w)\}).
                                                        \tag{PBC.7}
\]
Third, let $\psi:R\to\mathcal O(D)$ be a ring map and
$r_j\in\mathcal O(D)$ an arbitrary family.  If
$\phi(a),\rho_j\in\operatorname{Fn}(Z)$ have on $D$ the germs
of $\psi(a),r_j$, respectively, then
\[
 \operatorname{Gen}\bigl(D,\phi(R)\cup\{\rho_j\}\bigr)
 \quad\Longleftrightarrow\quad
 R[X_j]\longrightarrow\mathcal O(D),\
 X_j\longmapsto r_j
 \ \hbox{is surjective}.
                                                        \tag{PBC.8}
\]

\begin{proof}
An element of a generated subring is obtained by finitely many
additions, negations and multiplications from the generators and
$0,1$.  Substituting the asserted new expressions proves the first
operation by induction on such an expression.

For (PBC.7), write $w^aF=H$ on $D$ using PBC2, then choose
$b\in\langle S\rangle$ agreeing with $H$ on $W$.
On $D$ the function $F$ agrees with
$\operatorname{inv}(w)^a b$.

For (PBC.8), an expression in $\phi(R)$ and the $\rho_j$ becomes a
polynomial over $R$ by the same induction; its evaluation has exactly
the indicated germs.  Applied to the germs of any section of
$\mathcal O(D)$, equality of germs proves surjectivity.
Conversely, a polynomial representative of the section representing
$F\in\operatorname{Reg}(D)$ supplies the element of the generated
subring required by $\operatorname{Gen}$.  There is no finiteness,
affineness, or Noetherianity hypothesis in (PBC.8).  The actual
coefficient map $\psi$ is retained.
\end{proof}
% END PBC5

% BEGIN PBC6
\paragraph{PBC6: local affine generation gives global affine generation.}
\label{claim:PBC6}
Let $B\subseteq\operatorname{Fn}(Z)$ be a subring with
$B\subseteq\operatorname{Reg}(W)$.  Let
\[
 K=\{b\in B:b(z)=0\text{ for every }z\in W\}.
                                                        \tag{PBC.9}
\]
Suppose $w_\lambda\in B$ satisfy $K+(w_\lambda)=B$.
For each $\lambda$ let
$\operatorname{IsLoc}(O_\lambda,W,w_\lambda)$, assume
$O_\lambda$ affine, and assume
$\operatorname{Gen}(O_\lambda,B\cup\{\operatorname{inv}(w_\lambda)\})$.
Then $W$ is affine and $\operatorname{Gen}(W,B)$.
The family of $\lambda$ may be infinite.

\begin{proof}
Assign to $b\in B$ the unique actual section with its germs on
$W$.  PBC1 shows that this defines a ring map
$\rho:B\to\mathcal O(W)$ with kernel $K$.
The images $\rho(w_\lambda)$ generate the unit ideal.  The
$O_\lambda$ are their principal opens, so PBC4 proves affineness
of $W$.

Fix $F\in\operatorname{Reg}(W)$ and define the ideal
\[
 J_F=\{b\in B:\text{ there is }b'\in B
                  \text{ with }bF=b'\text{ on }W\}.
                                                        \tag{PBC.10}
\]
It contains $K$.  On $O_\lambda$, local generation expresses
$F$ as a polynomial in elements of $B$ and $w_\lambda^{-1}$.
Clearing the finitely many denominators gives
$w_\lambda^aF=b'$ there, with $b'\in B$.
The difference is regular on $W$.  As $W$ is now affine and in
particular quasi-compact, PBC2 gives $b\ge0$ with
$w_\lambda^b(w_\lambda^aF-b')=0$ on $W$.
Hence $w_\lambda^{a+b}\in J_F$.

Thus $\sqrt{J_F}$ contains $K$ and all $w_\lambda$; it is the
unit ideal.  Consequently $J_F$ itself contains $1$.
The definition of $J_F$ says that $F$ agrees on $W$ with an
element of $B$.  This is the required global generation.
\end{proof}
% END PBC6

% BEGIN PBC7
\paragraph{PBC7: the finite one-plus-an-ideal annihilator.}
\label{claim:PBC7}
Let $B$ be any commutative ring, $u_t$ a finite family,
$I=(u_t)$, and $K$ any ideal.  Suppose that for each $t$ there is
$e_t\ge0$ such that
\[
 u_t^{e_t}\in K+I^{e_t+1}.
                                                        \tag{PBC.11}
\]
There are $c\in B$ and $d\ge0$ with
\[
 c-1\in I,\qquad c u_t^d\in K\quad\hbox{for every }t.
                                                        \tag{PBC.12}
\]
This is the full original \texttt{exists\_\allowbreak one\_\allowbreak add\_\allowbreak kill}.

\begin{proof}
If the family is empty, take $c=1,d=0$.  If some $e_t=0$,
(PBC.11) says $1\in K+I$.  Write $1=c+i$ with $c\in K,i\in I$
and take $d=0$.
Otherwise every $e_t\ge1$.  Put
$D=1+\sum_t(e_t-1)$.  Every monomial of degree $D$ in the
$u_t$ is divisible by some $u_t^{e_t}$; multiplying its
(PBC.11) relation by the remaining factors gives
\[
 I^D\subseteq K+I^{D+1}.
                                                        \tag{PBC.13}
\]
Consider the finitely generated $B$-module
$M=(I^D+K)/K$, generated by the degree-$D$ monomials.
Equation (PBC.13) says $M=IM$.
Choose generators $m_1,\ldots,m_r$ and write
$m_i=\sum_j a_{ij}m_j$ with every $a_{ij}\in I$.
For the matrix $A=(a_{ij})$, the adjugate identity applied to
$(1-A)(m_j)_j=0$ shows that
$c=\det(1-A)$ kills every generator, hence all of $M$.
Expanding the determinant gives $c-1\in I$.
Thus $cI^D\subseteq K$ and $d=D$ works.
If $M=0$, the same conclusion follows from the empty determinant
$c=1$.  This is a finite-matrix argument, so no Noetherian
assumption on $B$ or finite generation of $K$ is involved.
\end{proof}
% END PBC7

% BEGIN PBC8
\paragraph{PBC8: mixed ratios and the extra-section ideal estimates.}
\label{claim:PBC8}
Let $T$ be finite and suppose sections in the stalk notation
$s_j$ have degree $N$ on $W_0$.  Their unit loci $D_j$ cover
$W_0$, are affine, and satisfy
\[
 \operatorname{Gen}\bigl(D_j,S_0\cup\{s_m/s_j:m\in T\}\bigr),
 \qquad S_0\subseteq\operatorname{Reg}(W_0).
                                                        \tag{PBC.14}
\]
Let $\operatorname{Sc}(W_0,N(d+1),\tau)$ with $d\ge0$, and
let $W$ be its unit locus in $W_0$.  Put
\[
 u_{ij}=s_j s_i^d/\tau,\quad
 B=\langle S_0\cup\{u_{ij}:i,j\in T\}\rangle,\quad
 I=(u_{ij})\subseteq B,
                                                        \tag{PBC.15}
\]
where quotients mean the prescribed total inverses.
Then $B\subseteq\operatorname{Reg}(W)$, the opens
$D_j\cap W$ are affine, and, for $K$ as in (PBC.9),
each $u_{ij}$ satisfies the hypothesis (PBC.11).

\begin{proof}
The numerator in $u_{ij}$ has the same degree as $\tau$, so
PBC1 makes $u_{ij}$ regular on $W$.  On $D_j$ the function
\[
 v_j=\tau/s_j^{d+1}
\]
is regular.  Its unit locus is $D_j\cap W$, hence this
intersection is a principal open in the affine $D_j$.
On $D_j\cap W$ one has
\[
 v_ju_{jj}=1,\qquad
 u_{jj}(s_m/s_j)=u_{jm}.
                                                        \tag{PBC.16}
\]
Since the finite family $D_j\cap W$ covers $W$, this last open
is quasi-compact.

By (PBC.14), $v_j$ agrees on $D_j$ with an expression $b_j$
in $S_0$ and the ratios $s_m/s_j$.  Induction on this expression,
using the second identity in (PBC.16), gives an integer $e\ge0$
and $x\in I^e$ with
\[
 u_{jj}^e b_j=x\quad\hbox{on }D_j\cap W.
                                                        \tag{PBC.17}
\]
Here a generator from $S_0$ uses exponent zero; a ratio uses
exponent one and $u_{jm}\in I$; products add exponents.
For a sum with exponents $a,b$, multiply its two numerators
by $u_{jj}^b,u_{jj}^a$, respectively.  Both new numerators
belong to $I^{a+b}$.  This proves (PBC.17), including sums
with coefficients and negative signs.

The first identity in (PBC.16) gives
$u_{jj}^e-u_{jj}x=0$ on $D_j\cap W$.
This intersection is also the unit locus of $u_{jj}$ in $W$.
PBC2, applied to the regular difference on the quasi-compact
open $W$, gives $m\ge0$ such that
\[
 u_{jj}^{m+e}
       \equiv u_{jj}^{m+1}x \pmod K.
                                                        \tag{PBC.18}
\]
The right side belongs to $I^{m+e+1}$.  Thus (PBC.11)
holds for each diagonal $u_{jj}$.

For all $i,j$, the pointwise identity
\[
 u_{ij}^{d+1}=u_{jj}u_{ii}^{d}
                                                        \tag{PBC.19}
\]
holds even outside $W$: it is a polynomial identity in the
functions $s_i,s_j,\operatorname{inv}(\tau)$.
If the diagonal relation uses exponent $e$, raising
(PBC.19) to $e$ and multiplying that relation by $u_{ii}^{de}$
shows
$u_{ij}^{(d+1)e}\in K+I^{(d+1)e+1}$.
This proves the assertion for every mixed ratio.
\end{proof}
% END PBC8

% BEGIN PBC9
\paragraph{PBC9: every section of a higher common degree has an affine generated unit locus.}
\label{claim:PBC9}
Under the full hypotheses of PBC8,
\[
 W\text{ is affine},\qquad
 \operatorname{Gen}\bigl(W,S_0\cup\{s_j s_i^d/\tau:i,j\in T\}\bigr).
                                                        \tag{PBC.20}
\]
This proves the original \texttt{extra\_\allowbreak section} for every
$d\ge0$ and every such $\tau$.

\begin{proof}
Apply PBC7 to $I,K$ and the mixed ratios.  We obtain
$c-1\in I$ and one $q$ with $c u_{ij}^q\in K$ for every pair.
At a point of $W$, choose $j$ with the point in $D_j$.
Then $u_{jj}$ is a unit by (PBC.16).  Evaluating
$c u_{jj}^q=0$ at that point shows $c=0$ there.
Thus $c\in K$.

Let $J=K+(u_{jj}:j\in T)$.  Formula (PBC.19) implies
$I\subseteq\sqrt J$.  Since $c\in J$ and $c-1\in I$,
we have $1\in\sqrt J$, whence $J=B$.

We can now apply PBC6 with the diagonal functions $u_{jj}$.
Their unit loci in $W$ are the already affine opens
$D_j\cap W$.  To check its generation condition, first localize
(PBC.14) at $v_j$ using PBC5.  On this open,
\[
 v_j^{-1}=u_{jj},\qquad
 s_m/s_j=u_{jj}^{-1}u_{jm}.
                                                        \tag{PBC.21}
\]
The localized generators are therefore expressions in $B$ and
$u_{jj}^{-1}$.  PBC6 proves that $W$ is affine and generated
by $B$.  Finally, substituting the defining generators of $B$
gives precisely (PBC.20).
If $T$ is empty the cover forces $W_0$, and hence $W$, to be
empty; the same argument uses the zero coordinate ring and remains
valid.  No nonemptiness condition has been added.
\end{proof}
% END PBC9

\subsection{One positive power and one global ample pack}

A \emph{pack} on $W$ with base generators $S$ consists of a positive
integer $N$, a nonempty finite tuple $s_0,\ldots,s_m$ satisfying
$\operatorname{Sc}(W,N,s_j)$, and their unit loci $D_j$ in $W$,
such that these opens cover $W$, are affine, and satisfy
$\operatorname{Gen}(D_j,S\cup\{s_b/s_j\}_b)$.
This is exactly the source's $\operatorname{Pack}(U,G,W,S)$.

% BEGIN PBC10
\paragraph{PBC10: synchronizing all local degrees before extension.}
\label{claim:PBC10}
Let $Z$ be quasi-compact and quasi-separated, with (PBC.2).
Let $\phi:R\to\operatorname{Fn}(Z)$ be a ring map whose values
are regular on $Z$.  Let $(a_\ell)_{\ell\in\Lambda}$ be a finite
family generating the unit ideal of $R$, and let $W_\ell$ be
the unit locus of $\phi(a_\ell)$ in $Z$.
Suppose a pack is given on each $W_\ell$ with base generators
\[
 S_\ell=\phi(R)\cup\{\operatorname{inv}(\phi(a_\ell))\}.
                                                        \tag{PBC.22}
\]
There is a single $N\ge1$ and, for each
$t=(\ell,a,b)$ with $0\le a,b\le m_\ell$, a global
degree-$N$ function $\sigma_t$, together with $k_t\ge0$, such that
\[
 N=N_\ell(d_\ell+1),\qquad
 \sigma_{\ell,a,b}
    =\phi(a_\ell)^{k_{\ell,a,b}}\,
       s_{\ell,b}s_{\ell,a}^{d_\ell}
       \quad\hbox{on }W_\ell .
                                                        \tag{PBC.23}
\]
The choices of $N,\sigma_t,k_t$ are made before any global
degree-$N$ section later used as a denominator.

\begin{proof}
Choose all the finitely many local pack data and put
$N=\prod_{\ell\in\Lambda}N_\ell$.  This is positive, with
the empty product equal to $1$.  Each $N_\ell$ divides $N$
with positive quotient, which we write $d_\ell+1$.
Multiplication and powers in (PBC.3) show that
$s_{\ell,b}s_{\ell,a}^{d_\ell}$ has degree $N$ on
$W_\ell$.  Apply PBC3 to this function with
$f=\phi(a_\ell)$.  It gives exactly (PBC.23).
There are only finitely many triples, but no common bound on their
extension exponents is needed.

The $W_\ell$ cover $Z$.  Indeed, if at a point all
$\phi(a_\ell)(z)$ were nonunits, they would all belong to the
maximal ideal of the local ring $\mathcal O_{Z,z}$.
The image of the equation $1\in(a_\ell)$ would put $1$ in that
maximal ideal, a contradiction.
\end{proof}
% END PBC10

% BEGIN PBC11
\paragraph{PBC11: affine generation for every denominator of the common degree.}
\label{claim:PBC11}
Keep the data chosen in PBC10.  For every $\tau$ satisfying
$\operatorname{Sc}(Z,N,\tau)$, and every open $E$ with
$\operatorname{IsLoc}(E,Z,\tau)$, one has
\[
 E\text{ affine},\qquad
 \operatorname{Gen}\bigl(E,\phi(R)\cup\{\sigma_t/\tau\}_t\bigr).
                                                        \tag{PBC.24}
\]

\begin{proof}
Let $B=\langle\phi(R)\cup\{\sigma_t/\tau\}_t\rangle$.
The ratio assertion of PBC1 gives $B\subseteq\operatorname{Reg}(E)$.
There is a ring map $\psi:R\to B$ induced by $\phi$.
The elements $\psi(a_\ell)$ generate the unit ideal of $B$.

Apply PBC9 on $W_\ell$ to its original pack and the restriction
of $\tau$, using $N=N_\ell(d_\ell+1)$.
The set $S_\ell$ is regular there: the only additional generator
is the inverse of the unit $\phi(a_\ell)$.
It follows that $W_\ell\cap E$ is affine and generated by
$S_\ell$ and the functions
$s_{\ell,b}s_{\ell,a}^{d_\ell}/\tau$.
On this intersection (PBC.23) gives the exact identity
\[
 \frac{s_{\ell,b}s_{\ell,a}^{d_\ell}}{\tau}
   =\operatorname{inv}(\phi(a_\ell))^{k_{\ell,a,b}}
               \frac{\sigma_{\ell,a,b}}{\tau}.
                                                        \tag{PBC.25}
\]
Consequently these generators are expressions in $B$ and
$\operatorname{inv}(\psi(a_\ell))$.  The open
$W_\ell\cap E$ is the unit locus of $\psi(a_\ell)$ in $E$.
PBC6 now proves that $E$ is affine and generated by $B$.
PBC5 substitutes the generators of $B$ and gives (PBC.24).
This argument applies to every $\tau$ after the numerator family
has been fixed.
\end{proof}
% END PBC11

% BEGIN PBC12
\paragraph{PBC12: the full original global pack constructor.}
\label{claim:PBC12}
Under the original hypotheses of PBC10, there are $N\ge1$,
$M\ge0$, and $\sigma_0,\ldots,\sigma_M$ such that
$\operatorname{Sc}(Z,N,\sigma_t)$ for every $t$, at every point
some $\sigma_t$ is a unit, and for every $t$ and every open $D$
\[
 \operatorname{IsLoc}(D,Z,\sigma_t)
 \ \Longrightarrow\
 D\text{ affine and }
 \operatorname{Gen}\bigl(D,\phi(R)\cup\{\sigma_m/\sigma_t\}_m\bigr).
                                                        \tag{PBC.26}
\]
This is the full original \texttt{glue\_\allowbreak pack}, including its
universal final open $D$.

\begin{proof}
Adjoin one zero function to the finite family in PBC10 and index
the resulting nonempty family by
$\{0,\ldots,M\}$, where $M$ is the number of triples
$(\ell,a,b)$.  Zero satisfies $\operatorname{Sc}$ in every degree.
At a point $z$, choose $\ell$ with $z\in W_\ell$ and then a local
section $s_{\ell,j}$ which is a unit at $z$.  The diagonal case
of (PBC.23) gives
\[
 \sigma_{\ell,j,j}(z)
   =\phi(a_\ell)(z)^{k_{\ell,j,j}}
                     s_{\ell,j}(z)^{d_\ell+1},
\]
which is a unit.  Thus the unit loci cover.

For the last assertion apply PBC11 with $\tau=\sigma_t$.
All its numerator functions occur in the enlarged family,
so PBC5 gives the displayed generation statement.
If the initial finite index set is empty, the unit-ideal
hypothesis and the local-ring argument of PBC10 force $Z$ to be
empty; the appended zero still gives the required nonempty tuple.
This accounts for the source's $\operatorname{Fin}(M+1)$
without imposing an extra nonemptiness hypothesis.
\end{proof}
% END PBC12

\subsection{Returning the pack to the actual line-bundle charts}

% BEGIN PBC13
\paragraph{PBC13: transport through an open immersion.}
\label{claim:PBC13}
For a scheme map $h:Y\to Z$, define
\[
 (h^\flat F)(y)=h_y^\#(F(h(y))).
                                                        \tag{PBC.27}
\]
This is a ring map, composes as pullback, preserves and reflects
the property of being a unit at the corresponding stalk, and
commutes with the total inverse.  It preserves regularity by
pulling actual sections.  If $h$ is an open immersion, every
function in $\operatorname{Fn}(Y)$ has a lift in
$\operatorname{Fn}(Z)$, and regularity and generation on an open
of $Y$ transport to its actual image in $Z$.
In particular, a pack for the pulled cover, pulled gauges and
pulled generators on $Y$ gives a pack on $h(Y)$ for the original
cover, gauges and generators.

\begin{proof}
The stalk maps of a scheme morphism are local homomorphisms.
A local homomorphism preserves units and takes nonunits to nonunits.
Therefore it commutes with inverse in both cases of the definition.
The ring and composition assertions follow from the stalk maps,
and a pulled section has exactly the germs in (PBC.27).

For an open immersion, $h$ identifies $Y$ with its image and
every stalk map is an isomorphism.  Given $F$ on $Y$, use the
inverse stalk maps to define its lift on $h(Y)$ and define the
lift to be zero elsewhere.  This is a lift as a function of
stalks; it is not claimed to be regular across the boundary.
Regularity on an open $V\subseteq Y$ transports because the
open-immersion section map identifies $\mathcal O(h(V))$ with
$\mathcal O(V)$.  To transport $\operatorname{Gen}$, pull a
regular function back, use its generating expression there, and
use the injective stalk maps to obtain the same expression on
$h(V)$.

Now lift each of the finitely many functions in the given pack.
On $h(Y)\cap U_i$, the regularity of its $i$th gauge follows
from regularity on $Y\cap h^{-1}U_i$ and the preceding section
identification.  Unit loci are images of the old unit loci,
so they cover and are affine.  Pullback commutes with their
ratios, and the generation assertion just proved transports
the old ratio generators to the lifted ones.  These are all
fields of the asserted pack.
\end{proof}
% END PBC13

% BEGIN PBC14
\paragraph{PBC14: chosen-chart values and the literal coordinate-ratio identity.}
\label{claim:PBC14}
Let $P$ be an original line-bundle chart datum on $Y$.
Its transition convention for a section is
$s_k=g_{kk'}s_{k'}$ on $P.U_k\cap P.U_{k'}$.
Choose a chart $\gamma(y)$ containing each $y$, and write
\[
 \operatorname{sv}_\gamma(s)(y)
       =\operatorname{germ}_y(s_{\gamma(y)}).
                                                        \tag{PBC.28}
\]
For every actual section $s$, its original nonzero locus is the
unit locus of $\operatorname{sv}_\gamma(s)$.  For every actual
section $t$ and $r\in\mathcal O(D(s))$, the identities
\[
 t_k|_{D(s)\cap P.U_k}
   =r|_{D(s)\cap P.U_k}\,s_k|_{D(s)\cap P.U_k}
       \quad\hbox{for every }k
                                                        \tag{PBC.29}
\]
are equivalent to
$\operatorname{germ}_y(r)=
\operatorname{sv}_\gamma(t)(y)/
\operatorname{sv}_\gamma(s)(y)$ for every $y\in D(s)$.

\begin{proof}
At any point in two charts the two coordinates of $s$ differ
by the unit germ of $g_{kk'}$.  One is a unit if and only if
the other is.  Taking the union of their ordinary principal
opens is therefore the unit locus in (PBC.28), exactly the
source's $\operatorname{nonzeroLocus}$.

Equation (PBC.29) in the chosen chart gives the asserted ratio
after division by the unit coordinate of $s$.
Conversely, multiply the ratio equality by that coordinate,
then multiply both sides by the same transition unit to obtain
the equality in any other chart $k$ containing the point.
Equality of all germs proves (PBC.29).
Only the given section compatibility and the chosen chart cover
are needed for these two assertions.

More generally, let $U_i\subseteq P.U_{\kappa(i)}$ and choose
$c(y)$ with $y\in U_{c(y)}$.  Suppose on $U_i$ that
\[
 \operatorname{germ}_y g_{\kappa(i),\kappa(c(y))}
      =G_i(y)^N.
                                                        \tag{PBC.30}
\]
Then every actual section $s$ gives
$\operatorname{Sc}(Y,N,\operatorname{sv}_{\kappa c}(s))$:
its $i$th gauged function is exactly the germ of
$s_{\kappa(i)}|_{U_i}$.  This is the forward
$\operatorname{sc\_\allowbreak sv}$ dictionary.
\end{proof}
% END PBC14

% BEGIN PBC15
\paragraph{PBC15: reconstructing an actual section from every gauged function.}
\label{claim:PBC15}
Assume $P.\operatorname{IsLB}$ and the precise data
$U_i,\kappa,c,G_i,N$ of (PBC.30).  For this assertion the
$G_i$ can be arbitrary stalk functions: no additional
global-unit or regular-ratio hypothesis is needed.
Every $\sigma$ with
$G_i^N\sigma\in\operatorname{Reg}(U_i)$ for all $i$ is
$\operatorname{sv}_{\kappa c}(s)$ for an actual section
$s$ on every original chart of $P$.
This is the full original \texttt{exists\_\allowbreak sv}.

\begin{proof}
For each original chart $k$, prescribe on $P.U_k$ the stalk function
\[
 F_k(y)=
   \operatorname{germ}_y g_{k,\kappa(c(y))}\,\sigma(y).
                                                        \tag{PBC.31}
\]
Near a fixed $z\in P.U_k$, use the fixed index $i=c(z)$.
On $P.U_k\cap U_i$, the cocycle identity and (PBC.30) give
\[
 F_k=
   \operatorname{germ}(g_{k,\kappa(i)})\,G_i^N\sigma.
                                                        \tag{PBC.32}
\]
The first factor is the germ of an actual transition unit on
this open, and the second is regular there by hypothesis.
Thus $F_k$ is locally regular at every point of $P.U_k$.
PBC1 gives one actual $s_k\in\mathcal O(P.U_k)$ with these germs.

On $P.U_k\cap P.U_{k'}$, the cocycle equation
$g_{k,\kappa(c(y))}=g_{kk'}g_{k',\kappa(c(y))}$
shows that $s_k=g_{kk'}s_{k'}$ at every germ, hence as
actual sections.  The family $(s_k)$ is therefore an
$\operatorname{LB.Sec}(P)$ in the original coordinates.
At $k=\kappa(c(y))$ the transition is $1$, so (PBC.31)
gives $\operatorname{sv}_{\kappa c}(s)=\sigma$.
The argument glues sections on their actual algebraic opens;
it does not assign a geometric regularity to the pointwise
choice $c$.
\end{proof}
% END PBC15

% BEGIN PBC16
\paragraph{PBC16: the exact base constants on a principal neighborhood.}
\label{claim:PBC16}
Let $B$ be affine, $a\in\mathcal O(B)$, and let the open
immersion $b:V\to B$ have image $D_B(a)$.  Suppose
$q:Y\to V$, $h:Y\to Z$, and $p:Z\to B$ satisfy
$p\circ h=b\circ q$.  Then every pulled base function
$q^*c$, $c\in\mathcal O(V)$, is, as a function of all
stalks of $Y$, an expression in
\[
 h^\flat(p^*\mathcal O(B))
       \quad\hbox{and}\quad
 h^\flat\bigl(\operatorname{inv}(p^*a)\bigr).
                                                        \tag{PBC.33}
\]

\begin{proof}
The actual open immersion identifies
$\mathcal O(V)$ with $\mathcal O(B)[a^{-1}]$.
Choose $r\in\mathcal O(B)$ and $n\ge0$ representing $c=r/a^n$.
Pullback gives
$q^*c=(p\circ h)^*r\,((p\circ h)^*a)^{-n}$.
The denominator is a unit everywhere on $Y$, since $q$ factors
through $D_B(a)$.  PBC13 therefore identifies its inverse with
the total inverse appearing in (PBC.33).
This proves literal membership in the generated subring of
$\operatorname{Fn}(Y)$, not just agreement on a smaller open.
\end{proof}
% END PBC16

% BEGIN PBC17
\paragraph{PBC17: a local relative-ample presentation gives the required principal pack.}
\label{claim:PBC17}
Let $A\to S$, an original chart datum $L$, an affine open
$U\subseteq S$, and an open $U_0\subseteq S$ be given.
Suppose the actual restriction of $L$ to $A_{U_0}$ is
relatively ample in the original sense.
Let $a\in\mathcal O(U)$ satisfy
$D_U(a)\subseteq U\cap U_0$.
Put $Z=A_U$, $p:Z\to U$, $j:Z\to A$, and choose $c(z)$ with
$j(z)\in L.U_{c(z)}$.
Set
\[
 U_i=j^{-1}(L.U_i),\qquad
 G_i(z)=\operatorname{germ}_z(j^*g_{i,c(z)})
                 \quad(z\in U_i),
                                                        \tag{PBC.34}
\]
using the source's value $1$ for this unit expression when the
indicated overlap condition fails.
There is a pack on $p^{-1}D_U(a)$ in these gauges with base
generators $p^*\mathcal O(U)\cup\{\operatorname{inv}(p^*a)\}$.
No additional global ampleness hypothesis is used.

\begin{proof}
View $D_U(a)$ also as the open $V$ of $U_0$ given by its image
in $S$.  This is an affine open, and its map $b:V\to U$
is the specified identification with $D_U(a)$.
Let $Y=(A_{U_0})_V$ with projection $q:Y\to V$.
There is a unique open immersion $h:Y\to A_U$ commuting with
the maps to $A$; its image is exactly $p^{-1}D_U(a)$ and
$p\circ h=b\circ q$.  These statements follow directly by
restricting the common point and structure-sheaf maps to the
same open of $A$.  No new model of $A$ is chosen.

Apply the given local relative ampleness on the affine open
$V\subseteq U_0$.  It supplies $N\ge1$, sections
$s_0,\ldots,s_m$ of the original repeated-factor power on
$Y$, affine nonzero loci covering $Y$, and ratio functions
whose polynomial evaluation over $\mathcal O(V)$ is surjective.
For this power use constant chart tuples: the same original
$L$-chart is used in each of the $N$ factors.
Its transition on two such tuples is the $N$th power of
the original transition, because its definition is the product
of those $N$ identical factors.  Thus (PBC.30) holds for the
pulled gauges $h^\flat G_i$ and the choice $c(h(y))$.

PBC14 sends each actual $s_j$ to its chosen-chart value,
with exactly its old nonzero locus.  The coordinate-ratio
identities and PBC5 turn the supplied polynomial surjections
into $\operatorname{Gen}$ with the base generators
$q^*\mathcal O(V)$.
PBC16 expresses every one of these base generators in the
pulled family
$h^\flat(p^*\mathcal O(U)\cup\{\operatorname{inv}(p^*a)\})$.
Replacing generators by PBC5 gives a pack with this latter
base family.  Finally PBC13 transports it through $h$.

For a point $z\in Z$, the local stalk map of $p$ reflects
units.  Hence $p^*a$ is a unit at $z$ precisely when
$p(z)\in D_U(a)$.  The pack open just constructed is
therefore the literal unit locus required in the source.
The forward dictionary used here requires only the supplied
sections and chart choice; the assertion has the original
\texttt{pack\_\allowbreak of\_\allowbreak rel\allowbreak Ample} scope.
\end{proof}
% END PBC17

% BEGIN PBC18
\paragraph{PBC18: full relative ampleness from neighborhoods of every base point.}
\label{claim:PBC18}
Let $S$ be an arbitrary scheme, $A\to S$ a quasi-compact,
quasi-separated morphism, and $L$ an actual line bundle on
$A$.  Suppose that for every $s\in S$ there is an open
$U_0$ containing $s$ such that the actual restriction
$L|_{A_{U_0}}$ is relatively ample.  Then
$\operatorname{RelAmple}(A,L)$ holds with the exact source
definition: for every affine open $U\subseteq S$ there are
$N\ge1$ and $m+1$ actual sections of
$\operatorname{powOver}(A,L,U,N)$ whose nonzero loci cover,
are affine, and have the stated coordinate-ratio polynomial
surjections over $\mathcal O(U)$.

\begin{proof}
Fix the affine open $U$ and put $Z=A_U$.
It is quasi-compact because $A\to S$ is quasi-compact.
It is quasi-separated because the restricted morphism is
quasi-separated and the affine base $U$ is separated.
Choose for every $z\in Z$ an original chart $c(z)$ of $L$
containing its image, and use (PBC.34).
These $G_i$ are units everywhere by their definition.
On $U_i\cap U_j$ the cocycle identity gives
\[
 G_i/G_j=j^*g_{ij}.
                                                        \tag{PBC.35}
\]
Thus the hypotheses (PBC.2) hold.

Call $a\in\mathcal O(U)$ good if $D_U(a)$ lies in a
neighborhood on which the assumed relative ampleness holds.
The principal-open basis on the affine scheme $U$ shows
that the $D_U(a)$ for good $a$ cover $U$:
at each point take its given neighborhood and then a
principal neighborhood contained in it.
Consequently the good elements generate the unit ideal.
Indeed, a proper ideal containing all of them would lie
in a prime ideal, whose point would be outside that cover.
Choose a finite expression for $1$ in this ideal.
It yields finitely many good $a_\ell$ which still generate $1$.
No finite-generation theorem for arbitrary ideals is being used.

For $\phi:\mathcal O(U)\to\operatorname{Fn}(Z)$ use the
actual pullback of base functions.  PBC17 gives the pack
required over each unit locus of $\phi(a_\ell)$.
Apply PBC12 to obtain one $N\ge1$ and global functions
$\sigma_0,\ldots,\sigma_m$ of degree $N$.

Now set $P=\operatorname{powOver}(A,L,U,N)$.
For $i$ let $\kappa(i)$ be the constant tuple of $N$
copies of $i$.  The original $L$-cocycle gives $P.\operatorname{IsLB}$;
each of its transition units is the product of the corresponding
pulled transition units.  On the constant tuples its transition
is $G_i^N$.  Hence PBC15 gives an actual section $s_t$ of
$P$ with chosen-chart value exactly $\sigma_t$.
PBC14 identifies its original nonzero locus with the unit
locus of $\sigma_t$.  PBC12 therefore gives the covering
and the affineness of every such open.

Fix $i$.  PBC1 gives a regular function
$r_{ij}$ on $D(s_i)$ whose germs are $\sigma_j/\sigma_i$.
By PBC14, on every original $P$-chart it satisfies
$s_j=r_{ij}s_i$.  The generation assertion (PBC.26),
combined with PBC5 for the actual coefficient map
$\mathcal O(U)\to\mathcal O(D(s_i))$, gives the surjection
\[
 \mathcal O(U)[X_0,\ldots,X_m]\longrightarrow
       \mathcal O(D(s_i)),\qquad X_j\longmapsto r_{ij}.
                                                        \tag{PBC.36}
\]
These are all fields in the original definition of
$\operatorname{RelAmple}$.  Since $U$ was arbitrary, the
assertion is the full \texttt{rel\allowbreak Ample\_\allowbreak of\_\allowbreak locally}.
\end{proof}
% END PBC18

\subsection{Affine base change with the prescribed bundle and ratios}

% BEGIN PBC19
\paragraph{PBC19: generated bundle comparisons, powers, and actual pulled sections.}
\label{claim:PBC19}
Let $F:X'\to X$, let $L,L'$ be actual line bundles, and let
$w$ be a generated comparison $L'\simeq F^*L$ in the source
sense.  It gives a pulled section of $L'$ for every actual
section of $L$.  This section satisfies the comparison at
every point of every test scheme, its nonzero locus is the
inverse image of the original one, and ratios of two such
sections are the pulled ratios.
The same comparisons exist for every finite
$\operatorname{LinComb}$ with arbitrary integer exponents,
and for the literal nested pullback and composed pullback.

\begin{proof}
On $F^{-1}(L.U_j)\cap L'.U_k$, write $w_{jk}$ for the
specified unit in the division bundle.  Its compatibility is
\[
 w_{jk}=F^*g_{jj'}(g'_{kk'})^{-1}w_{j'k'}
                                                        \tag{PBC.37}
\]
on each common overlap.  For a section $s$ of $L$, define
locally on this open
$s'_k=w_{jk}^{-1}F^*s_j$.
When $j$ changes and $k$ stays fixed, (PBC.37) and
$s_j=g_{jj'}s_{j'}$ show that these functions agree.
They glue to an actual function on $L'.U_k$.
When $k$ changes, the same identity gives
$s'_k=g'_{kk'}s'_{k'}$.
Thus they form an actual section of $L'$.
Every equality is an equality of regular functions, so after
pullback along any $r:Z\to X'$ it remains the exact germ
identity
\[
 \operatorname{germ}(F\circ r)^*s_j
    =\operatorname{germ}(r^*w_{jk})\,
                      \operatorname{germ}(r^*s'_k).
                                                        \tag{PBC.38}
\]
This is the source's all-test \texttt{sec\_\allowbreak pull\_\allowbreak down}
conclusion.  A unit multiplier preserves the property of a
germ being a unit, and a local stalk homomorphism reflects
units.  The asserted nonzero-locus equality follows.
For two sections the same unit multiplies both coordinates,
so it cancels from their ratio.

For maps $\alpha_\ell:Y\to X'$ and integers $\epsilon_\ell$,
on a common refinement multiply the pulled units as
\[
 \prod_\ell \alpha_\ell^*(w_{j_\ell k_\ell})^{\epsilon_\ell}.
                                                        \tag{PBC.39}
\]
Negative exponents use inverses of units.
Multiplying (PBC.37) to these powers gives exactly the ratio
of the two $\operatorname{LinComb}$ transition functions.
Hence (PBC.39) is a generated comparison on the original
refined charts.  It carries a generating tuple to the
product of that tuple with (PBC.39), proving the
generator-transport statement as well.
For completeness, fix a chart of the target
$\operatorname{LinComb}$ and intersect it with all charts
of the source.  They cover that target chart, since each
line bundle covers and the product index is finite.
On these intersections the prescribed product units agree,
by (PBC.37) and the old generator's compatibility.
Glue the units and their inverses to obtain a unit on the
whole target chart.  These units satisfy its transition
identity by the same equation.  After every test map $r$,
their germs equal the product (PBC.39) times the germs of
the old generator.  This is also the full all-test formula
used by the source's \texttt{exists\_\allowbreak up}.

For nested versus composed pullbacks, restriction and pullback
of the same original coordinate function agree by the
composition rule for the structure sheaf.  On matching
refined charts use the unit $1$ and on nonmatching charts
use the corresponding original transition.  The cocycle
identity proves compatibility and gives the comparison
between the literal chart records.  In particular the
identity comparison on a bundle $M$ is represented, on
the pair chart $(i,j)$ of its division bundle, by $g^M_{ij}$:
the identity
$g^M_{ij}=(g^M_{ii'}(g^M_{jj'})^{-1})g^M_{i'j'}$
is exactly its required compatibility.
Pulling any generating tuple along a chart refinement is
proved by this same equality and needs no new existence theorem.
\end{proof}
% END PBC19

% BEGIN PBC20
\paragraph{PBC20: relative ampleness when the base map lands in one affine open.}
\label{claim:PBC20}
Suppose the square
\[
 \begin{array}{ccc}
 A'&\xrightarrow{\phi}&A\\
 p'\downarrow&&\downarrow p\\
 S'&\xrightarrow{f}&S
 \end{array}
                                                        \tag{PBC.40}
\]
is Cartesian.  Let $L,L'$ be actual line bundles, with a
specified generated comparison $L'\simeq\phi^*L$.
Assume $\operatorname{RelAmple}(A,L)$, and let $U\subseteq S$
be an affine open containing $f(S')$.
Then $\operatorname{RelAmple}(A',L')$.
Here $S'$ and $f$ are arbitrary; $f$ need not be flat.
This is the full original \texttt{rel\allowbreak Ample\_\allowbreak pull}.

\begin{proof}
Fix any affine open $U'\subseteq S'$.  Restrict (PBC.40)
to obtain the actual Cartesian square
\[
 \begin{array}{ccc}
 A'_{U'}&\xrightarrow{\psi}&A_U\\
 \downarrow&&\downarrow\\
 U'&\xrightarrow{f|_{U'}}&U .
 \end{array}
                                                        \tag{PBC.41}
\]
The specified inclusion $f(U')\subseteq U$ defines this
restricted map.  The universal property of (PBC.40) and
the two open restrictions proves (PBC.41) is Cartesian.

Apply relative ampleness on $U$.  It provides $N\ge1$
and actual sections $s_0,\ldots,s_m$ of
$\operatorname{powOver}(A,L,U,N)$ with affine nonzero
loci $D_i$, ratio functions $r_{ij}$, and surjections
\[
 R[X_0,\ldots,X_m]\twoheadrightarrow Q_i,
 \quad R=\mathcal O(U),\quad Q_i=\mathcal O(D_i).
                                                        \tag{PBC.42}
\]
PBC19 compares their power bundle with the literal
$\operatorname{powOver}(A',L',U',N)$.
Pull each section through this comparison, obtaining $s'_i$.
Its nonzero locus is
$D'_i=\psi^{-1}D_i$, so the new loci cover.

Since $U,U'$ are affine, $U'\to U$ is affine, and its
base change $\psi$ is affine.  Hence each $D'_i$ is
affine.  More explicitly, put $R'=\mathcal O(U')$.
The actual affine square on $D_i$ identifies
\[
 \mathcal O(D'_i)\simeq R'\otimes_R Q_i.
                                                        \tag{PBC.43}
\]
This is the defining affine fiber-product construction:
maps out of the tensor product are pairs of ring maps
agreeing on $R$, which is precisely the square's universal
property.  Thus (PBC.43) is valid without any flatness
assumption.

Set $r'_{ij}=\psi^*r_{ij}$.  PBC19 gives on every original
power-bundle chart the equality $s'_j=r'_{ij}s'_i$.
Every pure tensor in (PBC.43) has its $Q_i$ factor represented
by a polynomial from (PBC.42).  It is therefore a polynomial
over $R'$ in the images of the $r_{ij}$.  Sums of pure
tensors give all of the tensor product.  Hence
\[
 R'[X_0,\ldots,X_m]\twoheadrightarrow\mathcal O(D'_i),
       \qquad X_j\longmapsto r'_{ij}.
                                                        \tag{PBC.44}
\]
The coefficient map and the ratio maps are the actual maps
of (PBC.41).  We have produced every field of relative
ampleness on $U'$, with the same $N$ and the pulled tuple.
As $U'$ was arbitrary, the result follows.
\end{proof}
% END PBC20

% BEGIN PBC21
\paragraph{PBC21: the literal restricted pullback bundle is locally relatively ample.}
\label{claim:PBC21}
In (PBC.40), suppose $L$ is an actual relatively ample
line bundle on $A$.  For every affine open $V\subseteq S$,
the bundle
\[
 (\phi^*L)|_{A'_{f^{-1}V}}
                                                        \tag{PBC.45}
\]
with the literal nested pullback charts is relatively ample
over $f^{-1}V$.

\begin{proof}
Let $i:A'_{f^{-1}V}\to A'$ and
$b:f^{-1}V\to S'$ be the actual open immersions.
Pasting the restriction square with (PBC.40) gives the
Cartesian square with upper map $\phi\circ i$ and lower
map $f\circ b$.  Its base image is contained in $V$.
PBC19 supplies a generated comparison
$i^*(\phi^*L)\simeq(\phi\circ i)^*L$ between the literal
two chart records.  Apply PBC20 with this comparison and
the affine open $V$.
The result is exactly (PBC.45); the possibly nonaffine
open $f^{-1}V$ has not been replaced by a smaller base.
\end{proof}
% END PBC21

\subsection{The actual group pullback and all the polarization fields}

For the remaining assertions, $A/S$ and $A'/S'$ are group
objects, and a \emph{group pullback} means (PBC.40) together
with the original all-test multiplication assertion.
Explicitly, for every $T\to S'$ and maps $x,y:T\to A'$,
composition with $\phi$ identifies their product with the
product of their two images in $A(T\to S'\to S)$.
The notation $\operatorname{toBase}$ below means exactly
this composition, with the induced structure map on $T$.

% BEGIN PBC22
\paragraph{PBC22: the pullback group, relative dimension, and the exact test maps.}
\label{claim:PBC22}
For arbitrary $f:S'\to S$, the specified scheme
$A'=A\times_S S'$ has its canonical pulled group structure
and makes (PBC.40) a group pullback.  If $A/S$ is an
abelian scheme, so is $A'/S'$.  If its original smooth
relative dimension is $g$, the pullback has the same
smooth relative dimension $g$.
For any group pullback, not only this chosen model,
$\operatorname{toBase}$ is injective on every test scheme,
preserves identities and inverses, and
\[
 \phi\circ e_{A'}=e_A\circ f,\qquad
 \phi\circ[-1]_{A'}=[-1]_A\circ\phi .
                                                        \tag{PBC.46}
\]

\begin{proof}
For a test $T\to S'$, the universal property identifies
$A'(T)$ with the maps $a:T\to A$ whose composition with
$p$ is the prescribed map $T\to S'\to S$.
This identification is natural in $T$.  Define multiplication,
identity and inverse by those of $A(T)$; their structure
maps to $S'$ remain the prescribed one.
Equivalently, lift the original multiplication and inverse
through the Cartesian square.  Uniqueness of these lifts
proves the group axioms and the original all-test product
identity.  The given categorical pullback group has exactly
these lifts, so the construction is the original
$\operatorname{Over.pullback}$ group.

For any group pullback, two maps $T\to A'$ with the same
composite to $A$ also have the same composite to $S'$
because they are over $S'$.  Cartesian uniqueness gives
their equality.  Thus $\operatorname{toBase}$ is injective.
It is multiplicative by the supplied all-test field.
A multiplicative map of groups preserves the identity
and inverse: cancel its image of $1$ in
$h(1)=h(1)^2$, then apply $h(xx^{-1})=1$.
Apply this to the identity test and to the universal point
of $A'$ to obtain (PBC.46).

Commutativity of the original group law is preserved by
the same lifts.  Smoothness and its relative dimension
are preserved by base change: on a standard smooth
presentation, base change replaces the coefficients by
their images, leaves the number of variables and relations
unchanged, and preserves the invertible Jacobian minor.
Such presentations cover the pulled scheme.
Properness is also preserved: finite type and separatedness
are stable under base change, and universal closedness
holds because every further base change is already a base
change of the original proper morphism.
Finally, a geometric fiber of $A'\to S'$ over an
algebraically closed field is the same fiber of $A\to S$
after composition of the field point with $f$.
It is connected by the original geometric-connectedness
field.  These are precisely the four fields
commutative, smooth, proper, and geometrically connected
in the original $\operatorname{IsAbelianScheme}$ record.
\end{proof}
% END PBC22

% BEGIN PBC23
\paragraph{PBC23: pulling the prescribed rigidification and local symmetry.}
\label{claim:PBC23}
Let (PBC.40) be any group pullback.
If $L$ is rigidified along $e_A$, then the literal $\phi^*L$
is rigidified along $e_{A'}$.
If $L$ satisfies the original locally-on-$S$ symmetry
condition, then $\phi^*L$ satisfies the original
locally-on-$S'$ symmetry condition.
The actual generating tuples are pulled from the given ones.

\begin{proof}
The given rigidification is a compatible unit tuple
trivializing $e_A^*L$.  Pull it by $f$ and use the first
identity in (PBC.46) to compare
$f^*e_A^*L$ with $e_{A'}^*\phi^*L$.
On each common chart the tuple is the pullback of the old
unit.  Its compatibility follows by pulling the old
compatibility equation.  The chart comparison of PBC19
puts it in the literal target chart record.

For symmetry, fix $s'\in S'$.  At $f(s')$ the original
condition supplies an open $U\subseteq S$ and a
generating tuple of
$[-1]_A^*L\otimes L^{-1}$ on $A_U$.
Use the open $f^{-1}U$, which contains $s'$.
The Cartesian square restricts to a map
$A'_{f^{-1}U}\to A_U$.
The second identity in (PBC.46) says this map commutes
with the inversion maps.  Pull the given tuple and
apply the finite two-factor comparison (PBC.39) with
exponents $(1,-1)$.  It is a generator of precisely the
source's $\operatorname{LinComb}$ on this restricted
space.  This is the required local symmetry at $s'$.
No global symmetry generator is inferred from the local ones.
\end{proof}
% END PBC23

% BEGIN PBC24
\paragraph{PBC24: the full affine-test translation-kernel condition survives base change.}
\label{claim:PBC24}
Let (PBC.40) be a group pullback, let $L,L'$ be actual
line bundles, and fix a generated comparison
$L'\simeq\phi^*L$.
If $\operatorname{KTriv}(A,L)$ holds, then
$\operatorname{KTriv}(A',L')$ holds, for every commutative
test ring in its original definition.

\begin{proof}
Take any commutative ring $R$, put $T=\operatorname{Spec}R$,
and fix $t:T\to S'$ and $x:T\to A'$ over $S'$.
Write $x_0=\phi\circ x$, regarded over $f\circ t$.
The two spaces on which the translation differences are defined are
\[
 Y'=T\times_{S'}A',\qquad Y=T\times_S A.
                                                        \tag{PBC.47}
\]
There is a canonical map
$\alpha:Y'\to Y$ with first component $\operatorname{pr}_T$
and second component $\phi\circ\operatorname{pr}_{A'}$.
It is an isomorphism.  To construct its inverse, on $Y$ lift
the two compatible maps
$\operatorname{pr}_A:Y\to A$ and
$t\circ\operatorname{pr}_T:Y\to S'$ through (PBC.40), and
then pair the resulting map to $A'$ with $\operatorname{pr}_T$.
Both composites are identities because their two projections
are identities.  This proves the assertion on the actual
fiber-product models.

Let $\tau_x:Y'\to A'$ be addition of $x$ to the second
component and let $\tau_{x_0}:Y\to A$ be the analogous map.
The all-test multiplication field, applied to the universal
points on $Y'$, gives
\[
 \phi\circ\tau_x=\tau_{x_0}\circ\alpha,\qquad
 \phi\circ\operatorname{pr}_{A'}
                         =\operatorname{pr}_A\circ\alpha .
                                                        \tag{PBC.48}
\]
By PBC19 the specified comparison for $L',L$ therefore
gives a generated comparison
\[
 \tau_x^*L'\otimes\operatorname{pr}_{A'}^*(L')^{-1}
 \simeq
 \alpha^*\bigl(
    \tau_{x_0}^*L\otimes\operatorname{pr}_A^*L^{-1}\bigr).
                                                        \tag{PBC.49}
\]
Its units on a common chart are the units of $w$ pulled
along $\tau_x$, divided by those pulled along
$\operatorname{pr}_{A'}$.  Thus (PBC.49) is a comparison
of the literal $\operatorname{translDiff}$ bundles.

Now suppose the antecedent in $\operatorname{KTriv}(A',L')$
holds.  For every $p\in T$ it gives an open $V\ni p$ and
an actual generating tuple of the left side of (PBC.49)
over $\operatorname{pr}_T^{-1}V$.
The isomorphism $\alpha$ commutes with $\operatorname{pr}_T$,
so it restricts to an isomorphism over this same $V$.
Transport the tuple by (PBC.49) and by this restricted
isomorphism.  It gives an actual generator of
$\operatorname{translDiff}(A,L,x_0)$ over
$\operatorname{pr}_T^{-1}V$ in $Y$.
Consequently $x_0$ satisfies the antecedent of the
original $\operatorname{KTriv}(A,L)$ for the same ring
$R$ and the composite map $f\circ t$.
That condition forces $x_0=e_A\circ f\circ t$.

PBC22 identifies this latter point as the image of the
identity point of $A'(T)$, and injectivity there gives
$x=e_{A'}\circ t$.  This is the exact original conclusion.
The argument used an isomorphism in (PBC.47), not any
unproved descent of sections along a nonflat map.
\end{proof}
% END PBC24

% BEGIN PBC25
\paragraph{PBC25: full polarized base change with all original fields.}
\label{claim:PBC25}
For every pair of schemes $S,S'$, every $f:S'\to S$, and
every original polarized record $P=\operatorname{PolAV}(S)$,
there are a polarized record $P'$ over $S'$ and a Cartesian
map $\phi:P'.A\to P.A$ over $f$ such that
\[
 \begin{gathered}
 P'.g=P.g,\\
 \operatorname{toBase}(x y)
       =\operatorname{toBase}(x)\operatorname{toBase}(y)
       \quad\text{for every }T/S'\text{ and every }x,y\in P'.A(T),\\
 \text{there is an actual generated }
       \operatorname{LB.Iso}(P.L.\operatorname{pull}(\phi),P'.L).
 \end{gathered}                                           \tag{PBC.50}
\]
In fact $P'.A$ is the specified categorical pullback
$P.A\times_S S'$, $P'.L$ is literally $P.L.\operatorname{pull}(\phi)$,
and $P'.g$ is defined to be $P.g$.
This proves the full original $\operatorname{PolBaseChangeSt}$.

\begin{proof}
Use the actual pullback from PBC22 with its pulled group
structure.  It is abelian and has the same smooth relative
dimension $g=P.g$.  Set $L'=\phi^*P.L$ with the original
pullback charts.  Pulling the transition units preserves
the cocycle and the covering, so $L'$ is an actual line bundle.
PBC23 supplies its rigidification and local symmetry from
$P.\operatorname{rigid}$ and $P.\operatorname{sym}$.
PBC24 supplies its $\operatorname{KTriv}$ from
$P.\operatorname{ktriv}$, using the identity comparison
on $L'$ constructed in PBC19.

It remains to check the relative-ample field on this same
bundle.  The pullback is proper by PBC22, hence its
structural morphism is quasi-compact and quasi-separated.
For every $s'\in S'$, choose an affine open $V\subseteq S$
containing $f(s')$.  PBC21 gives relative ampleness of the
literal restriction of $L'$ over $f^{-1}V$.
PBC18 then proves relative ampleness on all of $S'$.

These constructions supply every field of the original
$\operatorname{PolAV}$ record: group, abelian property,
$g$, smooth relative dimension, actual bundle, line-bundle
compatibility, rigidification, symmetry, relative ampleness,
and $\operatorname{KTriv}$.  Define $P'$ using them.
The Cartesian square is the original pullback square.
The dimension equality in (PBC.50) is equality by the
definition of $P'.g$.  The universal multiplication assertion
is exactly PBC22, for every test $T$, without an affine or
field restriction on that test.

Finally, since $P'.L=P.L.\operatorname{pull}(\phi)$,
use the actual identity division-bundle generator from
PBC19.  Its coordinate on the pair chart $(i,j)$ is the
pulled bundle's transition $g^{L'}_{ij}$; it satisfies
the generator compatibility by the cocycle identity.
Thus the last output is the source's generated
$\operatorname{LB.Iso}$, not only an isomorphism class
of line bundles.  This completes (PBC.50) for arbitrary
$S,S'$ and $f$.
\end{proof}
% END PBC25

% BEGIN PBC26
\paragraph{PBC26: the complete all-affine-test very-ample producer.}
\label{claim:PBC26}
For every original $P/S$, every commutative ring $R$, every
$t:\operatorname{Spec}R\to S$, and every specified Cartesian
model
\[
 \begin{array}{ccc}
 X&\xrightarrow{\pi}&P.A\\
 f_X\downarrow&&\downarrow\\
 \operatorname{Spec}R&\xrightarrow{t}&S,
 \end{array}
\]
there are $N\ge1$, $m\ge0$ and actual sections
$s_0,\ldots,s_m$ of
$\operatorname{pw}(P.L.\operatorname{pull}(\pi),N)$
such that this power is an actual line bundle and
$\operatorname{VeryAmpleBy}(f_X,\operatorname{pw}(\pi^*P.L,N),s)$.
This is the original \texttt{pol\allowbreak AV\_\allowbreak very\allowbreak Ample\allowbreak By} with its
$\operatorname{PolBaseChangeSt}$ premise now supplied.

\begin{proof}
PBC25 provides the polarized pullback over
$\operatorname{Spec}R$.  The prescribed $X$ is uniquely
isomorphic to its underlying pullback, commuting with both
projections.  Apply its relative ampleness to the whole affine
base.  It gives one positive power, the finite tuple of
sections, affine nonzero loci, actual ratio identities and
the surjective polynomial maps over
$\mathcal O(\operatorname{Spec}R)$.

The already proved GSL15 transports precisely these data
through the unique Cartesian isomorphism and the generated
bundle comparison to $X$.  Its argument uses GSL14 for
the actual sections, PBC19's unit cancellation for the
ratio functions, and the canonical coefficient-ring
isomorphism $R\simeq\mathcal O(\operatorname{Spec}R)$ for
the polynomial maps.  It also changes the repeated-factor
power to the literal $\operatorname{pw}$ chart record.
The resulting data are exactly the asserted
$\operatorname{VeryAmpleBy}$, including the line-bundle
compatibility of that power.

The quantifiers allow $N$ and the tuple to depend on this
specified affine test and Cartesian model.  They do not
assert one uniform ample power over all tests.
The ring $R$ is arbitrary, with no reducedness, locality,
Noetherianity or flatness condition.
\end{proof}
% END PBC26

% BEGIN PBC27
\paragraph{PBC27: actual section finiteness and the return to the original power-section module.}
\label{claim:PBC27}
Let $K$ be any field, $f:X\to\operatorname{Spec}K$, $M$
an original line-bundle chart datum, and
$Q=\operatorname{Perf}(f,M)$.
If the degree-zero matrix size of $Q$ is $a$, every
$a+1$ actual sections $s_0,\ldots,s_a$ satisfy a nonzero
$K$-linear relation in their coordinates on every original
$M$-chart.
Consequently, for every algebraically closed $K$, every
original $P/\operatorname{Spec}K$ and every actual line
bundle $M$ on $P.A$, there is one $n$ such that every
$n$-tuple of actual $M$-sections has a nonzero relation
over $\mathcal O(\operatorname{Spec}K)$, using the original
structural maps on every chart.  Moreover
\[
 \operatorname{SecM}(P.A,P.L,j,\operatorname{baseOver}(\operatorname{Spec}K))
 \quad\hbox{is finite over }\mathcal O(\operatorname{Spec}K)
       \quad\hbox{for every }j\in\mathbb Z.
                                                        \tag{PBC.51}
\]
More generally, over any field $K$, conclusion (PBC.51)
follows whenever the displayed finite-relation property is
supplied for every actual line bundle on $P.A$; algebraic
closedness is not a premise of this last implication.

\begin{proof}
Use the identity base change in $Q.\operatorname{sec}$.
For each $s_k$, first express it on the literal identity-pullback
charts of $M$.  Its preimage under the provided section
bijection is a vector $v_k\in\ker D_0\subseteq K^a$.
The actual coordinate identity in $\operatorname{Perf}$ says
that on its fixed cover $W_\alpha\subseteq M.U_{c(\alpha)}$,
\[
 s_{k,c(\alpha)}|_{W_\alpha}
    =\sum_{\ell=1}^{a} f^*(v_{k\ell})u_{\ell\alpha}.
                                                        \tag{PBC.52}
\]
This identity, rather than any unstated linearity of the
chosen bijection, is all that is needed.
The $a+1$ vectors in $K^a$ have a relation
$\sum_k d_kv_k=0$ with $(d_k)\ne0$.
For $a=0$, the sole vector is zero and $d_0=1$ works.
Multiply (PBC.52) by $f^*d_k$, sum in $k$, and interchange
the finite sums.  Every coefficient of every
$u_{\ell\alpha}$ is zero.  Thus
$\sum_k f^*(d_k)s_{k,c(\alpha)}=0$ on $W_\alpha$.

For any original chart $M.U_i$, restrict to
$M.U_i\cap W_\alpha$ and use the transition unit
$g_{i,c(\alpha)}$ to transfer this same relation to the
$i$th coordinates.  The opens $W_\alpha$ cover $X$.
Equality of sections on a cover proves
$\sum_k f^*(d_k)s_{k,i}=0$ on all of $M.U_i$.
This is the full original $\operatorname{Perf.exists\_\allowbreak relation}$.

For the asserted polarized variety, use PBC26 on the
identity Cartesian model of $P.A$.  It supplies the
original $\operatorname{IsProjBy}$ presentation.
The structure map is smooth by $P.\operatorname{isAb}$.
ASL21 therefore provides an actual
$Q=\operatorname{Perf}(P.A\to\operatorname{Spec}K,M)$.
Choose $n=Q.a+1$ before the tuple of sections.
Apply the first part and send its coefficients through the
canonical ring isomorphism
$K\simeq\mathcal O(\operatorname{Spec}K)$.
It is injective, so the coefficient tuple remains nonzero.
Its composite with the structural map to each chart is
exactly the scalar map in (PBC.52).
This proves the complete original
\texttt{sections\_\allowbreak finite\_\allowbreak of\_\allowbreak perf} conclusion.

To obtain (PBC.51), put $T=\operatorname{baseOver}(\operatorname{Spec}K)$.
The second projection
$\beta:T\times_{\operatorname{Spec}K}P.A\to P.A$
is an isomorphism with inverse the pair of the structural
map and the identity of $P.A$; the two projection identities
prove this directly.
On $P.A$ take the actual integer-power bundle
$M_j=\operatorname{LinComb}(P.L,\operatorname{id},j)$.
Apply the just proved finite relation statement to $M_j$.
A tuple in the source's $\operatorname{SecM}$ is a tuple
of actual sections on the identity fiber product.
Pull it back by $\beta^{-1}$, identifying its finite-singleton
power charts with those of $M_j$ by PBC19.
The resulting nonzero relation pulls back to the same
relation on the fiber product, because the scalar maps
commute with both projections.  Thus every tuple of the
chosen length in $\operatorname{SecM}$ is linearly dependent.

Finally, over a field, if every $n$-tuple is dependent then
the whole vector space is finitely generated: starting with
the empty independent family, repeatedly adjoin a vector
outside its span; this either terminates with a spanning
family or would produce an independent $n$-tuple.
Here the coefficient ring is a field by
$K\simeq\mathcal O(\operatorname{Spec}K)$.
This proves finiteness with the original coefficient action.
The identity-product argument works over any field whenever
the displayed original finite-relation hypothesis is supplied,
which is also the precise scope of the source's
\texttt{finite\_\allowbreak sec\allowbreak M}.
\end{proof}
% END PBC27

\subsection{Return to the section-law construction and exact remaining scope}

PBC25 supplies the original input denoted $\mathcal P$ in GSL15--GSL17
and \texttt{GSL\_\allowbreak POL\_\allowbreak BASE\_\allowbreak CHANGE} in the existing contracts.
The construction retains a fixed $P$ and pulls its actual group,
line bundle and specified generators.  It therefore can be inserted
in the previous conditional arguments without changing their sections
or their coefficient maps.
PBC26 supplies their all-test projective presentation.
ASL21's original smooth-projective perfect-complex theorem can now
be applied with that presentation on the specified affine pullback.
The completed ASL and SHV proofs are not repeated.

The two Serre conclusions of SHV17 and SHV21 remain available at their
original Noetherian-base scopes.  They do not impose Noetherianity on
PBC18 or PBC25, which were proved independently over arbitrary bases.
Conversely, supplying polarized base change does not prove the
remaining geometric inputs to GSL17: the multiplication-pullback
generator, the finite/flat/surjective and away-characteristic etale
multiplication geometry, and invariant-function descent on every
preimage of an open still have to be supplied.

PBC27 supplies actual section finiteness, including every integer
power in the original $\operatorname{SecM}$ model.
The geometric rank-one theorem, torsion-indexed section decomposition,
two-sided section growth, and the square/translation generators are also
separate original producers.  Until their exact remaining hypotheses
are discharged, full global
$\operatorname{SecBaseChangeSt}$,
$\operatorname{SecBaseChangeAnySt}$ and
$\operatorname{SecTwistSt}$ are not asserted by this chapter.
No change is made to the retained moduli or finite-level dictionary
and class-sum obligations, and no whole selected obligation is closed.

\paragraph{Source map.}
The following references identify the mathematical producer bodies;
the accompanying provenance binds their exact bytes and line ranges.
PBC1--PBC9 reconstruct
\texttt{Exists\allowbreak One\allowbreak Add\allowbreak Kill.\allowbreak lean}: the regular-function calculus,
\texttt{kill}, \texttt{frac}, \texttt{local\_\allowbreak global},
\texttt{exists\_\allowbreak one\_\allowbreak add\_\allowbreak kill}, \texttt{Sc.\allowbreak extend}, and
\texttt{extra\_\allowbreak section}.
PBC10--PBC18 reconstruct
\texttt{Rel\allowbreak Ample\allowbreak Of\allowbreak Locally\allowbreak On\allowbreak Base.\allowbreak lean}: the full
\texttt{glue\_\allowbreak pack}, open-immersion transport,
\texttt{gen\_\allowbreak iff}, the actual-section dictionary,
\texttt{exists\_\allowbreak sv}, \texttt{cst\_\allowbreak mem},
\texttt{pack\_\allowbreak of\_\allowbreak rel\allowbreak Ample}, and
\texttt{rel\allowbreak Ample\_\allowbreak of\_\allowbreak locally}.

PBC19--PBC24 reconstruct the relevant complete chart and pullback
arguments in \texttt{Ample\allowbreak Pt.\allowbreak lean}, \texttt{Sab\allowbreak Pull.\allowbreak lean},
\texttt{Theta\allowbreak Core.\allowbreak lean}, \texttt{Is\allowbreak Mon\allowbreak Hom\allowbreak Of\allowbreak Mul.\allowbreak lean},
\texttt{Const\allowbreak Pullback.\allowbreak lean}, and
\texttt{Pol\allowbreak Base\allowbreak Change\allowbreak Of\allowbreak Rel\allowbreak Ample\allowbreak Local.\allowbreak lean}.
The last of these supplies the constructor returned by PBC25.
The literal polarization and all-test output contracts are in
\texttt{Nonzero\allowbreak Locus.\allowbreak lean} and
\texttt{Global\allowbreak Sections\allowbreak Rank\allowbreak Alg\allowbreak Closed.\allowbreak lean}.
PBC26 returns to
\texttt{Very\allowbreak Ample\allowbreak By\allowbreak Of\allowbreak Rel\allowbreak Ample.\allowbreak lean}, using the already reviewed
GSL14--GSL15 proof blocks at their exact scopes.
PBC27 proves the original producers in
\texttt{Abelian\allowbreak Variety\allowbreak Line\allowbreak Bundle\allowbreak Sections\allowbreak Finite.\allowbreak lean} and
the \texttt{finite\_\allowbreak sec\allowbreak M} argument of
\texttt{Finite\allowbreak Of\allowbreak Not\allowbreak Linear\allowbreak Independent.\allowbreak lean}, using ASL21.


\section{The actual local deformation presentation and its relation count}
\label{sec:arithmetic-presentations}

\paragraph{The local ring that will be presented.}
Fix the original complete coefficient DVR \(O\), its finite residue
field \(k\) of characteristic \(\ell\), a frame
\(F.\mathrm{Ok}(C_0)\), and one supplied
\(S:\mathrm{Space}(O,F,n)\), where \(n=2\) or \(3\).
Write \(\mathcal A=S.A\), \(\mathcal B=S.B\), and
\(u=S.u\).  For a supplied original local object
\(L:S.\mathrm{Loc}(x)\), put
\[
 O'=L.O',\qquad R=L.R,\qquad
 \iota=L.\iota:\mathcal A\longrightarrow R,
 \qquad \rho=\iota_*u.                                      \tag{APD.1}
\]
All constructions below use these rings and this map.  In particular,
the final assertion is the literal presentation
\[
 R\simeq O'[[X_1,\ldots,X_N]]/(f_1,\ldots,f_{N'}),
 \qquad N=N'+nF.d+1.                                      \tag{APD.2}
\]
It holds for every supplied \(L\); it is not obtained by choosing a
more favorable point or a further coefficient extension.

The argument has three parts.  First we prove the commutative algebra
that converts an obstruction space into a bound on relations.  Next
we show that the actual local normalized representation represents its
deformation problem, including points with triangular, rather than
diagonal, reduction.  Finally we construct its continuous obstruction
cocycles, normalize their corrections, and compute the difference
between tangent and obstruction dimensions.  The frame's original
Euler field is used exactly as a field of \(F.\mathrm{Ok}(C_0)\).
Its arithmetic construction was supplied by LFC20; the original
\(P\) hypothesis at the arithmetic frame entry is not changed here.

\paragraph{Conventions and available inputs.}
For any local ring \(\mathcal O\), an
\(\mathrm{IsCO}(\mathcal O,C)\)-algebra is a complete Noetherian
local \(\mathcal O\)-algebra whose specified coefficient map induces
an isomorphism of residue fields.  Its residue map to
\(\kappa=\mathcal O/\mathfrak m_{\mathcal O}\) is denoted
\(\operatorname{res}_C\).  Every homomorphism between these coefficient
algebras is local, preserves this residue map, and is adically
continuous: the unit argument in ALC1 proves the first two facts,
and preservation of powers proves the third.  The general statements
below explicitly say when \(\mathcal O\) must itself be Noetherian
or complete; the arithmetic DVR assumptions are not imposed on the
broader statements.

We use the complete Noetherian algebra and Artin--Rees proofs in
ALC5--ALC7, the actual local chart and representation in ALC11 and
ALC17, the original weighted entry algebra in RTC7, its fixed-weight
continuity in RTC9, and the ordered flag and characteristic-polynomial
arguments in ACP3--ACP4.  We also use the elementary Nakayama lemma,
the Noetherian formal-power-series theorem, and compactness of an
inverse limit of finite discrete sets.  The constructions and all
deformation-specific applications of these facts are proved here.

\subsection{Minimal generators and obstruction spaces}

% BEGIN APD1
\paragraph{APD1. Tangent vectors, relative cotangents, and dual numbers.}
\label{claim:APD1}
Let \(\mathcal O\) be any commutative local ring and let
\(R\) be an \(\mathrm{IsCO}(\mathcal O,R)\)-algebra.  Put
\[
 T=\operatorname{Der}_{\mathcal O}(R,\kappa),\qquad
 V=\mathfrak m_R/(\mathfrak m_R^2+
                         \mathfrak m_{\mathcal O}R).
                                                               \tag{APD.3}
\]
Here the target of a derivation is an \(R\)-module by
\(\operatorname{res}_R\).  Restriction gives
\(T\simeq\operatorname{Hom}_\kappa(V,\kappa)\), so \(T\)
is finite dimensional.  There is a bijection
\[
 \operatorname{Hom}_{\mathcal O\text{-alg}}
       (R,\kappa[\epsilon]/\epsilon^2)\simeq T,
 \qquad f_D(r)=\operatorname{res}_R(r)+\epsilon D(r).
                                                               \tag{APD.4}
\]
This proves the full original \texttt{pr\_\allowbreak tan}, including its
cardinality assertion when the residue field is not finite.

\begin{proof}
A derivation kills \(\mathcal O\), and the Leibniz rule makes it
zero on \(\mathfrak m_R^2\) and on
\(\mathfrak m_{\mathcal O}R\).  Its restriction therefore gives
the indicated linear functional.  Conversely let \(a\) be such
a functional.  For \(r\in R\), choose \(c\in\mathcal O\)
with \(r-c\in\mathfrak m_R\), and set
\(D_a(r)=a([r-c])\).  A different choice changes \(r-c\) by an
element of \(\mathfrak m_{\mathcal O}R\), so the definition is
independent of the choice.  Choosing lifts for two residues shows
additivity and \(\mathcal O\)-linearity.  Expanding
\((c+y)(d+z)\), with \(y,z\in\mathfrak m_R\), shows
\(D_a(rs)=\bar rD_a(s)+\bar sD_a(r)\), since \(yz\) dies in
\(V\).  These two constructions are inverse.

The ideal \(\mathfrak m_R\) is finitely generated because \(R\)
is Noetherian.  Its images generate \(V\) over the residue field;
thus \(V\), and hence its dual, are finite dimensional.  No
Noetherian hypothesis on \(\mathcal O\) was used.

The residue field and the dual-number ring are both
\(\mathrm{IsCO}(\mathcal O,-)\)-algebras.  For the latter the
maximal ideal is \((\epsilon)\), its square is zero, and every
ideal is a finite-dimensional \(\kappa\)-space.  Thus any algebra
map to it has first coordinate \(\operatorname{res}_R\).
Its second coordinate is a derivation, by multiplying dual numbers.
Conversely the Leibniz rule proves that \(f_D\) in (APD.4) is
an algebra homomorphism.  This gives the claimed bijection, not
merely equality of two finite counts.
\end{proof}
% END APD1

% BEGIN APD2
\paragraph{APD2. The minimal power-series presentation is constructed.}
\label{claim:APD2}
Suppose in addition that \(\mathcal O\) is complete and
Noetherian.  If \(t=\dim_\kappa T\), there is a surjective
continuous \(\mathcal O\)-algebra map
\[
 p:P=\mathcal O[[X_1,\ldots,X_t]]\longrightarrow R,
 \qquad I=\ker p\subset
        \mathfrak m_{\mathcal O}P+\mathfrak m_P^2.             \tag{APD.5}
\]
Equivalently, every linear coefficient of every element of \(I\)
belongs to \(\mathfrak m_{\mathcal O}\).  The number of variables
is exactly the tangent dimension.

\begin{proof}
Choose \(z_1,\ldots,z_t\in\mathfrak m_R\) lifting a basis
of \(V\) in (APD.3).  Nakayama applied to the finite module
\(\mathfrak m_R/(\mathfrak m_{\mathcal O}R+(z_1,\ldots,z_t))\)
gives
\(\mathfrak m_R=\mathfrak m_{\mathcal O}R+(z_1,\ldots,z_t)\).
Evaluation of a power series at the \(z_i\) converges in \(R\)
and defines \(p\).  To prove surjectivity, first lift the residue
of a prescribed \(r\in R\) by an element of \(\mathcal O\).
At stage \(a\), an error in \(\mathfrak m_R^a\) is a sum of
products of \(a\) of the indicated maximal-ideal generators.
Replace the remaining coefficients by lifts of their residues from
\(\mathcal O\).  This lifts the error modulo
\(\mathfrak m_R^{a+1}\) by an element of \(\mathfrak m_P^a\).
Successive corrections converge in \(P\) and their image is \(r\).

If \(f\in I\), write its constant and linear terms as
\(c+\sum_i a_iX_i\).  The constant coefficient \(c\) lies
in \(\mathfrak m_{\mathcal O}\), because (p) preserves
residues.  Reducing \(p(f)=0\) in \(V\) gives
\(\sum_i\overline{a_i}[z_i]=0\).  The chosen basis forces
\(a_i\in\mathfrak m_{\mathcal O}\) for every \(i\), proving
(APD.5).  Conversely this coefficient condition is exactly the
displayed ideal containment, by separating terms of degrees zero,
one, and at least two.

The same approximation argument gives a useful criterion used below:
if \(f:R_1\to R_2\) is a local map of complete Noetherian local
rings, induces a surjection on residue fields, and
\(\mathfrak m_{R_2}=f(\mathfrak m_{R_1})R_2\), then \(f\)
is surjective.  A correction of order \(a\) is a sum of degree
\(a\) products of images; its coefficients can be lifted modulo
the maximal ideal, and completeness sums the resulting corrections.
This criterion does not require an injection of coefficient fields.
\end{proof}
% END APD2

% BEGIN APD3
\paragraph{APD3. The space of necessary relations and its coordinates.}
\label{claim:APD3}
For (APD.5), set \(C=P/\mathfrak m_PI\) and
\(K=I/\mathfrak m_PI\).  Then \(C\to R\) is a surjective
map of \(\mathrm{IsCO}(\mathcal O,-)\)-algebras,
\(\mathfrak m_CK=0\), and \(K\) is a finite-dimensional
\(\kappa\)-space.  A basis of \(K\) lifts to a generating
family of the ideal \(I\).

More generally, let \(C\) be any such coefficient algebra and
let \(K\subset C\) be any ideal annihilated by
\(\mathfrak m_C\).  Suppose a finite family \((e_i)\)
generates \(K\), and every relation
\(\sum_i a_ie_i=0\), with \(a_i\in\mathcal O\), has all
\(a_i\in\mathfrak m_{\mathcal O}\).  Then there are
\(\mathcal O\)-linear coordinates \(c_i:K\to\kappa\) such
that every \(y\in K\) can be written
\[
 y=\sum_i a_ie_i,\qquad \overline{a_i}=c_i(y).
                                                               \tag{APD.6}
\]
This is the full original \texttt{pr\_\allowbreak dualb}.

\begin{proof}
The finite-variable power-series ring \(P\) is complete
Noetherian local.  Its ideals are closed, and its quotients are
complete, by exactness of finite-module completion in ALC7.
Both quotient rings therefore have the asserted coefficient
structures.  Multiplication by \(\mathfrak m_C\) on \(K\)
is zero by construction.  Since \(I\) is a finite \(P\)-module,
\(I/\mathfrak m_PI\) is finite dimensional.  Nakayama applied
to \(I\) shows that lifts of any basis generate \(I\).

In the more general assertion multiplication by an element of \(C\)
on \(K\) depends only on its residue, so \(K\) is a
\(\kappa\)-space.  Any ideal expression for \(y\) in terms
of the \(e_i\) can replace its coefficients in \(C\) by
residue lifts from \(\mathcal O\).  Thus the \(e_i\) span
over \(\kappa\).  A linear dependence over \(\kappa\)
can also be lifted to \(\mathcal O\); its value is zero because
\(\mathfrak m_CK=0\).  The stated relation hypothesis makes
every coefficient zero.  Hence \((e_i)\) is a basis.
Take its coordinate functionals and lift their values back to
\(\mathcal O\); this gives (APD.6) and proves linearity as well.
The empty-family case gives \(K=0\) and is included.
\end{proof}
% END APD3

% BEGIN APD4
\paragraph{APD4. A nonzero minimal relation cannot split off.}
\label{claim:APD4}
Let \(C\twoheadrightarrow R\) be the small extension of APD3,
and let \(u:K\to\kappa\) be a nonzero linear functional.
The quotient
\[
 C_u=C/\ker u\longrightarrow R                            \tag{APD.7}
\]
does not admit a \(\mathcal O\)-algebra section.

\begin{proof}
Every \(\kappa\)-subspace of \(K\) is an ideal of \(C\),
because its scalar action factors through \(C/\mathfrak m_C\).
Thus \(C_u\) is a complete Noetherian local coefficient algebra,
and its kernel over \(R\) is the nonzero one-dimensional space
\(K/\ker u\).  Write this kernel as \(J\).
Minimality (APD.5), after passing to \(C_u\), gives
\[
 J\subset\mathfrak m_{\mathcal O}C_u+\mathfrak m_{C_u}^2.
                                                               \tag{APD.8}
\]
Suppose \(s:R\to C_u\) were a section.  Every
\(a\in\mathfrak m_{C_u}\) differs from \(s(p(a))\) by an
element of \(J\).  The first summand on the right of (APD.8)
already belongs to \(s(\mathfrak m_R)C_u\).  Consequently
\[
 \mathfrak m_{C_u}=s(\mathfrak m_R)C_u+\mathfrak m_{C_u}^2.
\]
Nakayama makes the last square unnecessary.  The surjectivity
criterion proved in APD2 now makes \(s\) surjective.  Since
\(p\circ s=1\), a surjective section makes \(p\) an
isomorphism, contradicting \(J\ne0\).  This argument proves
the obstruction to splitting, rather than assuming a minimal
presentation has a universal obstruction theory.
\end{proof}
% END APD4

% BEGIN APD5
\paragraph{APD5. Full obstruction-to-presentation algebra.}
\label{claim:APD5}
Let \(\mathcal O\) be complete Noetherian local with finite
residue field \(\kappa\), let \(R\) be an
\(\mathrm{IsCO}(\mathcal O,R)\)-algebra, and let \(H\) be
a finite-dimensional \(\kappa\)-space.  Suppose the original
\(\mathrm{HasObs}(\mathcal O,R,H)\) holds, in the following
literal sense.  For every surjection \(p:C\to R\) of coefficient
algebras with \(\mathfrak m_C\ker p=0\), there is a linear map
\[
 \mathrm{ob}_p:\operatorname{Hom}_{\mathcal O}(\ker p,\kappa)
                      \longrightarrow H                 \tag{APD.9}
\]
such that, whenever \(\mathrm{ob}_p(u)=0\), every factorization
\(C\overset q\twoheadrightarrow C'\overset{p'}\longrightarrow R\)
with \(p'q=p\) and
\(u(y)=0\Rightarrow q(y)=0\) for \(y\in\ker p\)
has an \(\mathcal O\)-algebra section of \(p'\).
Then there is a surjection
\[
 \mathcal O[[X_1,\ldots,X_t]]\longrightarrow R,
 \qquad t=\dim_\kappa T,
                                                               \tag{APD.10}
\]
whose kernel is generated by a specified family of exactly
\(h=\dim_\kappa H\) elements.  This proves the full original
\texttt{pr\_\allowbreak mazur}.

\begin{proof}
Use APD2 and the small extension of APD3.  If a nonzero
\(u\in K^\vee\) lay in the kernel of (APD.9), apply its
splitting conclusion to the particular quotient \(C_u\) in
APD4.  Its hypotheses hold exactly: the quotient kills
\(\ker u\), is surjective, and factors the map to \(R\).
The resulting section contradicts APD4.  Therefore
\(K^\vee\hookrightarrow H\), and
\(r=\dim K\le h\).  Choose a basis of \(K\) and lift it
to \(f_1,\ldots,f_r\in I\).  APD3 says these generate
\(I\).  Append \(h-r\) zero elements to obtain the required
family of length \(h\).  This proves both the exact variable
count and the stated bound on generators of the kernel.

No compatibility of the maps \(\mathrm{ob}_p\) for different
small extensions was assumed or needed.  The proof in fact works
over an infinite residue field too; the finite-residue restriction
in the named original statement is certainly preserved.
\end{proof}
% END APD5

% BEGIN APD6
\paragraph{APD6. The strict first-cohomology counting lemma.}
\label{claim:APD6}
Let \(K\) be any field, \(G\) a group with any topology,
\(W,T\) finite-dimensional \(K\)-spaces, \(r\) a
representation of \(G\) on \(W\), and
\(\operatorname{loc}_v:G_v\to G\) any family of group
homomorphisms.  Let \(b_v\subset W\) and
\(\phi:T\to\operatorname{Map}(G,W)\) be linear.  A strict
one-cocycle is a pair \((z,(a_v))\) with \(z\) locally
constant, \(z(st)=z(s)+r(s)z(t)\), and
\[
 z(\operatorname{loc}_v s)-(r(\operatorname{loc}_v s)-1)a_v
                                                        \in b_v.
                                                               \tag{APD.11}
\]
The boundaries are the pairs with \(z(s)=(r(s)-1)X\) and
\(a_v-X\in b_v\).  Quotienting the strict cocycle space by
its intersection with this boundary space defines \(H^1_C\); put
\(H^0_C=W^G\cap\bigcap_v b_v\).
Suppose
\[
 \begin{split}
 &(r(\operatorname{loc}_v s)-1)Y\in b_v\ \text{for all }s
                         \quad\Longrightarrow\quad Y\in b_v,\\
 &\text{for every }(z,a)\text{ as above, there are }D,X
        \text{ with }\phi(D)=z+(r-1)X.
 \end{split}                                                  \tag{APD.12}
\]
Then \(\dim H^1_C\le\dim T\).  Moreover if
\(X_0\ne0\) and \(\phi(D_0)=(r-1)X_0\), then
\[
           \dim H^1_C+1\le\dim T+\dim H^0_C.                 \tag{APD.13}
\]
This is full \texttt{pr\_\allowbreak count}; neither profiniteness nor
finiteness of the field or of the family of local groups is imposed.

\begin{proof}
Let \(C\) be the image of the coboundary map
\(W\to\operatorname{Map}(G,W)\), \(X\mapsto(r-1)X\).
Send a strict cocycle to \([z]\) in
\(\operatorname{Map}(G,W)/C\).  If \(z=(r-1)X\), equation
(APD.11) and the first clause of (APD.12) imply
\(X-a_v\in b_v\) for every \(v\).  Thus the kernel is
exactly the strict boundaries.  The map induced on \(H^1_C\)
is injective, and the second clause of (APD.12) puts its image
inside the image of
\(f:T\to\operatorname{Map}(G,W)/C\).
This proves the first dimension inequality and also proves that
\(H^1_C\) is finite dimensional.

If \(X_0\) is invariant, the first clause of (APD.12) puts it
in every \(b_v\); it is a nonzero element of \(H^0_C\).
The first inequality therefore implies (APD.13).  If it is not
invariant, \(D_0\ne0\), since \(\phi(D_0)\ne0\), but
\(D_0\in\ker f\).  Rank--nullity now gives
\(\dim\operatorname{im}f\le\dim T-1\), again proving
(APD.13).  Local constancy is only used to specify the original
cocycle space; no unstated topology on the coefficient field is used.
\end{proof}
% END APD6

\subsection{Realizing the original local deformation problem}

% BEGIN APD7
\paragraph{APD7. Residue constants in a general complete local ring.}
\label{claim:APD7}
Let \(\mathcal O\to C\) be a local homomorphism, with \(C\)
complete Noetherian local, and suppose its residue field \(k_C\)
is finite as an \(\mathcal O\)-module.  Let \(\kappa\) be
the residue field of \(\mathcal O\).  The subring
\[
 C_0=\{c\in C:\bar c\text{ belongs to the specified image of }
                         \kappa\text{ in }k_C\}
       =\operatorname{im}(\mathcal O)+\mathfrak m_C          \tag{APD.14}
\]
is an \(\mathrm{IsCO}(\mathcal O,C_0)\)-algebra.  The ring
\(C\) is finite as a \(C_0\)-module, and, as subsets of \(C\),
\[
 \mathfrak m_{C_0}=\mathfrak m_C,\qquad
 \mathfrak m_{C_0}^a=\mathfrak m_C^a\quad(a\ge1).
                                                               \tag{APD.15}
\]
For \(a=0\) the corresponding assertion is the intersection
identity \(C_0\cap C=C_0\).  This proves the full original
\texttt{cl\_\allowbreak sub}, without assuming that \(C\) is a DVR or that
its residue constants have a field section in \(C\).

\begin{proof}
Locality of the coefficient map gives an injection
\(\kappa\hookrightarrow k_C\).  The finite-module assumption
therefore says exactly that \([k_C:\kappa]\) is finite.
The inverse image in (APD.14) is a subring.  A member with nonzero
residue has an inverse in \(C\), whose residue is again in
\(\kappa\), so this inverse lies in \(C_0\).  This proves
locality, the asserted maximal ideal, and the prescribed residue
isomorphism.  For \(a\ge1\), a product of \(a\) elements
of \(\mathfrak m_C\), multiplied by an arbitrary coefficient
of \(C\), can absorb that coefficient into its first factor.
All the resulting factors lie in \(\mathfrak m_{C_0}\).
This proves (APD.15) in both directions.

Choose lifts \(b_1=1,b_2,\ldots,b_e\) of a
\(\kappa\)-basis of \(k_C\).  Every element of \(C\)
differs from a combination of the \(b_i\), with coefficients
lifted from \(\mathcal O\), by an element of
\(\mathfrak m_C\subset C_0\).  Thus the \(b_i\) generate
\(C\) over \(C_0\).  If \(x_j\) generate \(\mathfrak m_C\)
as a \(C\)-ideal, the finite family \(b_ix_j\) generates it
as a \(C_0\)-module.

Here is also a direct proof of Noetherianity of \(C_0\).
For an ideal \(I\subset C_0\), choose finitely many elements
of \(I\) generating \(CI\) as a \(C\)-ideal.  This makes
\(CI\) finite over \(C_0\).  The product
\(\mathfrak m_{C_0}CI\) is a finite \(C_0\)-module and is
contained in \(I\): indeed \(\mathfrak m_{C_0}C\subset C_0\).
The quotient \(CI/\mathfrak m_{C_0}CI\) is a finite-dimensional
\(\kappa\)-space.  Its subspace
\(I/\mathfrak m_{C_0}CI\) has a finite basis, and lifts of
that basis, together with generators of
\(\mathfrak m_{C_0}CI\), generate \(I\).  Hence every ideal
is finite.  Finally \(C_0\) is closed in \(C\), being the
inverse image of a subset of the discrete residue field.
Equation (APD.15) identifies its adic topology with the induced
one; completeness and separation follow.  All its coefficient
structures are restrictions of the original map.
\end{proof}
% END APD7

% BEGIN APD8
\paragraph{APD8. Broad residual points have one triangular direction.}
\label{claim:APD8}
Let \(K\) be any field, \(f:k\to K\), and let
\(r,\psi\) be an original normalized rank-two or rank-three
representation satisfying \(\mathrm{IsBar}(f)\).  This condition
prescribes the diagonal entries and ordinary characters; it does
not prescribe the off-diagonal entries.  Nevertheless either all
matrices \(r(g)\) are upper triangular or all are lower triangular.
Consequently a weight-zero monomial in their entries, containing
an off-diagonal factor, is zero.

\begin{proof}
The diagonal entries are characters
\(f(\Psi(g)^{w_i})\), where
\(w=(1,-1)\) or \((1,0,-1)\).  Comparing the \((0,0)\)
entry of \(r(gh)=r(g)r(h)\) in rank two gives
\(r(g)_{01}r(h)_{10}=0\) for all \(g,h\).  Over a field,
either the entire first function is zero or the entire second one
is zero.  These are the two alternatives.

In rank three use the normalized marked difference
\[
 Z=r(g_0)-r(g_0^*)=\operatorname{diag}(q,0,-q),
 \qquad q=f(\Psi(g_0)-\Psi(g_0)^{-1}).                       \tag{APD.16}
\]
The frame gives \(q\ne0\), and its odd characteristic gives
\(q\ne-q\).  For every \(i\), the diagonal character identity
applied to \(gh\) and to
\(gg_0h,gg_0^*h\) gives
\[
 \sum_{j\ne i}r(g)_{ij}r(h)_{ji}=0,
 \qquad
 \sum_{j\ne i}Z_{jj}r(g)_{ij}r(h)_{ji}=0.                 \tag{APD.17}
\]
There are only two summands, with distinct coefficients in the
second equation.  Subtraction after multiplying the first equation
therefore makes each summand zero.  In particular the pairs of
entry functions \((01,10)\), \((02,20)\), and \((12,21)\)
cannot both be nonzero.

Polarization gives \(r(g^*)_{01}=r(g)_{12}\) and
\(r(g^*)_{10}=r(g)_{21}\), so the two adjacent upper functions
vanish together, as do the two adjacent lower functions.  If
\(01\) vanishes, so does \(12\).  If \(02\) also vanishes
the representation is lower triangular.  Otherwise \(20\)
vanishes; the \((0,1)\) entry of the product identity, with
\(01=12=20=0\), gives
\(r(g)_{02}r(h)_{21}=0\).  It follows that \(21=10=0\),
so the representation is upper triangular.  The branch in which
\(10\) vanishes is the same argument with upper and lower
indices exchanged.  Thus one direction works for all group
elements at once.

The weights strictly decrease with the index.  In an upper
triangular monomial all factors that can be nonzero have
nonnegative weight, and every off-diagonal such factor has
positive weight.  A weight-zero monomial containing an
off-diagonal factor therefore has a zero factor.  The lower
triangular case reverses the signs.  This proof works over every
specified target field and divides only by nonzero differences
and by the integer \(2\); the allowed Host case \(\ell=3\)
is unchanged.
\end{proof}
% END APD8

% BEGIN APD9
\paragraph{APD9. The graded Laurent vertex is an actual separated algebra.}
\label{claim:APD9}
Suppose \(C\) is an \(\mathrm{IsCO}(O,C)\)-algebra,
\(C\hookrightarrow K\) is an injective map of commutative
rings with the specified \(O\)-scalar tower, and
\(\sigma=(r,\psi):\mathrm{NRep}(F,n,K)\) satisfies the original
\(\mathrm{Cone}(O,C,\sigma)\).  Explicitly, its diagonal
entries and full ordinary character values lie in \(C\) with
the prescribed residues; every weight-zero entry monomial with
an off-diagonal factor lies in \(\mathfrak m_C\); each entry
span \(M_{ij}=\langle r(g)_{ij}:g\in G_0\rangle_C\) is
finite; the entries are locally constant modulo
\(\mathfrak m_C^aM_{ij}\); and the character functions in
\(C\) are adically continuous.

Let \(P_j\subset K\) be the \(C\)-span of weight-\(j\)
entry monomials and define
\[
 D=\bigoplus_{j\in\mathbb Z}P_jt^j\subset K[t,t^{-1}],
 \qquad
 \mathfrak n=\{d\in D:[t^0]d\in\mathfrak m_C\}.            \tag{APD.18}
\]
Then \(D\) is Noetherian, \(\mathfrak n\) is a maximal ideal,
\(D/\mathfrak n=\kappa=k\), and
\(\bigcap_a\mathfrak n^a=0\).  There is a normalized
representation \(\sigma_D\) over \(D\) with
\[
 r_D(g)_{ij}=t^{w_i-w_j}r(g)_{ij},\qquad
 \psi_{D,v,i}(s)=\psi_{v,i}(s).                            \tag{APD.19}
\]
It is \(\mathfrak n\)-adically continuous and has the original
diagonal residual lift.

\begin{proof}
The products of the \(P_j\)'s satisfy \(P_iP_j\subset P_{i+j}\).
The degree-zero part is exactly \(C\): a diagonal-only monomial
lies in \(C\), and any other degree-zero monomial lies in
\(\mathfrak m_C\) by the cone condition.  For \(j\ne0\),
a product from \(P_jP_{-j}\) has an off-diagonal factor, so
lies in \(\mathfrak m_C\).  Thus taking constant coefficient
modulo \(\mathfrak m_C\) is a ring homomorphism
\(D\to k\), with kernel exactly \(\mathfrak n\).

For each pair \((i,j)\), choose finitely many generators of
\(M_{ij}\) and multiply them by \(t^{w_i-w_j}\).  This finite
family generates \(D\) as a \(C\)-algebra: every weighted
entry is in its span, every weighted monomial is a product of
these entries, and a Laurent polynomial has finite support.
Hilbert's basis theorem makes \(D\) Noetherian.

We spell out separation, since the vertex ideal need not be in
the Jacobson radical of the whole Laurent subalgebra.  It is a
homogeneous ideal, so its powers and their intersection
\(J=\bigcap_a\mathfrak n^a\) are homogeneous.  Artin--Rees
for the finite ideal \(J\subset D\) gives \(J=\mathfrak nJ\).
The determinant trick applied to finite generators of \(J\)
then gives \(b\in\mathfrak n\) with \((1-b)J=0\): writing
each generator as an \(\mathfrak n\)-linear combination of
the generators, the determinant of the corresponding
identity-minus-coefficient matrix is \(1-b\) and annihilates
them.  If \(at^j\in J\) is homogeneous, the degree-\(j\)
part of \(b(at^j)=at^j\) says \(ca=a\), where
\(c=[t^0]b\in\mathfrak m_C\).  The image of \(1-c\) is
a unit in \(K\), hence \(a=0\).  Every homogeneous part of
every element of \(J\) is zero, so \(J=0\).

Formula (APD.19) is conjugation by
\(\operatorname{diag}(t^{w_i})\) in \(K[t,t^{-1}]\).
Its inverse matrices come from \(g^{-1}\), and each ordinary
character has its inverse at \(s^{-1}\) already in \(C\).
The normalized marking, polarization, and the ordered ordinary
product identities therefore hold in \(D\); the relation
\(w_{i^{\mathrm{rev}}}=-w_i\) checks polarization explicitly.
An entry difference in \(\mathfrak m_C^aM_{ij}\), after
weighting, is a sum of elements of \(\mathfrak n^a\) times
weighted generators of \(M_{ij}\).  The cone continuity
condition consequently gives continuity modulo every
\(\mathfrak n^a\).  Characters use the same argument in
degree zero.  Every off-diagonal weight is nonzero, so its
weighted entry lies in \(\mathfrak n\); the diagonal and
character residues are exactly the specified ones.
\end{proof}
% END APD9

% BEGIN APD10
\paragraph{APD10. Cone points and broad local points are realized.}
\label{claim:APD10}
At the hypotheses of APD9 there are actual \(O\)-algebra maps
\[
 \xi:\mathcal A\longrightarrow K,
 \qquad\beta:\mathcal B\longrightarrow C,
 \qquad \xi_*u=\sigma,\qquad
 \xi(\mathrm{inc}(b))=\beta(b)\text{ in }K.                \tag{APD.20}
\]
This proves full original \texttt{alg}.  Consequently, for every
complete Noetherian local \(O\)-algebra \(C\) satisfying the
original \(\mathrm{IsCL}(O,C)\), every compatible residue map
\(f:k\to k_C\), and every adically continuous normalized
\(\sigma\) with \(\bar\sigma.\mathrm{IsBar}(f)\), there is
\[
 \xi:\mathcal A\longrightarrow C,
 \qquad \xi_*u=\sigma,
 \qquad \xi(\mathfrak m_{\mathcal B}\mathcal A)
                                  \subset\mathfrak m_C. \tag{APD.21}
\]
Thus the full original \texttt{point\_\allowbreak loc} is supplied, with its
broader residual condition intact.

\begin{proof}
Use the vertex algebra of APD9.  Each
\(D_a=D/\mathfrak n^{a+1}\) is Noetherian local, has nilpotent
maximal ideal, and has the prescribed residue field \(k\).
It is therefore a complete coefficient algebra over \(O\).
The universal property included in \(S\) gives a unique map
\(f_a:S.R\to D_a\) carrying the universal representation to
\(\sigma_D\bmod\mathfrak n^{a+1}\).  Uniqueness makes these
maps compatible.  They kill \(\mathfrak m_{S.R}^{a+1}\), since
coefficient-algebra maps are local.

Let \(\widehat D=\varprojlim D_a\).  The resulting map
\(f:S.R\to\widehat D\) sends every matrix entry to its actual
element of \(D\).  It sends every weight-zero monomial and
ordinary character value to \(C\).  We must also handle their
closure defining \(\mathcal B\).  The complete Noetherian ring
\(C\) has finite residue field, hence each
\(C/\mathfrak m_C^a\) is finite and \(C\) is compact in its
adic topology.  Its map to the Hausdorff ring \(\widehat D\)
is continuous and injective, the latter by APD9.  Its image is
therefore closed.  Approximating an element of \(\mathcal B\)
by polynomials in its defining generators, and using that \(f_a\)
kills \(\mathfrak m_{S.R}^{a+1}\), shows that its image lies
in this closed copy of \(C\).  This is a concrete coefficient
limit argument; it does not presume that \(D\) is complete.

Since \(\mathcal A\) is generated by \(\mathcal B\) and
the actual matrix entries, \(f(\mathcal A)\subset D\).
Separation of \(D\) makes this a unique algebra map into
\(D\), and its restriction to \(\mathcal B\) gives \(\beta\).
Evaluation \(t=1\) is a ring homomorphism on finite Laurent
polynomials and sends (APD.19) to \(\sigma\).  Composing
gives (APD.20).  In particular no evaluation of an unrestricted
completed Laurent series has been used.

For (APD.21), take \(C_0\subset C\) from APD7.  The
\(\mathrm{IsCL}\) hypothesis supplies precisely its finite
residue-extension premise.  The residual diagonal and character
conditions put all those values, and their specified residues,
in \(C_0\).  APD8 puts each weight-zero monomial containing an
off-diagonal entry in \(\mathfrak m_C=\mathfrak m_{C_0}\).
Each entry span over \(C_0\) is finite because \(C_0\) is
Noetherian and \(C\) is finite over it.

For clarity, its required entry-span continuity does not follow
just by renaming the topology.  If
\(M\subset C\) is one of the finite entry spans, Artin--Rees
(ALC6) gives, for each \(a\), an integer \(b\) with
\[
 M\cap\mathfrak m_{C_0}^{b}C
                         \subset\mathfrak m_{C_0}^{a}M.
                                                               \tag{APD.22}
\]
An entry difference in \(\mathfrak m_C^b\) belongs to the
left side, by (APD.15), so adic continuity in \(C\) gives the
cone's exact quotient-module continuity.  Character continuity
in \(C_0\) follows directly from (APD.15), including their
inverse values.  All five original cone clauses are now proved.
Apply (APD.20) with \(C_0\hookrightarrow C\).
The map \(\beta:\mathcal B\to C_0\) preserves residue
maps, since RTC7 makes \(\mathcal B\) an
\(\mathrm{IsCO}(O,-)\)-algebra.  It sends
\(\mathfrak m_{\mathcal B}\) into \(\mathfrak m_{C_0}\).
This proves the last assertion of (APD.21) on the original
inclusion \(\mathcal B\to\mathcal A\).
\end{proof}
% END APD10

% BEGIN APD11
\paragraph{APD11. Entries determine the map on the actual entry algebra.}
\label{claim:APD11}
Let \(C\) be any Noetherian local \(O\)-algebra; it need not
be complete.  If two maps \(\xi,\xi':\mathcal A\to C\)
carry \(u\) to the same normalized representation and both send
\(\mathfrak m_{\mathcal B}\mathcal A\) into \(\mathfrak m_C\),
then \(\xi=\xi'\).  This is the full original \texttt{ext\_\allowbreak A}.

\begin{proof}
The maps agree on matrix entries, ordinary character values and
their inverses, and hence on the algebra generated by weight-zero
entry monomials and character values.  Fix \(b\in\mathcal B\)
and \(a\ge1\).  RTC9's fixed-weight comparison at weight zero
gives \(b_0\) in that generating algebra with
\(b-b_0\in\mathfrak m_{\mathcal B}^{a}\mathcal B\): choose
the required ambient \(\mathfrak m_{S.R}\)-adic approximation
and then apply the comparison.  Both maps send this difference
into \(\mathfrak m_C^a\).  Their values on \(b_0\) agree,
so \(\xi(b)-\xi'(b)\in\mathfrak m_C^a\) for every \(a\).
Krull intersection in the Noetherian local ring \(C\) makes
the difference zero.  They thus agree on \(\mathcal B\) and
on all matrix-entry generators of \(\mathcal A\), proving
the assertion.  This uses the actual fixed-weight topology; no
continuity of an arbitrary map out of \(S.R\) is assumed.
\end{proof}
% END APD11

% BEGIN APD12
\paragraph{APD12. The actual local representation is universal.}
\label{claim:APD12}
For the supplied \(L\) in (APD.1), let \(C\) be every
\(\mathrm{IsCO}(O',C)\)-algebra and let \(\sigma\) be every
continuous normalized representation with
\[
 \operatorname{res}_{C,*}\sigma
                     =\operatorname{res}_{R,*}\rho
                    \quad\text{over }k'=\operatorname{Res}(O').
                                                               \tag{APD.23}
\]
There is a unique \(O'\)-algebra map \(g:R\to C\) with
\(g_*\rho=\sigma\).  Together with ALC17's continuity,
this is the full original \texttt{def\_\allowbreak iso} and the literal
\(\rho.\mathrm{Univ}(O')\) required by \texttt{pres}.

\begin{proof}
The map \(O\to O'\) is local: the original unramified
identity \(\mathfrak m_OO'=\mathfrak m_{O'}\) proves this
directly.  Its composite to \(C\) is therefore local.
The ring \(C\) is complete Noetherian, and its residue field
is \(k'\), finite over \(O\) because \(O'\) is a finite
\(O\)-module.  Thus it satisfies the precise
\(\mathrm{IsCL}(O,C)\) input of APD10.

The compatible maps \(k\to k'\to\operatorname{Res}(C)\)
and \(k\to k'\to\operatorname{Res}(R)\) are those induced
by the specified scalar towers.  ALC17 gives the residual
\(\mathrm{IsBar}\) condition for \(\rho\).  Composing
(APD.23) with the second residue identification gives the same
condition for \(\sigma\).  APD10 produces
\(\xi:\mathcal A\to C\), with \(\xi_*u=\sigma\), sending
\(x_0=\mathfrak m_{\mathcal B}\mathcal A\) into
\(\mathfrak m_C\).

Compare \(\operatorname{res}_C\xi\) and
\(\operatorname{res}_R\iota\) as maps \(\mathcal A\to k'\).
Their normalized representations agree by (APD.23).  Both kill
\(x_0\): for the latter this follows from
\(x_0\subset x=\iota^{-1}(\mathfrak m_R)\), proved in
ALC11.  APD11 therefore says these two residue maps are equal.
The original universal field of the supplied \(L\) now gives
a unique \(O'\)-map \(g:R\to C\) extending \(\xi\).
It sends \(\rho\) to \(\sigma\).

If \(g'\) is any other map sending the representation correctly,
the maps \(g'\iota\) and \(\xi\) agree on the representation
and both send \(x_0\) into \(\mathfrak m_C\), because
\(g'\) is a coefficient-algebra map.  APD11 makes them equal,
and the uniqueness field of \(L\) gives \(g'=g\).
The same \(O',R,\iota,x\) was fixed before \(C\) and
\(\sigma\) were chosen.  No diagonal residual hypothesis was
added to (APD.23).
\end{proof}
% END APD12

\subsection{Continuous boxes and square-zero corrections}

\paragraph{Scope and notation for the deformation lemmas.}
In APD13--APD20 the coefficient ring \(\mathcal O\) is an
arbitrary local ring, \(\kappa=\operatorname{Res}(\mathcal O)\),
and the rank \(n\) is arbitrary unless explicitly restricted.
The frame is any original \(F\) over a field of characteristic
\(\ell\), with \(F.\mathrm{Ok}(C_0)\).  This broader scope is
the scope of the original matrix lemmas.  It is not limited to
the arithmetic \(O',R\) of (APD.1).
Put
\[
 \theta(A)=J_n(A^{-1})^tJ_n,\qquad
 \theta_{\mathrm{Lie}}(X)=-J_nX^tJ_n,
 \qquad \mathrm{GG}_n(C)=\mathrm{GL}_n(C)\rtimes\{\pm1\},
                                                               \tag{APD.24}
\]
where the negative sign acts by \(\theta\).  A box
\(\mathfrak B\) consists of a homomorphism
\(b:\Gamma\to\mathrm{GG}_n(C)\), with sign \(F.\delta\),
and matrices \(g_v\) making the restriction to each \(G_v\)
upper triangular.  It is \emph{of} a normalized representation
\(\rho\) if its matrix at \(c\) is \(1\), its restriction
to \(G_0\) is \(\rho.r\), and the diagonal in the \(g_v\)
basis is the specified ordered character family \(\rho.\psi\).
This is exactly \(\mathrm{Box.Of}\), not a change of deformation
problem.  On \(W=M_n(\kappa)\) its residual adjoint action is
\[
 \overline{\operatorname{Ad}}(s)X=
 \begin{cases}
  \bar b(s)X\bar b(s)^{-1},&\delta(s)=1,\\
  \bar b(s)\theta_{\mathrm{Lie}}(X)\bar b(s)^{-1},&\delta(s)=-1,
 \end{cases}
 \quad
 \mathfrak b_v=\{X:\bar g_vX\bar g_v^{-1}
                                    \text{ is upper triangular}\}.
                                                               \tag{APD.25}
\]
Here \(\bar b(s)\) denotes the matrix part.  Bars on boxes
and representations use the specified residue map of the
coefficient algebra.

% BEGIN APD13
\paragraph{APD13. Boxes are constructed and recovered with their characters.}
\label{claim:APD13}
For every commutative ring \(C\), ideal \(J\subset C\), and
\(J\)-adically continuous normalized representation \(\rho\)
with an ordered ordinary flag at every \(v\), there is a box
\(\mathfrak B\) of \(\rho\), with \(J\)-adically continuous
matrix entries on \(\Gamma\).  Conversely a continuous box
with matrix \(1\) at \(c\) and diagonal marked difference
comes from a continuous normalized representation, with its full
ordered ordinary characters.  These are full original
\texttt{pr\_\allowbreak box} and \texttt{pr\_\allowbreak ofbox}.

\begin{proof}
Since \(c^2=1\) and \(\delta(c)=-1\), every \(s\in\Gamma\)
has a unique expression \(h\) or \(hc\), with \(h\in G_0\).
Define the matrix part of \(b(s)\) to be \(\rho.r(h)\) and
its sign to be \(\delta(s)\).  Polarization gives
\(\rho.r(chc^{-1})=\theta(\rho.r(h))\); this proves the
semidirect-product homomorphism law in the four possible sign
cases.  Its value at \(c\) has matrix \(1\).  Choose the
given ordinary bases \(g_v\); they give the required local
upper triangularity and ordered diagonal entries.

The map \(s\mapsto h\) is continuous: it is \(s\) on the
open subgroup \(G_0\) and \(sc\) on its open other coset.
Thus continuity of \(\rho\) proves the stated entry continuity
of \(b\), at every power of \(J\).

Conversely restrict the matrix part of the box to \(G_0\),
where the sign is positive.  Conjugate locally by \(g_v\),
and take its diagonal entries as characters.  Diagonal entries
of a product of upper triangular matrices multiply, and those
of the inverse are inverse units, so these really are characters.
An upper triangular matrix with its \(i\)-th diagonal entry
subtracted takes the \((i+1)\)-st standard flag space into the
\(i\)-th.  Starting with the rightmost factor proves the
ordered ordinary-product identity for every tuple of local
group elements.  The normalization at \(c\) gives
\(r(h^*)=J_nr(h)^tJ_n\), and the given marked difference
gives the remaining normalization.  Restriction and finite matrix
sums prove continuity of both matrix entries and character values.
Reduction and arbitrary coefficient maps preserve all of these
constructions.
\end{proof}
% END APD13

% BEGIN APD14
\paragraph{APD14. Continuous set-theoretic lifts across every small extension.}
\label{claim:APD14}
Let \(p:C'\twoheadrightarrow C\) be a surjection of
\(\mathrm{IsCO}(\mathcal O,-)\)-algebras, with
\(K=\ker p\) and \(\mathfrak m_{C'}K=0\).  Every
adically continuous function from any topological space into
\(C\) has an adically continuous lift to \(C'\).
Consequently every continuous box \(\mathfrak B\) over \(C\)
has the original \(\mathrm{SetLift}(p,\mathfrak B)\): continuous
matrix lifts \(r_s\) of its matrices, lifts \(g'_v\) of its
ordinary bases, and continuous upper triangular lifts \(t_{v,s}\)
of \(g_vb(\operatorname{loc}_v s)g_v^{-1}\).
This proves full \texttt{pr\_\allowbreak sec} and \texttt{pr\_\allowbreak setlift}.

\begin{proof}
The map preserves residues, so
\(p^{-1}(\mathfrak m_C)=\mathfrak m_{C'}\) and
\(p(\mathfrak m_{C'}^a)=\mathfrak m_C^a\) for every \(a\).
Artin--Rees applied to \(K\subset C'\), together with
\(\mathfrak m_{C'}K=0\), gives an integer \(a_0\ge1\)
with \(K\cap\mathfrak m_{C'}^{a_0}=0\).  Thus
\(p:\mathfrak m_{C'}^{a_0}\to\mathfrak m_C^{a_0}\) is an
additive bijection.  Choose a representative in \(C'\) for
each class of \(C/\mathfrak m_C^{a_0}\).  Lift a general
\(z\in C\) by its chosen representative plus the unique
element of \(\mathfrak m_{C'}^{a_0}\) lifting the remaining
difference.  This defines a section \(s:C\to C'\).

If \(z-z'\in\mathfrak m_C^a\) with \(a\ge a_0\), their
chosen representatives coincide.  The difference \(s(z)-s(z')\)
belongs to \(\mathfrak m_{C'}^{a_0}\) and lifts \(z-z'\).
There is also a lift in \(\mathfrak m_{C'}^a\); uniqueness
in \(\mathfrak m_{C'}^{a_0}\) makes these equal.  Hence
\(s(z)-s(z')\in\mathfrak m_{C'}^a\).  Applying this to a
continuous function at level \(\max(a,a_0)\) proves its lifted
continuity at level \(a\).  No compactness or profiniteness of
its domain is required.

Apply these lifts entry by entry to an invertible matrix-valued
function.  Its lifted determinant is a unit, since its image is
a unit under the local map \(p\), so the lifted matrix is
invertible.  If an entry function is identically zero, choose
its lift identically zero.  This preserves upper triangularity
for the local matrices.  Lift each fixed \(g_v\) in the same
way.  The local conjugated matrices are continuous because the
local group maps and finite matrix operations are continuous
modulo every power.  These choices give every field of
\(\mathrm{SetLift}\).  Neither the lifted \(r_s\) nor the
lifted \(t_{v,s}\) is asserted to be a homomorphism.
\end{proof}
% END APD14

% BEGIN APD15
\paragraph{APD15. The two explicit defects are continuous cone cocycles.}
\label{claim:APD15}
Let \(p:C'\to C\) be any map of coefficient algebras with
\(\mathfrak m_{C'}K=0\), and suppose a set lift
\(D=(r_s,g'_v,t_{v,s})\) of a box is supplied.  Define
\[
 \begin{split}
 \Omega(s,t)&=r_s\,\theta^{\delta(s)}(r_t)r_{st}^{-1}-1,\\
 A_v(s)&=r_{\operatorname{loc}_v s}
               ((g'_v)^{-1}t_{v,s}g'_v)^{-1}-1,
 \end{split}                                                  \tag{APD.26}
\]
where \(\theta^{+1}\) means the identity and
\(\theta^{-1}\) means \(\theta\).  All their entries belong
to \(K\), and both functions are locally constant.  For every
\(\mathcal O\)-linear \(u:K\to\kappa\), the entrywise
images \((\omega_u,a_{u,v})\) satisfy
\[
 \begin{split}
 \overline{\operatorname{Ad}}(s)\omega_u(t,r)+\omega_u(s,tr)
                 &=\omega_u(st,r)+\omega_u(s,t),\\
 \omega_u(\operatorname{loc}_v s,\operatorname{loc}_v t)
      -d a_{u,v}(s,t)&\in\mathfrak b_v,\\
 d\beta(s,t)&=\overline{\operatorname{Ad}}(s)\beta(t)
                                  -\beta(st)+\beta(s).
 \end{split}                                                  \tag{APD.27}
\]
The assignment \(u\mapsto(\omega_u,a_u)\) is linear.
This proves full \texttt{pr\_\allowbreak omk},
\texttt{borel\allowbreak Projection\_\allowbreak is\allowbreak Two\allowbreak Cocycle}, and \texttt{pr\_\allowbreak cocv}.

\begin{proof}
Mapping (APD.26) by \(p\) gives zero, by the homomorphism law
of the box and the three defining map identities of a set lift.
Also \(K\subset\mathfrak m_{C'}\), so \(K^2=0\).
Artin--Rees as in APD14 gives \(a_0\) with
\(K\cap\mathfrak m_{C'}^{a_0}=0\).  The functions in (APD.26)
are continuous modulo every power, since inversion in the reduced
matrix group is a function of the reduced invertible matrix and
\(\delta\) is locally constant.  Equality modulo
\(\mathfrak m_{C'}^{a_0}\) for two values in \(M_n(K)\)
is actual equality.  Their functions are therefore locally
constant, without any assumption that \(\kappa\) is finite.

The action of a coefficient \(c\in C'\) on \(K\) is its
residue action: subtract a residue lift from \(\mathcal O\)
and use \(\mathfrak m_{C'}K=0\).  Thus \(u(cy)=\bar c\,u(y)\).
For matrices \(X,Y\) with entries in \(K\),
\((1+X)(1+Y)=1+X+Y\) and \((1+X)^{-1}=1-X\).
Conjugation by a lifted signed matrix acts on these entries
through (APD.25), since
\(\theta(1+X)=1+\theta_{\mathrm{Lie}}(X)\).

Put \(x_s=(r_s,\delta(s))\).  Associativity gives the exact
identity
\[
 (x_sx_tx_{st}^{-1})(x_{st}x_rx_{str}^{-1})
 =x_s(x_tx_rx_{tr}^{-1})x_s^{-1}(x_sx_{tr}x_{str}^{-1}).
\]
Substitution of \(1+\Omega\), followed by \(K^2=0\), gives
the first equation of (APD.27), after applying \(u\).
For the local equation put
\(T_v(s,t)=t_{v,s}t_{v,t}t_{v,st}^{-1}-1\).
It has entries in \(K\) and is upper triangular.  Substitution
of the definitions and cancellation of the invertible factors gives
\[
 \Omega(\operatorname{loc}_v s,\operatorname{loc}_v t)
 -\bigl(r_{\operatorname{loc}_v s}A_v(t)
             r_{\operatorname{loc}_v s}^{-1}
                 -A_v(st)+A_v(s)\bigr)
                 =(g'_v)^{-1}T_v(s,t)g'_v.                \tag{APD.28}
\]
One can check this without commuting matrices: multiply
\((1+\Omega)(1+A_v(st))\), replace the lifted local matrices
by \((1+A_v(s))(g'_v)^{-1}t_{v,s}g'_v\), and cancel.
All products of two kernel matrices disappear.  Applying \(u\)
to (APD.28) makes its right side an element of
\(\mathfrak b_v\), proving the second equation.  Finally all
applications of \(u\) are to actual kernel entries; their
linearity proves linearity of the entire assignment.
\end{proof}
% END APD15

% BEGIN APD16
\paragraph{APD16. A vanishing two-cocycle gives the actual corrected box.}
\label{claim:APD16}
Let \(\mu:\kappa\to C'\) be an additive map satisfying
\[
 \mu(\bar c a)=c\mu(a),\qquad
 \mu(\kappa)\subset\mathfrak m_{C'},\qquad p\mu=0.
                                                               \tag{APD.29}
\]
Here \(p:C'\to C\) is any map between the original
\(\mathrm{IsCO}\)-algebras, and the set lift is any original
\(\mathrm{SetLift}(p,\mathfrak B)\).  We do not require
\(\mathfrak m_{C'}\ker p=0\), nor surjectivity of \(p\).
Only the image of \(\mu\) must be square zero, as follows
from (APD.29).  Suppose \(\beta:\Gamma\to W\)
is locally constant, \(Y_v\in W\), and functions \(a_v\)
satisfy
\[
 \Omega=\mu(d\beta),\qquad A_v=\mu(a_v),\qquad
 a_v(s)-\beta(\operatorname{loc}_v s)
       +(\overline{\operatorname{Ad}}(\operatorname{loc}_v s)-1)Y_v
                                                        \in\mathfrak b_v.
                                                               \tag{APD.30}
\]
Entrywise application of \(\mu\) is understood.  The formulas
\[
 b'(s)=\bigl(1-\mu(\beta(s))\bigr)(r_s,\delta(s)),
 \qquad g'_v{}^{\mathrm{new}}=g'_v(1-\mu(Y_v))             \tag{APD.31}
\]
define a continuous box over \(C'\) whose image is the given box
over \(C\).  These are the actual corrections in full original
\texttt{pr\_\allowbreak sqz} and \texttt{pr\_\allowbreak correct}.

\begin{proof}
Equation (APD.29) implies \(\mu(a)\mu(b)=0\), since the
residue of \(\mu(a)\) is zero.  Consequently
\(E(X)=1+\mu(X)\) is invertible,
\(E(X)E(Y)=E(X+Y)\), and \(E(X)^{-1}=E(-X)\).
For any signed invertible matrix \(z\), (APD.29) gives
\[
 zE(X)z^{-1}=E(\operatorname{Ad}(\bar z)X).                \tag{APD.32}
\]
This follows from matrix multiplication for positive sign; for
negative sign use \(\theta(E(X))=E(-J_nX^tJ_n)\).
The same coefficient calculation gives, for every matrix \(A\)
over \(C'\),
\(A\mu(X)=\mu(\bar A X)\) and
\(\mu(X)A=\mu(X\bar A)\), by expanding each matrix entry.
These identities and (APD.32) use only the first two conditions
of (APD.29).  They therefore give all outputs of
\texttt{pr\_\allowbreak sqz} for any \(\mathrm{IsCO}(\mathcal O,C')\),
any rank, and any original \(\mathrm{IsSq}\) map, without a
frame, a map \(p\), or its kernel conditions.

Multiply the two expressions for \(b'(s)\) and \(b'(t)\)
using (APD.32).  Their correction exponent is
\(-\beta(s)-\overline{\operatorname{Ad}}(s)\beta(t)
 +d\beta(s,t)=-\beta(st)\).  Thus \(b'\) is a group
homomorphism, with the unchanged sign character.  For a local
element, conjugating by the new \(g'_v\) gives
\[
 E\bigl(\bar g_v[
     a_v(s)-\beta(\operatorname{loc}_v s)
       +(\overline{\operatorname{Ad}}(\operatorname{loc}_v s)-1)Y_v
                ]\bar g_v^{-1}\bigr)t_{v,s}.              \tag{APD.33}
\]
The bracket condition in (APD.30) says its exponential is upper
triangular.  The product is therefore upper triangular as well.
Both the representation and bases reduce to the given ones because
\(p\mu=0\).  The new matrix entries are finite sums of
continuous entries times values of the locally constant \(\beta\).
They are continuous modulo every maximal-ideal power.  This
verifies every box field and the exact formula (APD.31), including
its signs.
\end{proof}
% END APD16

% BEGIN APD17
\paragraph{APD17. Both normalization equations are solved explicitly.}
\label{claim:APD17}
Let \(p:C'\to C\) be any map of coefficient algebras with
\(\mathfrak m_{C'}\ker p=0\).  Let \(\rho\) be normalized
over \(C\), and suppose its residue satisfies
\(\rho.\mathrm{Good}\): \(2\ne0\), the diagonal entries
of its marked difference are pairwise distinct, and each local
character list has a separating group element.  If a continuous
box \(\mathfrak B'\) maps to a box of \(\rho\), then there
are \(h\in\mathrm{GL}_n(C')\), with \(p(h)=1\), and a
continuous normalized \(\rho'\) over \(C'\), with
\(p_*\rho'=\rho\), such that the conjugated box has matrix
\(1\) at \(c\) and restricts to \(\rho'\) on \(G_0\).
This proves full \texttt{pr\_\allowbreak norm1}, \texttt{pr\_\allowbreak norm2}, and
\texttt{pr\_\allowbreak norm}; no surjectivity of \(p\) is required.
More precisely, the two elementary assertions below also allow
an arbitrary target \(\mathcal O\)-algebra \(C\), with no
\(\mathrm{IsCO}\) structure: assume explicitly
\(\ker p\subset\mathfrak m_{C'}\).  The first takes any
signed involution of negative sign with matrix reducing to
\(1\) and \(2\ne0\) in \(\kappa\).  The second takes
any \(Z\) fixed by \(\theta_{\mathrm{Lie}}\), with all
off-diagonal entries in \(\ker p\) and pairwise distinct
residual diagonal entries.  Both retain
\(\mathfrak m_{C'}\ker p=0\).  These are the literal broader
premises of \texttt{pr\_\allowbreak norm1} and \texttt{pr\_\allowbreak norm2}.

\begin{proof}
Write \(K=\ker p\).  It lies in \(\mathfrak m_{C'}\), so
\(K^2=0\); moreover \(2\) is a unit in \(C'\).  If the
matrix at \(c\) is \(1+X\), then \(X\in M_n(K)\).
The involution identity at \(c\) gives
\(J_nX^tJ_n=X\).  Put \(h_1=1-X/2\).  Its inverse is
\(1+X/2\), and signed conjugation gives
\[
 h_1(1+X)\theta(h_1^{-1})
                 =(1-X/2)(1+X)(1-X/2)=1.                 \tag{APD.34}
\]
Thus the first normalization is solved by an explicit kernel
matrix, with no lifting of an orthogonal group assumed.

In the resulting box put \(Z=b(g_0)-b(g_0^*)\).  Polarization
now gives \(\theta_{\mathrm{Lie}}(Z)=Z\).  Its off-diagonal
entries lie in \(K\), and the differences \(Z_{ii}-Z_{jj}\)
are units for \(i\ne j\), by the residual separation condition.
Define
\[
 A_{ii}=0,\qquad A_{ij}=Z_{ij}(Z_{ii}-Z_{jj})^{-1}\ (i\ne j),
 \qquad h_2=1+A.                                         \tag{APD.35}
\]
Every entry of \(A\) belongs to \(K\), so \(A^2=0\).
The relation \(Z_{j^{\mathrm{rev}},i^{\mathrm{rev}}}=-Z_{ij}\)
gives \(A_{j^{\mathrm{rev}},i^{\mathrm{rev}}}=-A_{ij}\);
therefore \(\theta(h_2)=h_2\).  This conjugation preserves
the matrix \(1\) at \(c\).  For \(i\ne j\), expansion gives
\[
 ((1+A)Z(1-A))_{ij}
  =Z_{ij}+A_{ij}Z_{jj}-Z_{ii}A_{ij}=0.                    \tag{APD.36}
\]
The omitted terms each contain two entries of \(K\) and are
zero; this explains why the formula works over an arbitrary
coefficient algebra, not just a field.  It proves the second
normalization equation.

Conjugate each local basis simultaneously by replacing it with
\(g'_v h_1^{-1}h_2^{-1}\).  Upper triangularity and its
ordered diagonal characters are thereby preserved.  APD13 now
recovers a normalized \(\rho'\).  All conjugating matrices
reduce to \(1\), so after applying \(p\) this box is the
original one; its restricted representation and diagonal
characters are precisely \(\rho\).  Constant conjugations
and the locally constant sign character preserve all adic
continuity assertions.  The two elementary calculations also
prove \texttt{pr\_\allowbreak norm1} and \texttt{pr\_\allowbreak norm2} at their
broader literal premises: they only need the stated kernel,
involution and diagonal conditions, respectively.
\end{proof}
% END APD17

% BEGIN APD18
\paragraph{APD18. The continuous cone is the original obstruction space.}
\label{claim:APD18}
Let \(R\) be an \(\mathrm{IsCO}(\mathcal O,R)\)-algebra,
let \(\rho.\mathrm{Univ}(\mathcal O)\) hold, and let
\(\mathfrak B\) be a continuous box of \(\rho\) with good
residue.  Define \(Z_C^2\) by local constancy and the two
equations (APD.27).  Its coboundaries are the pairs
\((d\beta,a)\), where \(\beta\) is locally constant and
there are \(Y_v\) with
\[
 a_v(s)-\bigl(\beta(\operatorname{loc}_v s)
       -(\overline{\operatorname{Ad}}(\operatorname{loc}_v s)-1)Y_v\bigr)
                                                    \in\mathfrak b_v.
                                                               \tag{APD.37}
\]
Put \(H_C^2=Z_C^2/(Z_C^2\cap B_C^2)\), exactly the original
\(\mathrm{coneH2}\).  Then
\(\mathrm{HasObs}(\mathcal O,R,H_C^2)\) holds in the full
sense of APD5.  Its vanishing-to-section assertion is also the
full original \texttt{pr\_\allowbreak obs}.

\begin{proof}
For every original small surjection \(p:C'\to R\), choose
a set lift by APD14.  APD15 gives the linear map
\[
 \mathrm{ob}_p:K^\vee\longrightarrow H_C^2,
 \qquad u\longmapsto[(\omega_u,a_u)].                     \tag{APD.38}
\]
Suppose its value at \(u\) is zero, and let
\(C'\overset q\twoheadrightarrow C''\overset{p''}\to R\)
be any factorization permitted in APD5.  The map \(q|_K\)
factors through \(u\).  If \(u=0\), it is zero and we set
\(\mu=0\).  Otherwise \(u\) is surjective onto the
one-dimensional \(\kappa\)-space, and define
\(\mu(a)=q(y)\) whenever \(u(y)=a\).  The kernel condition
on \(q\) makes this independent of \(y\) and additive.
For every coefficient \(c''\in C''\), lift it by \(q\)
and reduce that lift modulo \(\mathfrak m_{C'}\).
The equality \(\mathfrak m_{C'}K=0\) then gives
\(\mu(\overline{c''}a)=c''\mu(a)\).
Its image lies in \(\mathfrak m_{C''}\), is contained in
\(q(K)\), and is killed by \(p''\).  This proves every
clause of the original \texttt{pr\_\allowbreak mu} construction.
Indeed the construction only uses an ideal
\(K\subset\mathfrak m_{C'}\) annihilated by
\(\mathfrak m_{C'}\), the surjection \(q\), and its stated
condition on \(\ker u\).  Without mentioning \(p,p''\), it
returns \(\mu(u(y))=q(y)\) and
\(\mu(\kappa)\subset q(K)\) at the full original general-ideal
scope of \texttt{pr\_\allowbreak mu}.

Push the chosen set lift through \(q\).  Its defects are
\(\mu(\omega_u)\) and \(\mu(a_u)\), because all original
defect entries are in \(K\).  Vanishing in the stated quotient
means there are actual \(\beta,Y\) satisfying (APD.37), with
\(\beta\) locally constant.  APD16 corrects the pushed set
lift to a continuous box.  Furthermore
\(\ker p''=q(K)\) and
\(\mathfrak m_{C''}\ker p''=0\), using surjectivity of \(q\).
APD17 therefore normalizes this box to a continuous
\(\rho''\) with \(p''_*\rho''=\rho\).

Residue preservation makes \(\bar\rho''=\bar\rho\).
Universality gives \(s:R\to C''\) with \(s_*\rho=\rho''\).
Both \(p''s\) and the identity of \(R\) induce \(\rho\),
so its uniqueness clause gives \(p''s=1\).  This is exactly
the section required in APD5, for every permitted quotient.
Nothing requires \(\ker q=\ker u\); the original implication
\(u(y)=0\Rightarrow q(y)=0\) is sufficient.

The argument from a supplied set lift and a supplied boundary
does not use surjectivity of \(p\).  Keeping that general map,
and only the stated surjectivity of \(q\), proves the original
\texttt{pr\_\allowbreak obs} as well.  In particular the proof establishes
the original obstruction predicate, not just liftability across
one chosen one-dimensional extension.
\end{proof}
% END APD18

% BEGIN APD19
\paragraph{APD19. Every strict tangent cocycle comes from the universal ring.}
\label{claim:APD19}
At the hypotheses of APD18, for every strict one-cocycle
\((z,(a_v))\in Z_C^1\) there are
\(D\in\operatorname{Der}_{\mathcal O}(R,\kappa)\) and
\(X\in M_n(\kappa)\) such that
\[
 \phi(D)(s):=D(b(s))\bar b(s)^{-1}
              =z(s)+(\overline{\operatorname{Ad}}(s)-1)X
                         \quad(s\in\Gamma).             \tag{APD.39}
\]
Before the universality step, the dual-number construction below
has the broader original \texttt{pr\_\allowbreak tbox} scope: it applies to
any box over \(\kappa\) with locally constant matrix part,
under the original frame and coefficient premises.  It needs no
\(\mathrm{Univ}\), \(\mathrm{Good}\), or
\(\mathrm{Box.Of}\) assumption.  Those additional hypotheses
enter only when the constructed box is normalized and represented
by an algebra map from \(R\).

The derivative acts entrywise on the matrix part.  This is full
original \texttt{pr\_\allowbreak tsurj}, not a surjectivity assumption on
the tangent map.

\begin{proof}
The residual box is locally constant.  Indeed its entries modulo
\(\mathfrak m_R\) are locally constant by continuity, and
there are finitely many matrix entries; the sign is locally
constant too.  Over \(\kappa[\epsilon]/\epsilon^2\) define
\[
 b_z(s)=(1+\epsilon z(s))\bar b(s),
 \qquad g_{z,v}=\bar g_v(1+\epsilon a_v).                  \tag{APD.40}
\]
The one-cocycle identity makes \(b_z\) multiplicative.  Its
local conjugate is
\((1+\epsilon\bar g_v[z(\operatorname{loc}_v s)
 -(\overline{\operatorname{Ad}}(\operatorname{loc}_v s)-1)a_v]
 \bar g_v^{-1})\bar t_{v,s}\), which is upper triangular
by (APD.11).  Thus (APD.40) is a continuous box reducing to
\(\bar{\mathfrak B}\).  It is precisely the original
\texttt{pr\_\allowbreak tbox} construction.

Apply APD17 to the reduction onto \(\kappa\).  Its kernel is
\((\epsilon)\), annihilated by the maximal ideal.  We obtain
a conjugating matrix \(h=1+\epsilon Y\) and a normalized
\(\rho_z\) with residue \(\bar\rho\).  Universality gives
an actual algebra map \(R\to\kappa[\epsilon]/\epsilon^2\)
inducing \(\rho_z\); APD1 expresses it as
\(\operatorname{res}_R+\epsilon D\).

Both the image of the original box and the conjugated box have
matrix \(1\) at \(c\) and agree on \(G_0\).  The decomposition
\(\Gamma=G_0\sqcup G_0c\) therefore makes their signed
representations equal on all of \(\Gamma\).  The first-order
coefficient of conjugation by \(1+\epsilon Y\) is
\(z(s)+Y-\overline{\operatorname{Ad}}(s)Y\).
Multiplying the matrix derivative by \(\bar b(s)^{-1}\)
therefore gives (APD.39) with \(X=-Y\).  This proves the
claimed tangent statement at arbitrary local \(\mathcal O\)
and at every rank satisfying the original good-residue premise.
\end{proof}
% END APD19

% BEGIN APD20
\paragraph{APD20. The weight torus supplies the extra tangent direction.}
\label{claim:APD20}
Let
\(T_0=\operatorname{diag}(w(n)_i)\in M_n(\kappa)\).
At the hypotheses of APD18, if
\(\theta_{\mathrm{Lie}}(T_0)=T_0\), there is a tangent
vector \(D_0\) with
\[
       \phi(D_0)(s)=(\overline{\operatorname{Ad}}(s)-1)T_0
                         \quad(s\in\Gamma).             \tag{APD.41}
\]
This is the full original \texttt{pr\_\allowbreak tor}.

\begin{proof}
Extend the residual box to the dual numbers and conjugate it by
\(h=1-\epsilon T_0\), changing every local basis to
\(\bar g_vh^{-1}\).  Its first-order cochain is exactly the
right side of (APD.41).  The condition
\(\theta_{\mathrm{Lie}}(T_0)=T_0\) ensures
\(\theta(h)=h\), so its matrix at \(c\) remains \(1\).
The marked difference is diagonal and \(h\) is diagonal;
thus this difference remains diagonal too.  APD13 consequently
recovers a normalized continuous representation from this box
without further normalization.  Its residual representation is
\(\bar\rho\).  Universality and APD1 give a tangent vector
\(D_0\), and equality of the two boxes on \(G_0\) and \(c\)
gives (APD.41) on \(\Gamma\), exactly as in APD19.

For the actual ranks, \(T_0=\operatorname{diag}(1,-1)\) or
\(\operatorname{diag}(1,0,-1)\).  It is nonzero over every
field, and reversing its diagonal entries and negating returns
itself.  These are the two properties needed by (APD.13).
Notice that \(D_0\) itself is allowed to be zero when
\(T_0\) is invariant; APD6 retains that branch and obtains
the extra dimension from \(H_C^0\) instead.
\end{proof}
% END APD20

\subsection{The residual calculation and the exact numerical difference}

% BEGIN APD21
\paragraph{APD21. The actual residue is good and has the required ordinary bases.}
\label{claim:APD21}
Return to the actual \(O',R,\rho\) of (APD.1), and write
\(k'=\operatorname{Res}(O')\).  There is the specified residue
map \(f:k\hookrightarrow k'\) such that
\(\bar\rho.\mathrm{IsBar}(f)\); the matrices of \(\bar\rho\)
are all upper triangular or all lower triangular.  Its residue is
good in the exact sense of APD17.  There is a continuous box
\(\mathfrak B\) of \(\rho\).  These statements supply full
\texttt{pr\_\allowbreak res}, the actual box input, and full original
\texttt{pr\_\allowbreak good} at every field extension of \(k\).

\begin{proof}
Take the residue map induced by the original local map \(O\to O'\).
ALC17 gives the residual diagonal and character formulas after
passing further to \(\operatorname{Res}(R)\).  The canonical
map \(k'\to\operatorname{Res}(R)\) is an isomorphism, so
injectivity descends the same formulas to \(k'\).  APD8 gives
the stated global triangular alternative at this exact residue.

Here is the good-residue calculation over any target field \(K\)
and any \(f:k\to K\).  The odd prime condition gives
\(2\ne0\).  If \(z=\Psi(g_0)\), then \(z^2\ne1\),
and the diagonal entries of the marked difference are
\(q,-q\), or \(q,0,-q\), where \(q=f(z-z^{-1})\ne0\).
They are pairwise distinct, since \(2q\ne0\).
For each local place the frame supplies \(s_v\) with
\(z_v=\Psi(\operatorname{loc}_v s_v)\) satisfying
\(z_v^2\ne1\).  The ordered character values at \(s_v\)
are \(f(z_v),f(z_v^{-1})\), or
\(f(z_v),1,f(z_v^{-1})\).  They are pairwise distinct:
equality of the outer pair gives \(z_v^2=1\), and equality
with \(1\) gives \(z_v=1\).  This proves all three
\(\mathrm{Good}\) fields.  The last two conclusions of
\texttt{pr\_\allowbreak good}, namely \(T_0\ne0\) and
\(\theta_{\mathrm{Lie}}T_0=T_0\), were explicitly checked
in APD20.

In \(R\), the differences of the full character values at
\(s_v\) are units, by the residual calculation.  The exact
characteristic-polynomial identity over \(\mathcal A\) is
ACP4 and survives \(\iota\).  Together with the ordered
ordinary-product identity it gives the ordered ordinary basis
over \(R\) by ACP3.  Continuity of \(\rho\) is ALC17.
APD13 constructs the required box from these bases.  Thus all
inputs here are supplied for the actual local representation;
no ordinary flag is left as a new hypothesis on it.
\end{proof}
% END APD21

% BEGIN APD22
\paragraph{APD22. The local and infinite adjoint spaces have their exact dimensions.}
\label{claim:APD22}
For every field \(K\), every rank, and every
\(g\in\mathrm{GL}_n(K)\), the subspace
\(\mathfrak b(g)=\{X:gXg^{-1}\text{ is upper triangular}\}\)
has dimension \(n(n+1)/2\).  For a box of a good normalized
representation, its local subspaces further satisfy
\[
 \dim_K\mathfrak b_v=\frac{n(n+1)}2,
 \qquad
 (\operatorname{Ad}(\operatorname{loc}_v s)-1)Y\in\mathfrak b_v
 \text{ for all }s\quad\Longrightarrow\quad Y\in\mathfrak b_v.
                                                               \tag{APD.42}
\]
The rigidity assertion requires only an arbitrary frame, a box
of the supplied normalized representation, and its
\(\mathrm{Good}\) condition; it does not require
\(F.\mathrm{Ok}\).  They are stable under the local adjoint
action.  This stability
requires only a box over any commutative ring and its defining
identities; neither good residue nor \(F.\mathrm{Ok}\) is
needed for that assertion.  If
\(x\in\mathrm{GG}_n(K)\), \(x^2=1\), its sign is negative,
and \(2\ne0\) in \(K\), then
\[
            \dim_K\ker(\operatorname{Ad}(x)-1)
                              =\frac{n(n-1)}2.            \tag{APD.43}
\]
For any continuous box over an \(\mathrm{IsCO}\)-algebra its
residual signed representation and its residual adjoint
representation have open kernels.  These assertions prove full
\texttt{pr\_\allowbreak bdim}, \texttt{pr\_\allowbreak brig}, \texttt{pr\_\allowbreak bstab},
\texttt{pr\_\allowbreak cfix}, and \texttt{pr\_\allowbreak open} at their respective
original scopes.

\begin{proof}
Conjugation by \(g_v\) identifies \(\mathfrak b_v\) with
the upper triangular matrices.  The entries on and above the
diagonal are free coordinates, giving \(n(n+1)/2\) dimensions
over any field.  Products and inverses of invertible upper
triangular matrices are upper triangular, so conjugation by a
local box matrix preserves this subspace, over any commutative
ring.

For the second assertion of (APD.42), take a separating local
element supplied by \(\mathrm{Good}\), and in the ordinary
basis write its matrix as \(A\) and the matrix of \(Y\) as
\(Z\).  Then \(A\) is upper triangular with distinct
diagonal entries.  The assumed condition, multiplied on the
right by \(A\), says that \(AZ-ZA\) is upper triangular.
For \(i>j\) its \((i,j)\) entry equals
\[
 (A_{ii}-A_{jj})Z_{ij}
       +\sum_{q>i}A_{iq}Z_{qj}-\sum_{q<j}Z_{iq}A_{qj}.
\]
Starting at the greatest distance below the diagonal and
descending, the summands in the two sums already vanish.
The distinct diagonal entries force \(Z_{ij}=0\).  This
proves the rigidity assertion for every rank, not only ranks
two and three.

For (APD.43), write the matrix part of \(x\) as \(A\).
The signed involution identity says \(A\theta(A)=1\),
equivalently \(S=AJ_n\) is an invertible symmetric matrix.
The adjoint action is \(X\mapsto-SX^tS^{-1}\).
Thus its fixed matrices are exactly those with
\((XS)^t=-XS\).  Right multiplication by \(S\) identifies
this space with all skew-symmetric matrices.  Since \(2\ne0\),
their diagonal entries are zero and each pair \(i<j\) supplies
one free coordinate.  There are \(n(n-1)/2\) such pairs.
This proves the original integer identity
\(2\dim\ker(\operatorname{Ad}(x)-1)+n=n^2\), including
the ranks zero and one.

Finally each residual matrix entry of a continuous box is locally
constant.  Their finite tuple is locally constant.  The kernel
of the signed representation is the intersection of its
identity-matrix fiber with the open kernel of \(\delta\),
so is open.  The adjoint kernel contains this open subgroup
and is a union of its cosets, so is open as well.  This also
proves the original local-constancy output, even when the
coefficient residue field is infinite.
\end{proof}
% END APD22

% BEGIN APD23
\paragraph{APD23. The cyclotomic dual of the residual adjoint module is zero.}
\label{claim:APD23}
Let \(n=2\) or \(3\), let \(K\) be any field of
characteristic \(\ell\), let \(f:k\to K\), and let
\(\rho\) satisfy \(\mathrm{IsBar}(f)\), with all its
matrices upper triangular or all lower triangular.  Let
\(\mathfrak B\) be a box of \(\rho\).  Every linear
functional \(\lambda:M_n(K)\to K\) satisfying
\[
 \lambda(\operatorname{Ad}(s)X)
                   =\omega(s)\lambda(X)
                         \quad(s\in\Gamma,X\in M_n(K)) \tag{APD.44}
\]
is zero.  Here \(\omega=F.\mathrm{om}(K)\) is the exact
cyclotomic character.  This is full original \texttt{pr\_\allowbreak dual}.

\begin{proof}
Write \(a_{ij}=\lambda(E_{ij})\).  Then
\(\lambda(X)=\sum_{i,j}X_{ij}a_{ij}\).  For \(g\in G_0\),
the adjoint formula and (APD.44), applied after multiplying by
\(\rho(g)\), give
\[
 \sum_p\rho(g)_{pi}a_{pj}
                  =\omega(g)\sum_q a_{iq}\rho(g)_{jq}.
                                                               \tag{APD.45}
\]
Put \(M(g)=\rho(g)^t\).  If these matrices are upper
triangular, induct on \(i\) from the last index downwards
and on \(j\) from the first index upwards.  Once the entries
already reached by that induction have vanished, (APD.45)
reduces to
\[
       \rho(g)_{ii}a_{ij}
                         =\omega(g)a_{ij}\rho(g)_{jj}
                                  \quad(g\in G_0).       \tag{APD.46}
\]
Indeed terms with \(p<i\) or \(q>j\) vanish by triangularity,
and the other off-diagonal terms are exactly the earlier ones.
If \(M\) is lower triangular, reverse both orders.  It remains
to show that (APD.46) forces \(a_{ij}=0\).

If \(a_{ij}\ne0\), substitute the prescribed diagonal
characters in (APD.46) and cancel this nonzero scalar.  The
cyclotomic character commutes with the field map from \(k\),
since it is induced from the prime field.  Injectivity of \(f\)
therefore gives
\[
          \omega_k(g)\Psi(g)^{w_j-w_i}=1
                                  \quad(g\in G_0).       \tag{APD.47}
\]
The possible differences are \(0,\pm1,\pm2\).  For zero,
the frame's \(f2a\) contradicts (APD.47); for \(\pm1\),
its \(f2b\) does.  If the difference is \(\pm2\), then
\(j=i^{\mathrm{rev}}\).  The matrix part of the box at
\(c\) is \(1\), and
\(\theta_{\mathrm{Lie}}(E_{i,i^{\mathrm{rev}}})
             =-E_{i,i^{\mathrm{rev}}}\).
Equation (APD.44) hence gives \(\omega(c)=-1\), because
the coefficient is nonzero.  With \(\Psi(c)=1\) and
\(\delta(c)=-1\), it follows that the character
\(\omega_k\Psi^{w_j-w_i}\delta_k\) is \(1\) both on
\(G_0\), by (APD.47), and on \(c\).  These generate
\(\Gamma\); it is identically \(1\), contradicting the
frame's \(f2c\).  This settles every weight difference.

The triangular induction now makes every \(a_{ij}\) zero,
so \(\lambda=0\).  In particular the two outer weights in
rank three have not been treated as adjacent weights: their
contradiction uses the negative-sign element and the precise
\(\delta_k\) factor.  No exclusion of \(\ell=3\) occurs.
\end{proof}
% END APD23

% BEGIN APD24
\paragraph{APD24. The tangent excess is at least \(nF.d+1\).}
\label{claim:APD24}
For the actual local ring and box supplied above, let
\[
 T=\operatorname{Der}_{O'}(R,k'),\quad
 h_i=\dim_{k'}H_C^i\ (i=0,1,2),\quad t=\dim_{k'}T.
\]
The obstruction space \(H_C^2\) is finite dimensional and
\[
                         h_2+nF.d+1\le t.                \tag{APD.48}
\]

\begin{proof}
All input scopes for the frame's Euler field can now be checked.
The field \(k'\) is finite of characteristic \(\ell\), by
the original coefficient-ring fields of \(L\).  The adjoint
space \(W=M_n(k')\) is finite dimensional.  Its kernel is
open and the local Borels are stable by APD22.  The residual
representation is globally triangular by APD21, and APD23
proves that its cyclotomic dual is zero.  Thus the literal
\(F.\mathrm{EC}\) field gives finite dimensionality of
\(H_C^2\) and
\[
 h_2+h_0+\sum_v F.dv(v)\dim\mathfrak b_v
        \le h_1+\sum_{w\in F.\mathrm{Inf}}
             \dim\ker(\overline{\operatorname{Ad}}(F.cw(w))-1).
                                                               \tag{APD.49}
\]
This is the previously reconstructed Euler field at exactly its
original finite-module premise, not a conjectural dimension formula.

Put \(a=n(n-1)/2\).  APD22 gives
\(\dim\mathfrak b_v=a+n\).  Each \(F.cw(w)\) is an
involution of negative sign, so its fixed adjoint dimension is
\(a\), again by APD22.  The original frame identities are
\(\sum_v F.dv(v)=F.d\) and
\(\#F.\mathrm{Inf}=F.d\).  Substitution in (APD.49)
and cancellation of \(F.d\,a\) yield
\[
                         h_2+h_0+nF.d\le h_1.             \tag{APD.50}
\]

APD12 supplies actual universality, APD21 supplies the good
residue, and APD1 makes \(T\) finite dimensional.  APD19
is the second premise of APD6; the local rigidity in APD22
is its first.  The nonzero \(T_0\) and the actual tangent
vector of APD20 supply its extra-direction premise.  Hence
\(h_1+1\le t+h_0\).  Adding one to (APD.50), combining
these inequalities, and cancelling \(h_0\) proves (APD.48).
This also displays exactly where the additional \(1\), beyond
\(nF.d\), enters the relation count.
\end{proof}
% END APD24

% BEGIN APD25
\paragraph{APD25. Full original \texttt{pres} on every supplied local object.}
\label{claim:APD25}
For every original \(O,F,C_0,n,S,x,L\) of (APD.1), there
are integers \(N,N'\), a surjective \(O'\)-algebra map
\[
 p:O'[[X_1,\ldots,X_N]]\longrightarrow R,
 \quad f_1,\ldots,f_{N'}\in O'[[X_1,\ldots,X_N]],
\]
such that
\[
            N=N'+nF.d+1,\qquad \ker p=(f_1,\ldots,f_{N'}).
                                                               \tag{APD.51}
\]
This is the full original \(L.\mathrm{Pres}\), with no
\(F.\mathrm{Big}\) or nontriangularity assumption.  In particular
it is available at every original rank-two and rank-three local
chart, including those arising from the permitted Host prime
\(\ell=3\).

\begin{proof}
APD18 supplies the original
\(\mathrm{HasObs}(O',R,H_C^2)\).  Set
\(q=nF.d+1\), \(t=\dim T\), and
\(h=\dim H_C^2\).  APD24 says \(h+q\le t\).
Let
\[
              H'=H_C^2\oplus(k')^{t-q-h}.
\]
For every small extension, compose the actual map (APD.38)
with the inclusion in the first summand of \(H'\).  This
inclusion is injective, so its zero locus is unchanged.
The same sections supplied in APD18 prove
\(\mathrm{HasObs}(O',R,H')\), with every quotient in its
original quantifier still allowed.

The ring \(O'\) is complete Noetherian local with finite
residue field.  APD5 therefore gives a surjection with exactly
\(t\) variables and \(\dim H'=t-q\) listed relations.
Take \(N=t\) and \(N'=t-q\).  Inequality (APD.48)
makes these natural-number differences legitimate, and
\(N=N'+q\) proves (APD.51).  If fewer nonzero generators
are necessary, the construction pads the family with zeros;
this preserves its literal generating-ideal assertion.  The
kernel is proper, and all these relations lie in the maximal
ideal of the power-series ring.

No replacement of the coefficient extension \(O'\), the
closed point \(x\), or the local map \(\iota\) has occurred.
The full ordinary characters entered point realization, the
box, normalization, and the obstruction corrections.  Thus
the result is a presentation of the actual original \(R\),
not of a deformation problem obtained by discarding characters
or by identifying only residual semisimplifications.
\end{proof}
% END APD25

\paragraph{The exact next consumers and the retained arithmetic frontier.}
The rank-two consumer \texttt{sharp2} in
\texttt{Chained\allowbreak Bot.\allowbreak lean} first extracts precisely (APD.51).
The power-series ring has dimension \(N+1\), so the next
component-chain input is used with \(h=N'\) and
\(D=2F.d+1\).  Its numerical threshold is
\(D+1+h=N+1\).  Exact numerical control of both counts
supplies this threshold.  The remaining original
\texttt{c1ch}, with its \texttt{Cl\allowbreak Comp}/\texttt{Cl\allowbreak Chain}
bodies, and the triangular and coincident-character dimension
suppliers remain separate work.  The present chapter does not
assert \texttt{sharp2}, \texttt{sharp3}, or a new component
dimension bound.

Rank three also uses the ordinary presentation in its appropriate
branch, but its even/odd \(\mathrm{EPres}\) record has additional
data and is not supplied by (APD.51).  Global closed-point and
seed incidence, the rank-two-to-rank-three corner, full
\(\mathrm{AllGood}\), actual \(\mathrm{OffFin}\), and the
automorphic propagation inputs remain visible.  The two spaces
and modularity ideals remain distinct on their common arithmetic
frame.  The arithmetic entry still retains the original \(P\),
and the Alt route still retains its own \(\mathrm{WitA}\) and
\(\mathrm{WitAQ}\) witnesses.  No whole selected terminal
obligation is closed by this local presentation chapter.


\section{The reference unipotent covolume and its exact discriminant}
\label{sec:reference-covolume-v56}

This chapter proves the original \texttt{Covol\allowbreak Ref}.  The calculation concerns
the prescribed reference measure obtained from the Euclidean volume of the
standard nilpotent Lie algebra.  It is a separate calculation from the
adelic volume normalization and the actual \texttt{In\allowbreak Volumes} established
in GMR and SVW.  In particular, none of the supplied measures
\(X.\mu_N\), \(X.\mu_H\), the global character, or the adelic bridge is
changed here.

Fix any original \(D7\) datum \(D\) and any supplied
\(\mathrm{HidaSetting}(D)\), denoted by \(X\).  We work on its original
Hermitian space \((V,h)\), with \(h\) conjugate-linear in its first
variable.  Write
\[
 \begin{gathered}
 E=D.host.E,\quad F=D.host.F,\quad r=D.R.rS,\quad k=D.R.k,\\
 q=k+1,\quad d=[F:\mathbb Q],\quad
 \Delta=|\operatorname{disc}(X.B.I.F)|.
 \end{gathered}
\]
The supplied coefficient isomorphism identifies \(X.B.I.F\) with \(F\),
so this is also \(|\operatorname{disc}F|\).  We retain the supplied
ordering of the \(d\) real places and its chosen complex embeddings
\(\sigma_v:E\hookrightarrow\mathbb C\).  Put
\[
 \mathfrak d=\mathfrak D_{F/\mathbb Q}\mathcal O_E,
 \qquad \mathfrak a_b=\mathfrak d^{-(r-b)}\quad(0\leq b<r),
 \qquad T_w=\mathrm{Tw}(w)>0.
\tag{RCV.1}
\]
These are the literal \texttt{dd\allowbreak S}, \texttt{aa\allowbreak D}, and \texttt{Tw} of the
source.  Its hyperbolic vectors will be written \(e_b,f_b\), its
orthogonal-complement lattice \(L_b\), and its full lattice \(L_0\).
Thus the supplied lattice identity is
\[
 L_0=\bigoplus_{b<r}\mathfrak a_b e_b\ \oplus L_b\ \oplus
       \bigoplus_{b<r}\mathfrak a_b^{-1}f_b,
 \qquad h(e_b,f_a)=\delta_{ba}.
\tag{RCV.2}
\]
The directness in this display will also be verified below.  The letter
\(b\) in \(L_b\) is part of the source's name for one lattice; it does
not index a family of lattices.

For each word \(w\), the supplied vectors
\((w_{w,v,1},\ldots,w_{w,v,k},z_w)\) form a basis of
\(W_v=\langle e,f\rangle^\perp_v\), with Gram matrix
\(\operatorname{diag}(1,\ldots,1,-T_w)\).  We shall use two words
independently: \(u\) specifies a reference measure, and \(w\) specifies
the coordinate box and the arithmetic fundamental domain.  No equality
between them is assumed.

All source statements are read at commit
\texttt{7cbf7a0c32fca1badf05931879a37c3f8a2b03b4}.  The accompanying
manifest gives exact declaration ranges.  The only inherited human
results used below are ASG2, for the original trace-one element and
unramified CM extension, and AMD8, for Haar uniqueness at the source's
weak \texttt{Is\allowbreak Haar\allowbreak On\allowbreak Set} scope.  The intervening arithmetic and
coordinate arguments are proved in this chapter.

\subsection{The four arithmetic coordinate lattices}

% BEGIN RCV1
\paragraph{RCV1. Discriminants, extended ideals, and conjugation.}
\label{claim:RCV1}
For the actual coefficient fields and every \(b<r\),
\[
 |\operatorname{disc}E|=\Delta^2,\qquad
 N_E(\mathfrak d)=\Delta^2,\qquad
 N_E(\mathfrak a_b)=\Delta^{-2(r-b)}.
\tag{RCV.3}
\]
All these ideals are nonzero and conjugation invariant.  There is an
integral \(\theta\in E\) with \(\theta+\bar\theta=1\).
Here \(N_E\) denotes the positive absolute norm of a fractional ideal.
\begin{proof}
ASG2 applies to the original host and gives an integral trace-one
element.  The source's \(\mathrm{theta}(D)\) chooses an element with
exactly these two properties when one exists; we take that actual
choice as \(\theta\).  ASG2 also gives conjugation invariance and
absence of finite ramification in
\(E/F\).  We recall the determinant consequence to specify the exact
normalization being used.  For a finite extension of number fields,
the discriminant of the integral trace pairing satisfies
\[
 |\operatorname{disc}E|
 =|\operatorname{disc}F|^{[E:F]}
       N_F(\mathfrak D_{E/F}^{\mathrm{disc}}).
\]
This follows by writing the trace as the composite of the two traces:
in local integral bases the determinant is the determinant from the
base field repeated \([E:F]\) times, multiplied by the norm of the
relative trace determinant.  Changing a local basis multiplies the
determinant by the square of a unit, so the local computations glue
as fractional ideals.  At a split or unramified quadratic place the
integral relative trace pairing is perfect.  Its determinant is a
unit, including at residue characteristic two.  Consequently the
relative discriminant ideal is the unit ideal, and \([E:F]=2\) gives
the first equality in (RCV.3).

For a nonzero ideal \(J\) of \(\mathcal O_F\),
\(N_E(J\mathcal O_E)=N_F(J)^2\).  To see this without a global
free-basis assumption, localize at each prime of \(\mathcal O_F\).
The integral algebra \(\mathcal O_E\) is a free module of rank two
there, and its quotient by a power of the prime has twice the
module length of the corresponding quotient of the base ring.
Multiplying residue-field cardinalities proves the norm identity.
It extends to fractional ideals by taking quotients.
The trace-dual definition of the different gives
\(N_F(\mathfrak D_{F/\mathbb Q})=|\operatorname{disc}F|\): in a
\(\mathbb Z\)-basis, the index of the trace dual is the absolute
determinant of the trace Gram matrix.  Apply the extended-ideal norm
identity and then take the prescribed negative powers.  This proves
the remaining equalities.  Finally, the supplied field isomorphism
preserves integral elements and the rational trace pairing; it
therefore preserves its discriminant.  Thus every occurrence of
\(\Delta\) in (RCV.3) is the discriminant of the field in the actual
consumer, not of a replacement field.
\end{proof}
% END RCV1

% BEGIN RCV2
\paragraph{RCV2. Euclidean covolumes with every measurable domain.}
\label{claim:RCV2}
Let \(\Lambda\) be the integer span of a real basis
\(v_1,\ldots,v_m\) of a finite-dimensional real space with its specified
Lebesgue measure.  It is countable and discrete, has a Borel strict
additive fundamental domain, and every measurable strict additive
fundamental domain has volume \(|\det(v_1,\ldots,v_m)|\).
In the chosen embeddings of a CM field this gives, for every nonzero
fractional ideal \(I\subset E\),
\[
 \operatorname{covol}(\sigma(I))
 =2^{-d}\sqrt{|\operatorname{disc}E|}\,N_E(I).
\tag{RCV.4}
\]
For a fractional ideal \(J\subset F\), its real embedding lattice has
covolume \(\sqrt{\Delta}\,N_F(J)\).
\begin{proof}
The image of \([0,1)^m\) under the basis matrix is a Borel fundamental
parallelotope: coordinatewise integer subtraction gives existence
and the half-open intervals give uniqueness.  The same invertible
matrix proves discreteness and the determinant formula.  If \(A\)
is any other measurable strict fundamental domain and \(P\) is this
parallelotope, partition \(A\) by the countable sets
\(A\cap(\lambda+P)\), \(\lambda\in\Lambda\).  Translating these
pieces by \(-\lambda\) partitions \(P\).  Translation invariance
and countable additivity give \(\operatorname{vol}A=
\operatorname{vol}P\).  This argument first takes place in
\([0,\infty]\); it does not assume that \(A\) has finite volume.

For an integral ideal, choose an integer basis and form the matrix
of all field embeddings.  The square of its determinant is the
discriminant of that basis; its absolute value is
\(\sqrt{|\operatorname{disc}E|}\,N_E(I)\), by the integral index
formula.  Replacing the two rows \((\sigma,\bar\sigma)\) by
\((\operatorname{Re}\sigma,\operatorname{Im}\sigma)\) divides
the absolute determinant by two.  There are exactly \(d\) pairs,
which proves (RCV.4).  Clearing an integer denominator proves the
fractional-ideal version.  The same argument with only real rows
proves the assertion for \(F\).

Permuting the places or conjugating any of the chosen representatives
only permutes real coordinates or changes the sign of an imaginary
coordinate.  Its absolute determinant is one.  These formulas
therefore hold for the actual embedding and place witnesses in
\(X\).  The zero-dimensional case has the one-point parallelotope,
with volume and empty determinant equal to one.
\end{proof}
% END RCV2

% BEGIN RCV3
\paragraph{RCV3. The actual orthogonal-complement lattice is self-dual.}
\label{claim:RCV3}
Put \(W=\langle e_b,f_b:b<r\rangle^\perp\).  Then
\(\dim_E W=q\), \(L_b=L_0\cap W\), and \(L_b\) is a full
\(\mathcal O_E\)-lattice in \(W\) which is self-dual for \(h|_W\).
\begin{proof}
Pairing a relation among the \(e_b,f_b\) with each \(f_a,e_a\)
proves their linear independence and identifies their Gram matrix
as the hyperbolic matrix.  Their span is nondegenerate.  Thus
\(V=\langle e,f\rangle\perp W\) and
\(\dim_E W=(2r+k+1)-2r=q\).  The projection onto \(W\) is
\[
 p_W(y)=y-\sum_b h(f_b,y)e_b-\sum_b h(e_b,y)f_b.
\tag{RCV.5}
\]
Apply it to the supplied sum decomposition of \(L_0\).  It vanishes
on the two hyperbolic summands and is the identity on \(L_b\),
so its image is precisely \(L_b\).  Pairing with \(e,f\) also
proves directness in (RCV.2), and proves \(L_0\cap W=L_b\).
Finite generation of \(L_0\) and its full span now imply the
same two properties for \(L_b\) in \(W\).

The original \(V\) record supplies Hermitian unimodularity of
\(L_0\).  Its hyperbolic summands are themselves perfect duals:
\(\overline{\mathfrak a_b}=\mathfrak a_b\), and the dual of
\(\mathfrak a_b e_b\) in the \(f_b\) direction is
\(\mathfrak a_b^{-1}f_b\).  More directly, if \(x\in W\)
pairs integrally with \(L_b\), it pairs as zero with all the
hyperbolic summands.  It therefore pairs integrally with \(L_0\),
belongs to \(L_0\) by unimodularity, and belongs to \(L_b\) by
the intersection identity.  The converse follows from integrality
of \(h\) on \(L_0\).  This proves the asserted self-duality,
without choosing a principal \(\mathcal O_E\)-basis for \(L_b\).
\end{proof}
% END RCV3

% BEGIN RCV4
\paragraph{RCV4. The complete trace Gram determinant, in any integer basis.}
\label{claim:RCV4}
For every nonzero fractional ideal \(a\) of \(E\) and every
\(\mathbb Z\)-basis \(m_1,\ldots,m_{2dq}\) of \(aL_b\),
\[
 \left|\det\big(\operatorname{Tr}_{E/\mathbb Q}
                         h(m_i,m_j)\big)\right|
       =|\operatorname{disc}E|^q N_E(a)^{2q}.
\tag{RCV.6}
\]
In particular this proves the full original \texttt{Tr\allowbreak Det\allowbreak Unimod}
when \(a=\mathfrak a_b\), for every supplied basis and its original
dimension witness.
\begin{proof}
First let \(L=L_b\), and let \(L^{\mathrm{tr}}\) be its dual for
the rational bilinear form \((x,y)\mapsto\operatorname{Tr}h(x,y)\).
If the trace pairing with every \(y\in L\) is integral, the same
is true with every \(cy\), \(c\in\mathcal O_E\).  By the
definition of the codifferent this says
\(h(x,L)\subset\mathfrak D_{E/\mathbb Q}^{-1}\).  Localize
at a prime of \(\mathcal O_E\).  There the different is principal
and \(L\) is free; RCV3 identifies the Hermitian dual of \(L\)
with \(L\).  Multiplying by a local generator of the different
therefore proves
\[
           L^{\mathrm{tr}}=\mathfrak D_{E/\mathbb Q}^{-1}L.
\tag{RCV.7}
\]
Equality at all primes proves equality of the original fractional
lattices.  Conjugation does not change this formula: the different
is stable under every field automorphism, and we are using an
absolute determinant.

A rank-\(q\) projective \(\mathcal O_E\)-lattice satisfies
\([L:IL]=N_E(I)^q\) for every nonzero integral ideal \(I\).
Indeed, after localization at a prime it is free of rank \(q\),
so the quotient by the \(s\)-th power of that prime has length
\(qs\).  Multiplication over the primes proves the formula.
Apply this to the different in (RCV.7).  Since its norm is
\(|\operatorname{disc}E|\),
\([L^{\mathrm{tr}}:L]=|\operatorname{disc}E|^q\).
In an integer basis of \(L\), the trace-dual basis is obtained
by the inverse Gram matrix.  Thus this index is exactly the
absolute Gram determinant.  This proves (RCV.6) for \(a=\mathcal O_E\).

For integral \(a\), a basis matrix for \(aL\) relative to a
basis of \(L\) has absolute determinant \([L:aL]=N_E(a)^q\).
A bilinear Gram matrix changes by \(B^tGB\), so its determinant
changes by \(N_E(a)^{2q}\).  For arbitrary fractional \(a\),
multiply it by a positive integer to reduce to this case and
cancel the corresponding scalar determinant.  This also proves
that every integer basis has the same answer: a change of such
basis has determinant \(\pm1\).  The rank is \(2dq\) because
the lattice is full in the \(q\)-dimensional \(E\)-space \(W\).
\end{proof}
% END RCV4

% BEGIN RCV5
\paragraph{RCV5. The precise real quadratic form on the row coordinates.}
\label{claim:RCV5}
For a word \(w\), define the additive row map
\[
 \rho_w(x)_v=
 \big(h_v(x,w_{w,v,1}),\ldots,h_v(x,w_{w,v,k}),h_v(x,z_w)\big)
       \quad(x\in W).
\tag{RCV.8}
\]
It extends to a real linear isomorphism
\(W\otimes_{\mathbb Q}\mathbb R\simeq\prod_v\mathbb C^q\).
Its being conjugate-linear in \(x\) over \(E\) is harmless and
is not replaced by complex linearity.  For \(x,y\in W\),
\[
 \frac12\operatorname{Tr}_{E/\mathbb Q}h(x,y)
 =\sum_v\operatorname{Re}\left(
    \sum_{j=1}^k \rho_w(x)_{vj}\overline{\rho_w(y)_{vj}}
       -\frac{\rho_w(x)_{vq}\overline{\rho_w(y)_{vq}}}{T_w}\right).
\tag{RCV.9}
\]
\begin{proof}
The vectors in (RCV.8) have nonsingular Gram matrix
\(G_w=\operatorname{diag}(1_k,-T_w)\) and are \(q\) in number,
so they are a basis of each \(W_v\).  For example the coefficient
expansion is
\[
 y=\sum_j h_v(w_{w,v,j},y)w_{w,v,j}
                         -T_w^{-1}h_v(z_w,y)z_w.
\]
Inserting this into \(h_v(x,y)\) gives the expression inside the
real part in (RCV.9).  The trace of an element of a CM field is
twice the sum of its real parts over one member of each conjugate
pair.  This proves the displayed identity.

The usual real scalar extension of the field-embedding map is an
isomorphism \(E\otimes_{\mathbb Q}\mathbb R\simeq\mathbb C^d\):
its matrix in a rational basis has nonzero discriminant.  Apply this
to coordinates on \(W\), and then apply the nonsingular Hermitian
coordinate maps just described at each place.  This proves the real
linear isomorphism assertion.  In particular the images of any
integer basis of a full lattice in \(W\) are a real basis.
\end{proof}
% END RCV5

% BEGIN RCV6
\paragraph{RCV6. Full row-block covolume.}
\label{claim:RCV6}
For every \(X,w,b\), the source's row lattice is exactly
\(\rho_w(\mathfrak a_b L_b)\).  It is countable and discrete,
has a measurable strict additive fundamental domain, and every
such domain has volume
\[
 c_{R_b}(w)=T_w^d\,2^{-dq}
                \Delta^{q(1-2(r-b))}.
\tag{RCV.10}
\]
This is the full original \texttt{Row\allowbreak Block\allowbreak Covol}.
\begin{proof}
The map \(\rho_w\) is additive and commutes with integer scalars.
Consequently its image of \(\mathfrak a_bL_b\) is already an
additive subgroup.  The \texttt{Add\allowbreak Subgroup.\allowbreak closure} in the source
means the subgroup generated by that image; it does not mean a
topological closure.  It therefore gives this same lattice.

Choose any integer basis \((m_i)\) of \(\mathfrak a_bL_b\),
and let \(B\) be its real-coordinate matrix under \(\rho_w\).
In real and imaginary coordinates the form on the right of (RCV.9)
has diagonal matrix \(H\), with each entry of
\((1_k,-T_w^{-1})\) repeated twice at each place.  Hence
\(|\det H|=T_w^{-2d}\).  Equations (RCV.6) and (RCV.9) imply
\[
 |\det B|^2 T_w^{-2d}
 =2^{-2dq}|\operatorname{disc}E|^q
                          N_E(\mathfrak a_b)^{2q}.
\]
Taking the positive square root and using (RCV.3) gives precisely
(RCV.10).  The determinant is nonzero; RCV2 now gives countability,
discreteness, existence, and the formula for every measurable
strict domain.  No special choice of row basis, orthonormal frame,
or domain survives in the answer.
\end{proof}
% END RCV6

% BEGIN RCV7
\paragraph{RCV7. The two complex ideal blocks.}
\label{claim:RCV7}
For \(0\leq b<a<r\), the source's \(X\)- and \(K\)-block
lattices and their covolumes are respectively
\[
\begin{array}{c|c|c}
 \text{block}&\text{lattice in }\mathbb C^d&\text{covolume}\\ \hline
 X_{ba}&\sigma(\mathfrak a_b\mathfrak a_a^{-1})
                 &2^{-d}\Delta^{1-2(a-b)}\\
 K_{ba}&\sigma(\mathfrak a_b\mathfrak a_a)
                 &2^{-d}\Delta^{1-2(r-b)-2(r-a)}.
\end{array}
\tag{RCV.11}
\]
Each assertion includes countability, existence of a measurable
strict additive domain, and equality for every such domain.
\begin{proof}
These images are additive subgroups, so again the generated-subgroup
notation does not enlarge them.  Apply (RCV.4) and
\(\sqrt{|\operatorname{disc}E|}=\Delta\).  Multiplicativity of
the fractional ideal norm gives, respectively,
\[
 N_E(\mathfrak a_b\mathfrak a_a^{-1})=\Delta^{-2(a-b)},
 \qquad
 N_E(\mathfrak a_b\mathfrak a_a)
       =\Delta^{-2(r-b)-2(r-a)}.
\]
The ideals are nonzero.  RCV2 supplies all the domain and
countability assertions, with the source's actual embeddings.
\end{proof}
% END RCV7

% BEGIN RCV8
\paragraph{RCV8. The imaginary diagonal block, including the factor at two.}
\label{claim:RCV8}
For \(0\leq a<r\), put
\(I=\mathfrak a_a^2\) and
\(I^- =\{t\in I:\bar t=-t\}\).  The lattice
\(\operatorname{Im}\sigma(I^-)\subset\mathbb R^d\) has covolume
\[
                  c_{D_a}=\Delta^{1/2-2(r-a)}.
\tag{RCV.12}
\]
It is the source's diagonal block, with the same countability,
existence, and every-domain conclusions as in RCV7.
\begin{proof}
Let \(J=\mathfrak D_{F/\mathbb Q}^{-2(r-a)}\), so
\(I=J\mathcal O_E\).  Extension and contraction give
\(I\cap F=J\).  For completeness, this can be checked after
localizing at a prime of \(\mathcal O_F\): write the fractional
ideal as one generator times the local ring, divide by that
generator, and use \(\mathcal O_E\cap F=\mathcal O_F\).

The relative trace carries \(I\) into \(J\).  It is onto,
because \(\theta j\in I\) has trace \(j\) for every \(j\in J\).
We thus have the split exact sequence of integer lattices
\[
 0\longrightarrow I^-\longrightarrow I
       \xrightarrow{\operatorname{Tr}_{E/F}}J\longrightarrow0,
 \qquad j\longmapsto\theta j.
\tag{RCV.13}
\]
This is an integral splitting, including at two.  In the chosen
complex embeddings, every element of \(I^-\) has real part zero,
whereas \(\operatorname{Re}\sigma(\theta j)=\sigma(j)/2\).
Choose integer bases of the kernel and of \(J\), and lift the
latter by \(j\mapsto\theta j\).  The resulting real-coordinate
basis matrix for \(I\) is block triangular.  Its determinant gives
\[
 \operatorname{covol}_{\mathbb C^d}\sigma(I)
 =\operatorname{covol}_{\mathbb R^d}\big(\tfrac12\sigma(J)\big)
       \operatorname{covol}_{\mathbb R^d}\operatorname{Im}\sigma(I^-).
\tag{RCV.14}
\]
Both determinant factors are nonzero; the sequence also shows that
the kernel has integer rank \(d\).  RCV2 and RCV1 give the first
factor on the left as \(2^{-d}\Delta N_F(J)^2\), and the
first factor on the right as \(2^{-d}\sqrt\Delta N_F(J)\).
Their quotient is \(\sqrt\Delta N_F(J)\), which is (RCV.12).
Thus the factors \(2^{-d}\) cancel here.  Replacing this kernel
by the imaginary parts of all of \(I\) would lose exactly the
integral condition that justifies this cancellation.
Finally apply RCV2 to its integer basis to obtain every measurable
strict-domain conclusion.
\end{proof}
% END RCV8

% BEGIN RCV9
\paragraph{RCV9. The complete original block-covolume package.}
\label{claim:RCV9}
For every supplied \(X\), every word \(w\), and every block in
\[
 \mathcal P=(\{(b,a):b<a<r\}\sqcup\{b<r\})
                \sqcup(\{(b,a):b<a<r\}\sqcup\{a<r\}),
\]
the source's \(cLat(X,w,p)\) is a full lattice in its specified
real or complex block space.  Its Lebesgue covolume is exactly
\(cVal(X,w,p)\), as displayed in (RCV.10)--(RCV.12).
This proves \texttt{Block\allowbreak Covol} with all its original quantifiers.
\begin{proof}
The four alternatives for \(p\) are exactly the \(X\), row,
\(K\), and diagonal blocks proved in RCV6--RCV8.  Their measures
are the standard Lebesgue measures on \(\mathbb C^d\),
\((\mathbb C^q)^d\), \(\mathbb C^d\), and \(\mathbb R^d\),
respectively.  Their orderings differ from real-imaginary coordinate
order only by coordinate permutations, whose absolute determinant
is one.  In each case we have proved all three fields of the original
package: countability of the subgroup, existence of a measurable
strict additive domain, and its stated volume for every supplied
measurable strict additive domain.  Every displayed covolume is
finite and strictly positive.  There is no remaining row-determinant
or diagonal-lattice premise.
\end{proof}
% END RCV9

\subsection{The actual group and its arithmetic subgroup in coordinates}

% BEGIN RCV10
\paragraph{RCV10. A complete polynomial chart for the unipotent group.}
\label{claim:RCV10}
At one real place fix a word \(w\), put
\(G_w=\operatorname{diag}(1_k,-T_w)\), and use the ordered basis
\((e,W_w,f)\) of \(V_v\).  Every element of the original
\(N_v\) has a unique expression
\[
 n(Q,B,K)=
 \begin{pmatrix}
 Q&B&QM\\
 0&1&-G_w^{-1}B^*Q^{-*}\\
 0&0&Q^{-*}
 \end{pmatrix},\qquad
 M=K-\tfrac12Q^{-1}BG_w^{-1}B^*Q^{-*},
\tag{RCV.15}
\]
where \(Q\) is upper unitriangular, \(B\) is an arbitrary
\(r\)-by-\(q\) complex matrix, and \(K^*=-K\).  Here
\(Q^{-*}=(Q^{-1})^*\).
The free coordinates are
\[
 Q_{ba}\ (b<a),\qquad B_{bj}\ (j\leq q),\qquad
 K_{ba}\ (b<a),\qquad \operatorname{Im}K_{aa}.
\tag{RCV.16}
\]
They are exactly the grouped source coordinates \(\operatorname{blk}\);
their real and imaginary parts, with the specified permutation, are
\(\kappa_w\).  They give a homeomorphism from \(N_\infty\) to a
real coordinate space and a measurable section into the ambient group.
Thus the original \texttt{Kap\allowbreak Surj} and \texttt{Coord\allowbreak Meas} hold.
\begin{proof}
The source defines \(N_v\) by the isometry condition,
\((n-1)e_a\in\langle e_b:b<a\rangle\), and
\((n-1)e^\perp\subset\langle e\rangle\).  Its additional
condition on \(z_w\) follows from the last one because
\(z_w\perp e\).  In particular the subgroup is independent of
the word.  The first two conditions force the upper block triangular
shape, with upper unitriangular \(Q\) and identity middle block.
Preserving the pairings of the \(e\)- and \(f\)-columns forces
the lower right block to be \(Q^{-*}\).  Pairing the middle
columns with the \(f\)-columns forces the middle right block in
(RCV.15).  Pairing two \(f\)-columns then gives
\[
 M+M^*=-Q^{-1}BG_w^{-1}B^*Q^{-*}.
\tag{RCV.17}
\]
Its solutions are exactly the stated \(M\), with arbitrary
skew-Hermitian \(K\).  Conversely these computations verify every
pairing of the displayed columns, so (RCV.15) is an isometry with
the required flag conditions.  Its diagonal blocks are invertible,
which also verifies that it lies in the actual unitary group.

The coordinate readings are literal: \(Q_{ba}=h(f_b,ne_a)\),
\(B_{bj}=h(f_b,nw_j)\) and \(B_{bq}=h(f_b,nz_w)\), while
\(M=Q^{-1}(h(f_b,nf_a))_{ba}\).  Its skew part is \(K\);
the diagonal entries of the Hermitian correction in (RCV.17)
are real, so \(\operatorname{Im}M_{aa}=\operatorname{Im}K_{aa}\).
These are precisely the four source blocks.

The inverse of an upper unitriangular matrix is the finite polynomial
\(1-U+U^2-\cdots+(-U)^{r-1}\), where \(U=Q-1\).
Consequently (RCV.15) and its coordinate readings on \(N_v\)
are continuous polynomial maps.  The subgroup is closed, since its
flag and isometry equations are closed finite-dimensional equations.
Taking the product over the supplied places yields the stated
homeomorphism and section.  Images of measurable subsets of the
group under this chart are measurable by taking inverse images under
the continuous inverse chart.  On ambient matrices, the source uses
the total matrix-inverse operation in the reading of \(M\); it is
Borel measurable, being the adjugate times the reciprocal determinant
with the specified value at determinant zero.  This verifies the
ambient measurability required by \texttt{Coord\allowbreak Meas} as well.
\end{proof}
% END RCV10

% BEGIN RCV11
\paragraph{RCV11. The full triangular multiplication law and its measure.}
\label{claim:RCV11}
Give the four blocks the source heights
\[
 h(X_{ba})=a-b,\quad h(R_b)=2r-b,\quad
 h(K_{ba})=4r+1-b-a,\quad h(D_a)=4r+1-2a.
\tag{RCV.18}
\]
In block coordinates multiplication has the form
\[
       (x*y)_p=x_p+y_p+\mathcal B_p(x_{<h(p)},y_{<h(p)}),
\tag{RCV.19}
\]
where each correction is a polynomial, depends only on strictly
lower-height blocks in both arguments, and vanishes if either
argument is zero.  This is the full original \texttt{Kap\allowbreak Law}.
Coordinate Lebesgue measure is both left and right invariant.
\begin{proof}
Multiply (RCV.15) for \((Q,B,K)\) and \((Q',B',K')\).
Writing \(U=Q^{-1}\), \(U'=Q'^{-1}\), one obtains
\[
 \begin{split}
 Q''&=QQ',\qquad B''=QB'+B,\\
 M''&=M'+U'MU'^*-U'UBG_w^{-1}B'^*U'^*,\\
 K''&=K'+U'KU'^*-\operatorname{skew}
                         (U'UBG_w^{-1}B'^*U'^*),
 \end{split}
\tag{RCV.20}
\]
where \(\operatorname{skew}(A)=(A-A^*)/2\).  The last line
uses that \(U'HU'^*\) is Hermitian whenever \(H\) is Hermitian,
so the Hermitian part of \(M\) contributes no skew part.

In \(Q''_{ba}\), every correction has indices \(b<c<a\),
so its two \(X\)-heights are strictly less than \(a-b\).
In row \(b\) of \(B''\), every correction uses \(Q_{bc}\)
and row \(c>b\) of \(B'\), again of strictly lower heights.
For the central block, every entry of \(U'KU'^*\) contributing
to position \((b,a)\) has central indices \(c\geq b,e\geq a\).
Apart from \((c,e)=(b,a)\), which contributes exactly \(K_{ba}\),
its central height is \(4r+1-c-e<4r+1-b-a\).  This statement
also covers conjugate lower entries and diagonal entries, since their
heights use the sum of the indices.  Entries of \(U,U'\) are
polynomials in the lower-height \(X\)-blocks.  Every row block
appearing in the remaining term in (RCV.20) has height at most
\(2r\), whereas every central height is at least \(2r+3\).
This proves strict dependence for both off-diagonal and diagonal
central coordinates.  All the corrections are polynomials, hence
jointly measurable.  Substitution of the identity element proves
their asserted vanishing; strict dependence makes the same statement
true whenever just the relevant lower coordinates vanish.

Order the finitely many blocks by increasing height, breaking ties
arbitrarily.  Left translation by a fixed point is a triangular
map with identity diagonal blocks: successively it translates each
block by a function of the preceding blocks.  Each such map is a
global measurable bijection and preserves product Lebesgue measure,
by the translation invariance in that block and Tonelli's theorem.
Their finite composition has the same property.  The identical
argument with the other argument fixed proves right invariance.
Thus the coordinate measure is a genuine Haar measure on the actual
closed unipotent subgroup.
\end{proof}
% END RCV11

% BEGIN RCV12
\paragraph{RCV12. The needed finite-level membership, on the supplied lattice.}
\label{claim:RCV12}
For every supplied \(X,w\) and every finite adelic element
\(n\in N_{\mathrm{fin}}(w)\cap K\),
\[
                         x_w^{-1}nx_w\in K_2.
\tag{RCV.21}
\]
Here \(K\), \(K_2\), \(x_w=X.xu(w)\), and all finite frame
vectors are the original ones.  This is the full source theorem
\texttt{unipotent\_\allowbreak conj\_\allowbreak mem\_\allowbreak level}; it assumes no character
identity.
\begin{proof}
Let \(A\) be the integral finite adelic coefficient ring and
\(L=X.B.I.Lf\).  Its form is integral.  Put \(g=X.B.g\).
The supplied \texttt{basis\allowbreak Approx} and \texttt{xu\_\allowbreak z} assertions
give one positive integer \(C\), coprime to \(g\), such that
\[
 \begin{split}
 C(e_b-e_b^0),\ C(f_b-f_b^0)&\in gL,\\
 C(x_w^{-1}z_w-z^0)&\in gL
 \end{split}
\tag{RCV.22}
\]
for all \(b\).  Take the product of their finitely many individual
clearing integers.  Work for the moment in \(L_C=L[1/C]\) over
\(A_C=A[1/C]\), modulo \(gL_C\).  The actual \(e_b,f_b\)
belong to \(L_C\) and are congruent to \(e_b^0,f_b^0\).
Moreover, the natural map \(L/gL\to L_C/gL_C\) is an
isomorphism: multiplication by \(C\) is invertible modulo \(g\)
by an integer Bezout identity.  This observation will descend every
congruence below to the original module, with no localization left
in the conclusion.

Set \(a=x_w^{-1}nx_w\).  The supplied \(x_w\in K\) and
\(n\in K\) imply \(a\in K\); all three maps preserve \(L_C\).
The supplied \(x_w\) fixes each \(e_b,f_b\), hence preserves
their orthogonal complement.  Conjugation by it therefore preserves
the flag unipotent conditions.  In particular
\((a-1)e_b\in\langle e_c:c<b\rangle\).
An element \(s\in L_C\cap\langle e_c:c\in S\rangle\)
has the exact expansion
\(s=\sum_{c\in S}h(f_c,s)e_c\), whose coefficients lie in
\(A_C\).  Modulo \(g\), its image belongs to the span of
\(e_c^0\), \(c\in S\).  Apply this with
\(s=(a-1)e_b\) and then use \(e_b\equiv e_b^0\) to obtain
the first family of congruences defining \(K_2\).

Next, \((n-1)z_w\in\langle e\rangle\).  Since \(x_w\)
fixes that span, \((a-1)x_w^{-1}z_w\) belongs to it and to
\(L_C\).  Equation (RCV.22) gives
\((a-1)z^0\in\langle e^0\rangle+gL_C\), the second
congruence of \(K_2\).

Finally let \(\beta_i\) be any supplied \(bw_i\).  Its original
Gram congruences give \(h(e_b^0,\beta_i)\in gA\), and hence
\(h(e_b,\beta_i)\in gA_C\).  Decompose
\[
 \beta_i=\sum_b h(f_b,\beta_i)e_b+p_W\beta_i
                              +\sum_b h(e_b,\beta_i)f_b.
\]
All these terms belong to \(L_C\), because the frame vectors
and their pairings are integral there.  The last sum belongs to
\(gL_C\).  The flag conditions put \((a-1)\) of the first
two terms in \(\langle e\rangle\); the first term uses the
upper flag condition and the second uses \(p_W\beta_i\in e^\perp\).
Thus \((a-1)\beta_i\) is, modulo \(gL_C\), an integral
combination of \(e_b^0\).  This is stronger than the third
congruence of \(K_2\), which permits earlier \(bw_j\) as well.

All the vectors \((a-1)e_b^0,(a-1)z^0,(a-1)\beta_i\)
belong to \(L\).  The isomorphism modulo \(g\) descends the
three span congruences to \(L/gL\): coefficients in \(A_C/gA_C\)
come from \(A/gA\).  Together with \(a\in K\), they are
exactly \(a\in K_2\).  Notice that the argument uses only
membership, so neither \(\mathrm{PsiDd}\) nor any equality of
finite characters has been introduced as an extra premise.
\end{proof}
% END RCV12

% BEGIN RCV13
\paragraph{RCV13. The exact rational-arithmetic dictionary.}
\label{claim:RCV13}
For every \(X,w\), the subgroup \(\Gamma_w=\mathrm{NOinf}(X,w)\)
is the archimedean image of
\[
            \{n\in N(E):nL_0=L_0\}.
\tag{RCV.23}
\]
The rational isometry and its finite and infinite components are
linked by the original \texttt{Rat\allowbreak Link}.  This proves the complete
original \texttt{NOrat\allowbreak C}, not only an inclusion in an auxiliary
arithmetic group.
\begin{proof}
By definition, \(s\in\Gamma_w\) means there is an original
rational adelic point \(\gamma\in X.B.I.GFset\) with
\(\gamma_\infty=s\), whose two components lie in the specified
unipotent groups, and such that
\(x_w^{-1}\gamma_f x_w\in K_2\).  Since \(K_2\subset K\)
and \(x_w\in K\), we have \(\gamma_f\in K\).

For a rational vector \(y\), membership in \(L_0\) is equivalent
to membership of its finite adelic image in \(L\).  One elementary
verification chooses an integer basis of the full lattice \(L_0\).
The rational coordinates of \(y\) are integral in every \(\mathbb Z_p\)
if and only if they are integers.  Scalar extension to the finite
adeles is exactly the original \(L\), by the supplied lattice
and coefficient-coordinate identities.  Consequently preservation of
\(L\) by a rational isometry and its inverse is equivalent to
preservation of \(L_0\).

The original rational point therefore gives an isometry \(n\)
preserving \(L_0\).  Its infinite flag conditions descend to
rational ones.  Indeed, decompose a rational vector using the
hyperbolic projection (RCV.5); vanishing of any unwanted coefficient
at an embedding forces its vanishing in \(E\), and the same
argument applied to a rational basis of \(e^\perp\) proves
the universal perpendicular-space condition.  Thus \(n\in N(E)\).
This proves one direction of (RCV.23).

Conversely, a rational isometry in (RCV.23) gives its actual pair of
adelic components through the supplied coordinates.  The isometry,
flag, and perpendicular-space equations extend to all completions;
the latter can be checked on a rational basis of \(e^\perp\).
Hence its components belong to \(N_\infty\) and \(N_{\mathrm{fin}}\).
The lattice argument above gives \(\gamma_f\in K\), and RCV12
gives \(x_w^{-1}\gamma_f x_w\in K_2\).  These are precisely
the defining witnesses for \(\Gamma_w\).  The coefficient map
is injective, so no rational point is lost or duplicated.
In particular \(\Gamma_w\) is countable and is a subgroup of
the actual \(N_\infty\).  The source's extra condition on
\(z_w\) in \(N(E)\) is redundant for the same reason as in
RCV10, so this dictionary also respects every word witness.
\end{proof}
% END RCV13

% BEGIN RCV14
\paragraph{RCV14. Rational coordinates and reconstruction.}
\label{claim:RCV14}
For \(n\in N(E)\), define
\[
 Q_{ba}=h(f_b,ne_a),\quad A_{ba}=h(f_b,nf_a),\quad
 M=Q^{-1}A,\quad x_b=p_W(n^{-1}f_b).
\tag{RCV.24}
\]
Then \(Q\) is upper unitriangular, \(x_b\in W\), and, putting
\(G(x)_{ba}=h(x_b,x_a)\),
\[
               M+M^*=-H,\qquad H=Q^{-1}G(x)Q^{-*}.
\tag{RCV.25}
\]
Conversely every such rational tuple \((Q,x,M)\) reconstructs a
unique element of \(N(E)\).  At every chosen place its block
coordinates are \(\sigma(Q_{ba})\), \(\rho_w(x_b)\), the
upper entries of \(\operatorname{skew}\sigma(M)\), and
\(\operatorname{Im}\sigma(M_{aa})\).
This proves the original \texttt{Dict\allowbreak Gram}, \texttt{Dict\allowbreak Ex}, and
\texttt{Dict\allowbreak Place}.
\begin{proof}
The group conditions give the assertions about \(Q\).  For
\(y\in W\), preservation of the form gives
\[
 h(f_b,ny)=h(n^{-1}f_b,y)=h(x_b,y).
\]
Thus \(ny=y+\sum_bh(x_b,y)e_b\).  Put \(U=Q^{-*}\).
The complete action on the three rational summands is forced to be
\[
 \begin{split}
 ne_a&=\sum_b Q_{ba}e_b,\\
 ny&=y+\sum_b h(x_b,y)e_b\quad(y\in W),\\
 nf_a&=\sum_b(QM)_{ba}e_b
                         -\sum_b U_{ba}x_b+\sum_b U_{ba}f_b.
 \end{split}
\tag{RCV.26}
\]
The middle component of the last line is obtained by pairing with
all vectors \(ny\), \(y\in W\); nondegeneracy on \(W\)
makes it unique.  Pairing two last-line columns gives
\(M+M^*+U^*G(x)U=0\), which is (RCV.25).

Conversely, (RCV.26) defines a rational linear map for every tuple
as in the statement.  Its three diagonal blocks are \(Q,1,U\),
so it is invertible.  Pairings of its \(e\)-columns and its
\(W\)-columns are unchanged; their pairings with the \(f\)-columns
are unchanged by the two correction terms in the last line.
The remaining \(f,f\) pairing is zero precisely by (RCV.25).
Thus it is an isometry and has the required unipotent flag behavior.
For this map, pairing \(n^{-1}f_b\) with any \(y\in W\)
recovers \(h(x_b,y)\), so its projected vector is indeed \(x_b\).
All its readings (RCV.24) are therefore the prescribed tuple, proving
uniqueness and the exact reconstruction assertion.

Finally extend (RCV.26) to each chosen complex embedding.  Its row
entries on the frame of \(W_v\) are \(\rho_w(x_b)\).
The remaining readings agree directly with (RCV.15)--(RCV.17).
This proves the exact placement statement on the original rational
and adelic components, including the order of conjugation in the
row map.
\end{proof}
% END RCV14

% BEGIN RCV15
\paragraph{RCV15. Exact preservation of the original lattice.}
\label{claim:RCV15}
For a tuple in RCV14, the associated \(n\) preserves \(L_0\)
if and only if
\[
 \begin{split}
 Q_{ba}&\in\mathfrak a_b\mathfrak a_a^{-1}\quad(b<a),\\
 x_b&\in\mathfrak a_b L_b,\\
 M_{ba}&\in\mathfrak a_b\mathfrak a_a\quad(a,b<r).
 \end{split}
\tag{RCV.27}
\]
This is the full original \texttt{Dict\allowbreak Lat}.
\begin{proof}
Suppose first that \(nL_0=L_0\).  Apply \(n\) to
\(c e_a\), \(c\in\mathfrak a_a\), and extract the \(e_b\)
coefficient by pairing with \(f_b\).  Equation (RCV.2) gives
\(Q_{ba}\mathfrak a_a\subset\mathfrak a_b\), equivalently
the first line of (RCV.27).  Apply \(n^{-1}\) to
\(c f_b\), \(c\in\mathfrak a_b^{-1}\), and project onto
\(W\).  This gives \(\mathfrak a_b^{-1}x_b\subset L_b\),
equivalently \(x_b\in\mathfrak a_bL_b\); the equivalence may
be checked after localizing the invertible ideal at each prime.
Likewise applying \(n\) to \(\mathfrak a_a^{-1}f_a\)
gives \(A_{ba}\in\mathfrak a_b\mathfrak a_a\).
The finite series for \(Q^{-1}\) has entries
\((Q^{-1})_{bc}\in\mathfrak a_b\mathfrak a_c^{-1}\)
for \(b\leq c\), because the ideals telescope along each
matrix-product path.  Hence \(M=Q^{-1}A\) satisfies the last
line of (RCV.27).

Conversely assume (RCV.27).  The same inverse-series argument
gives the stated ideal bounds on \(Q^{-1}\), and therefore
\(U_{ba}\in\mathfrak a_b^{-1}\mathfrak a_a\) whenever it
is nonzero.  Formula (RCV.26) now preserves each lattice generator
in (RCV.2).  For \(y\in L_b\), self-duality and conjugation
invariance give \(h(x_b,y)\in\mathfrak a_b\).
For \(c f_a\), \(c\in\mathfrak a_a^{-1}\), its three
components have respectively coefficients in \(\mathfrak a_b\),
vectors in \(L_b\), and coefficients in \(\mathfrak a_b^{-1}\):
use \((QM)_{ba}\in\mathfrak a_b\mathfrak a_a\) and
\(U_{ba}x_b\in\mathfrak a_aL_b\).  Thus \(nL_0\subset L_0\).

Because \(n\) is an isometry and \(L_0\) is self-dual,
\((nL_0)^{\#}=nL_0\), where \(\#\) denotes Hermitian
duality.  Taking duals reverses the inclusion and gives
\(L_0\subset nL_0\).  Hence it is equality.  This final step
proves preservation, rather than only integral inclusion, for the
actual possibly nonprincipal lattice.
\end{proof}
% END RCV15

% BEGIN RCV16
\paragraph{RCV16. The complete shifted-fiber description of the arithmetic subgroup.}
\label{claim:RCV16}
In the block coordinates of RCV10, \(\Gamma_w\) has exactly the
source's \texttt{NOfib} description
\[
 t\in\operatorname{blk}_w(\Gamma_w)
 \quad\Longleftrightarrow\quad
 t_p-\tau_p(t_{<h(p)})\in cLat(X,w,p)\quad\hbox{for every }p.
\tag{RCV.28}
\]
The shifts of the \(X\) and row blocks are zero.  On every real
coordinate tuple form the upper unitriangular \(Q\) from its
\(X\)-blocks and the matrix \(B\) from its row blocks, and put,
at each place, \(H=Q^{-1}BG_w^{-1}B^*Q^{-*}\).
When these blocks are rationally recoverable, this is exactly
\(\sigma(Q^{-1}G(x)Q^{-*})\) by RCV5.  The other shifts are
\[
 \tau_{K_{ba}}=\tfrac12 H_{ba},\qquad
 \tau_{D_a}=-\operatorname{Re}H_{aa}
                                  \operatorname{Im}\sigma(\theta).
\tag{RCV.29}
\]
They are the literal \(\mathrm{tauQ}\) of the source, depend only
on lower-height blocks, and vanish at zero.  Both directions of
(RCV.28), including existence of a rational lift of every admissible
coordinate tuple, hold.  More precisely, for every rational
\(\mathrm{DataOK}\) tuple the shifted-lattice conditions on its
\(\mathrm{blkQ}\) are equivalent to \(\mathrm{LatQ}\), and
every real tuple satisfying those conditions is the \(\mathrm{blkQ}\)
of a rational \(\mathrm{DataOK}\) tuple.  These are the full
original \texttt{Fib\allowbreak Q} and \texttt{Fib\allowbreak Onto}.
\begin{proof}
The ideal bounds of (RCV.27) imply
\(G(x)_{ba}\in\mathfrak a_b\mathfrak a_a\), and the
inverse-matrix bounds used in RCV15 then imply the same bound for
\(H_{ba}\).  From (RCV.25),
\[
 \operatorname{skew}(M)_{ba}=M_{ba}+H_{ba}/2.
\tag{RCV.30}
\]
Thus the upper off-diagonal lattice condition on \(M\) is
exactly the \(K\)-condition in (RCV.28).  Its lower entries are
forced by \(M_{ab}=-H_{ab}-\overline{M_{ba}}\), and obey
the corresponding ideal condition as well.

On the diagonal, \(H_{aa}\in F\), and
\[
                  \xi_a=M_{aa}+H_{aa}\theta
\]
is anti-invariant because \(M_{aa}+\bar M_{aa}=-H_{aa}\)
and \(\theta+\bar\theta=1\).  It lies in \(\mathfrak a_a^2\)
exactly when \(M_{aa}\) does, since \(\theta\) is integral.
Its imaginary coordinate is
\(\operatorname{Im}\sigma(M_{aa})+
 \operatorname{Re}\sigma(H_{aa})\operatorname{Im}\sigma(\theta)\).
This gives exactly the diagonal shift (RCV.29), with its stated
sign.  RCV13--RCV15 prove the forward implication.

For the converse, the unshifted \(X\) blocks uniquely recover
rational entries of \(Q\) in the specified ideals.  The row
blocks uniquely recover \(x_b\in\mathfrak a_bL_b\), by
injectivity of (RCV.8).  They give the same rational \(H\) as
above.  For each \(b<a\), the shifted \(K\)-coordinate
recovers a unique \(M_{ba}\in\mathfrak a_b\mathfrak a_a\).
Set \(M_{ab}=-H_{ab}-\overline{M_{ba}}\).  Each diagonal
condition recovers an anti-invariant \(\xi_a\in\mathfrak a_a^2\);
set \(M_{aa}=\xi_a-H_{aa}\theta\).  These formulas give
every entry of a rational matrix satisfying (RCV.25) and (RCV.27).
RCV14 reconstructs its isometry, RCV15 proves lattice preservation,
and RCV13 places its actual infinite component in \(\Gamma_w\).
Its coordinates are exactly the supplied tuple, by (RCV.30) and
the diagonal computation.  This proves surjectivity onto the entire
shifted-fiber set, not merely onto coordinates of points already
known to be rational.

For the stated \texttt{Fib\allowbreak Q} equivalence, a rational tuple failing
an \(X\)- or row-lattice condition fails the same unshifted
condition on both sides.  Otherwise the computations (RCV.30) and
the diagonal identity prove the equivalence entry by entry, for
every \(\mathrm{DataOK}\) tuple.  The reconstruction in the
preceding paragraph is exactly the universal \texttt{Fib\allowbreak Onto}
return.  All \(X\)- and row heights are below all central heights.
Hence the actual polynomial shifts have the required strict
lower-height dependence.  They are measurable as well, although
the original \texttt{NOfib} does not require this extra property.
Finally, subgroup closure and inclusion in \(N_\infty\) were
proved in RCV13.  These supply the other fields of \texttt{NOfib}.
\end{proof}
% END RCV16

\subsection{All fundamental domains and all supplied Haar measures}

% BEGIN RCV17
\paragraph{RCV17. A triangular lattice has the product covolume.}
\label{claim:RCV17}
For the actual group of RCV11 and subgroup of RCV16, choose any
half-open fundamental parallelotope \(P_p\) for each block lattice.
The Cartesian product \(P=\prod_p P_p\) is a strict fundamental
domain in block coordinates.  Every measurable strict domain
\(A\) for this subgroup has
\[
                     \operatorname{vol}A=\prod_p cVal(X,w,p).
\tag{RCV.31}
\]
Consequently, for every original ambient
\(\operatorname{IsSFund}(\Gamma_w,\Omega)\),
\(\kappa_w(\Omega\cap N_\infty)\) is measurable and satisfies
(RCV.31).
\begin{proof}
Order the blocks by increasing height.  The shifted-fiber condition
and its surjectivity in RCV16 show that, for every block \(p\)
and every \(\lambda\in cLat_p\), there is \(\gamma\in\Gamma_w\)
whose preceding blocks are zero and whose \(p\)-th block is
\(\lambda\).  To construct it, set those earlier coordinates to
zero, the current coordinate to \(\lambda\), and choose each
later coordinate in its nonempty shifted lattice.  All earlier
conditions hold because the shifts vanish at zero.

Given an arbitrary point of the real coordinate space, reduce its
first coordinate into \(P_1\) by left multiplication by such a
\(\gamma\).  At the next step use an element with all preceding
coordinates zero.  Equation (RCV.19) and its zero-argument property
say that this leaves those coordinates fixed and translates the
current one by exactly \(\lambda\).  After finitely many steps
the point lies in \(P\).  For uniqueness, suppose \(x,y\in P\)
and \(y=\gamma*x\).  At the first coordinate the difference
is in its lattice, so the strict parallelotope condition gives
\(\gamma_1=0\) and \(x_1=y_1\).  Inductively the earlier
coordinates of \(\gamma\) vanish, so its current coordinate
belongs to the unshifted lattice and the correction in (RCV.19)
vanishes.  The same argument gives its vanishing.  Thus
\(\gamma=1\) and \(x=y\).  This proves that \(P\) is a
strict domain even though the subgroup itself has shifted fibers.

Its product Lebesgue volume is \(\prod_pcVal_p\), by RCV9 and
Tonelli.  The subgroup is countable by RCV13, and left translations
preserve Lebesgue measure by RCV11.  Partition any other measurable
strict domain \(A\) by its intersections with the translates of
\(P\), and translate the pieces back.  They partition \(P\),
exactly as in RCV2.  This proves (RCV.31) before assuming finiteness
of \(A\).

An original \texttt{Is\allowbreak SFund} is a measurable set in the ambient
archimedean group; its defining translates cover that whole group
and meet it uniquely.  Because \(\Gamma_w\subset N_\infty\),
its intersection with \(N_\infty\) has the same strict-domain
properties on \(N_\infty\).  Its image is measurable by the
inverse chart in RCV10.  Finally \(\kappa_w\) only separates
complex entries into real and imaginary coordinates and permutes
them.  It preserves precisely the block product Lebesgue measure.
Thus (RCV.31) applies to every original ambient domain witness.
\end{proof}
% END RCV17

% BEGIN RCV18
\paragraph{RCV18. Multiplication of the four exact block constants.}
\label{claim:RCV18}
Put
\[
 e_N=-\frac{8r^3+6r^2+6kr^2+r}{6},\qquad
 C_w=T_w^{rd}2^{-d(rk+r^2)}\Delta^{e_N}.
\tag{RCV.32}
\]
Then \(\prod_p cVal(X,w,p)=C_w\), and hence the full
original \texttt{Coord\allowbreak Covol} holds with its literal return
\(\mathrm{covolN}(X,w)/\mathrm{twoU}=C_w\).
\begin{proof}
Only the \(r\) row blocks contain \(T_w\); their product gives
\(T_w^{rd}\).  There are two families of \(r(r-1)/2\)
complex ideal blocks, and \(r\) row blocks of complex dimension
\(dq\).  Thus the power of two in the product is
\(-d(r(r-1)+rq)=-d(rk+r^2)\).

We record the discriminant sum, so that no root count is hidden.
Using \(\sum_{a=0}^{r-1}a=r(r-1)/2\) and
\(\sum_{a=0}^{r-1}a^2=r(r-1)(2r-1)/6\), the four sums are
\[
\begin{array}{c|c}
 X&-r(r-1)(2r-1)/6\\
 R&-qr^2\\
 K&r(r-1)(2r-1)/2-2r^2(r-1)\\
 D&-r^2-r/2.
\end{array}
\tag{RCV.33}
\]
For example the first is
\(\sum_{a=0}^{r-1}\sum_{b<a}(1-2(a-b))\); the third is
\(\sum_{a=0}^{r-1}\sum_{b<a}(1-4r+2a+2b)\).
Adding the four displayed sums and substituting \(q=k+1\)
gives exactly \(e_N\).  They are identities also at \(r=0\).
Positivity of \(T_w,2,\Delta\) permits ordinary addition of
the real exponents.

The source defines
\[
 \mathrm{covolN}(X,w)
 =T_w^{rd}2^{-d(rk+r(r-1)/2)}\Delta^{e_N},\qquad
 \mathrm{twoU}=2^{dr(r+1)/2}.
\]
Since \(r(r-1)/2+r(r+1)/2=r^2\), their quotient is \(C_w\).
Combine this arithmetic with RCV17, including its measurability
and every-domain statements, to obtain the complete original
\texttt{Coord\allowbreak Covol}.  Its right side is finite and strictly positive.
\end{proof}
% END RCV18

% BEGIN RCV19
\paragraph{RCV19. The original Haar formula with every normalization.}
\label{claim:RCV19}
Let \(\mu\) be any ambient measure satisfying the original
\(\mathrm{IsHaarOnSet}(N_\infty,\mu)\): it is sigma-finite,
nonzero, supported on \(N_\infty\), and invariant under its
left translations.  Let \(\mathrm{Box}_w\) be the original
\(\mathrm{boxN}(X,w)\), obtained by placing every real coordinate
of \(\kappa_w\) in \([0,1)\).  Then
\[
 (\kappa_w)_*\mu=\mu(\mathrm{Box}_w)\,d\kappa,
 \qquad 0<\mu(\mathrm{Box}_w)<\infty.
\tag{RCV.34}
\]
For every original \(\operatorname{IsSFund}(\Gamma_w,\Omega)\),
\[
                 \mu(\Omega)=C_w\,\mu(\mathrm{Box}_w).
\tag{RCV.35}
\]
These are the full \texttt{Coord\allowbreak Haar} and \texttt{Covol\allowbreak N},
with their original \([0,\infty]\)-valued equalities understood
via \(\operatorname{ofReal}\) on the finite constants.
\begin{proof}
The coordinate Lebesgue measure transported through RCV10 is a
nonzero locally finite Haar measure by RCV11.  The subgroup is
closed and second countable.  AMD8 applies at exactly the weak
\texttt{Is\allowbreak Haar\allowbreak On\allowbreak Set} scope above: restriction to the supported
closed subgroup is a genuine Haar measure, and uniqueness gives
one finite strictly positive scalar multiplying the coordinate
measure.  The coordinate unit cube has volume one, including
in dimension zero.  Evaluating at its inverse image identifies
that scalar as \(\mu(\mathrm{Box}_w)\), proving (RCV.34).
Support on \(N_\infty\) justifies this assertion for the ambient
map \(\kappa_w\), including its values off the subgroup.

The same support gives \(\mu(\Omega)=\mu(\Omega\cap N_\infty)\).
Apply (RCV.34) and RCV18 to the measurable coordinate image of
this set.  This gives (RCV.35).  In particular its finiteness and
positivity are conclusions.  The proof has fixed neither a Haar
normalization nor a fundamental domain; it applies to every
witness in the original statement.
\end{proof}
% END RCV19

\subsection{The prescribed Lie measure and the independent reference word}

% BEGIN RCV20
\paragraph{RCV20. The Euclidean Lie-algebra volume factor.}
\label{claim:RCV20}
For arbitrary nonnegative integers \(r,k\), take the source's
standard group \(N_{k,r}\) with middle Gram matrix
\(D_0=\operatorname{diag}(1_k,-1)\).  Its actual realified
Lie space consists exactly of
\[
 L(A,B,K)=
 \begin{pmatrix}
 A&B&K\\ 0&0&-D_0B^*\\0&0&-A^*
 \end{pmatrix},\qquad
 A\text{ strictly upper triangular},\quad K^*=-K.
\tag{RCV.36}
\]
If the independent complex coordinates are the upper entries of
\(A,K\) and all entries of \(B\), and the real coordinates
are \(\operatorname{Im}K_{aa}\), its Euclidean volume is
\[
        2^{r^2+rk}\,L_*\!\left(dA\,dB\,dK_{\mathrm{upper}}
                                      \,d\operatorname{Im}K_{\mathrm{diag}}\right).
\tag{RCV.37}
\]
This proves the complete original \texttt{Const\allowbreak Step}.
\begin{proof}
Differentiate the standard upper-flag and isometry equations for
\(\exp(tZ)\) at zero.  They force (RCV.36).  Conversely a matrix
of that form is nilpotent and skew-adjoint for the standard
hyperbolic Hermitian matrix.  Its finite exponential preserves
that form, preserves the upper flag and acts trivially on its
specified quotients.  Thus \(\exp(tZ)\in N_{k,r}\) for every
real \(t\).  This proves equality with the source's
\(\mathrm{lieOf}(N_{k,r})\), and hence with its real linear
span \(\mathrm{lieE}(N_{k,r})\); no larger tangent space has
been substituted.

The source's ambient Euclidean norm is the sum of the squares of
the real and imaginary parts of every matrix entry.  Therefore
\[
 \|L(A,B,K)\|^2
  =2\left(\sum_{b<a}|A_{ba}|^2+
           \sum_{b,j}|B_{bj}|^2+
           \sum_{b<a}|K_{ba}|^2\right)
                          +\sum_a(\operatorname{Im}K_{aa})^2.
\tag{RCV.38}
\]
Indeed the \(A\) entries occur again in \(-A^*\), the \(B\)
entries again in \(-D_0B^*\), and every off-diagonal entry of
\(K\) has its conjugate partner.  A diagonal imaginary entry
occurs once.  Polarization shows that distinct real coordinate
directions are orthogonal, with length \(\sqrt2\) for each
real or imaginary part of a complex parameter and length one
for a diagonal parameter.  The number of complex parameters is
\[
              2\binom r2+r(k+1)=r^2+rk.
\]
The real volume factor is consequently
\((\sqrt2)^{2(r^2+rk)}=2^{r^2+rk}\), proving (RCV.37).
The coordinate linear map is a measurable isomorphism onto the
actual Lie space, as the entries in (RCV.36) recover every parameter.
This supplies both the map and the volume equality demanded by
\texttt{Const\allowbreak Step}, including the empty-coordinate case.
\end{proof}
% END RCV20

% BEGIN RCV21
\paragraph{RCV21. The exponential-to-polynomial chart has determinant one.}
\label{claim:RCV21}
Let \(\Phi\) denote (RCV.15) in the standard normalized frame
\(T=1\).  Then the source's prescribed measure satisfies
\[
 \mathrm{haarLie}(N_{k,r},1)
   =2^{r^2+rk}\,\Phi_*d\kappa_0
                     \quad(r,k\geq0).
\tag{RCV.39}
\]
In particular the full original \texttt{Shear\allowbreak Step} and
\texttt{Exp\allowbreak Chart} hold, with no free proportionality constant.
\begin{proof}
Assign degrees \(b\), \(r\), and \(2r-b\) to the standard
basis vectors \(e_b,W_j,f_b\), respectively.  A potentially
nonzero entry of (RCV.36) increases this degree strictly from its
row to its column.  Therefore it is nilpotent, and its exponential
and logarithm are finite polynomial sums.  Assign the independent
coordinates the corresponding positive weights
\[
 a-b\text{ for }A_{ba},\quad r-b\text{ for }B_{bj},\quad
 2r-b-a\text{ for }K_{ba},\quad 2r-2a\text{ on the diagonal}.
\tag{RCV.40}
\]
Every term of degree at least two in an exponential entry follows
a path of at least two positive steps.  Each factor therefore has
strictly smaller weight than the resulting entry.  It follows that
the \(Q\)- and \(B\)-coordinates read from \(\exp L\) are
their linear parameters plus polynomials in smaller-weight parameters.
For the central coordinate, reading \(M=Q^{-1}C\) introduces
only products of positive-weight entries.  Subtracting its Hermitian
part, equivalently reading \(\operatorname{skew}(M)\), again
introduces only such products.  Its linear term is exactly \(K\),
and on the diagonal it is exactly its imaginary parameter.

Thus the map \(s=\Phi^{-1}\circ\exp\circ L\), written in
real and imaginary coordinates and ordered by (RCV.40), has the form
\(s_j(t)=t_j+P_j(t_{\text{smaller weight}})\).  This statement
also holds across equal weights, since every nonlinear factor has
strictly smaller weight.  Successive subtraction gives a polynomial
inverse.  As in RCV11, repeated translation in one coordinate block
and Tonelli show that \(s\) preserves Lebesgue measure exactly.
It is the measurable, measure-preserving map required by
\texttt{Shear\allowbreak Step}; by construction \(\Phi(s(t))=\exp L(t)\).

The definition of \(\mathrm{haarLie}(N_{k,r},1)\) is the
pushforward of the Euclidean volume of its actual realified Lie
space under matrix exponential.  Insert (RCV.37), then use
\(\exp L=\Phi\circ s\) and \(s_*d\kappa_0=d\kappa_0\).
This gives (RCV.39) with the exact coefficient shown.  It also
explains why replacing the exponential chart by the polynomial
chart does not introduce an additional Jacobian factor.
\end{proof}
% END RCV21

% BEGIN RCV22
\paragraph{RCV22. The full local reference-measure formula.}
\label{claim:RCV22}
Take any integers \(n,r,k\) with \(n=2r+k+1\), any signature
\(\epsilon\), and any supplied frame \(e,f,w_1,\ldots,w_k,z\)
in \(\mathbb C^n\) satisfying the hyperbolic pairings,
orthonormal positive \(w_j\), all stated orthogonalities, and
\(h_\epsilon(z,z)=-T\) with \(T>0\).
Let \(\nu\) be the literal source measure
\(\mathrm{refNg}(n,\epsilon,e,f,w,z/\sqrt T)\).
Then the full source coordinate map \(\kappa_{e,f,w,z}\) is
\(\nu\)-almost-everywhere measurable and
\[
       (\kappa_{e,f,w,z})_*\nu
                =2^{r^2+rk}T^{-r}\,d\kappa.
\tag{RCV.41}
\]
Furthermore \(\nu\) is an original \texttt{Is\allowbreak Haar\allowbreak Loc} measure
on \(N(\epsilon,e,z')\) for every \(z'\perp e\), without
a norm or nonzero assumption on \(z'\).
\begin{proof}
The normalized frame \((e,w,z/\sqrt T,f)\) has the source's
standard Gram matrix.  Its column matrix is invertible by
nondegeneracy of that matrix.  The corresponding conjugation
\(g\mapsto P^{-1}gP\) is a closed measurable embedding of the
actual unitary group into standard matrices, with the expected
inverse on its image.  The source defines \(\nu\) as the
pullback of \(\mathrm{haarLie}(N_{k,r},1)\) along precisely
this map.  Equation (RCV.39) is supported on \(N_{k,r}\), all
of which lies in this image.  Thus it may equivalently be pushed
forward along the inverse conjugation, with no discarded mass.

In the normalized chart, the readings of \(Q\) and \(K\)
are unchanged.  The final entry in each row block is
\(h(f_b,gz)=\sqrt T\,h(f_b,g(z/\sqrt T))\).
There are \(r\) such complex entries.  Multiplying each by
\(\sqrt T\) has real determinant \((\sqrt T)^{2r}=T^r\).
The pushforward of Lebesgue measure under this invertible linear
map is \(T^{-r}\) times Lebesgue measure.  Combining this
with (RCV.39) proves (RCV.41).  The coordinate identities hold
on the supporting closed subgroup; hence they give the required
almost-everywhere measurability for the original ambient reading.

The measure in (RCV.41) is finite on compact subsets, nonzero,
and sigma-finite.  It is supported on the actual unipotent subgroup,
and RCV11 proves left invariance in these coordinates.  The extra
condition \((g-1)z'\in\langle e\rangle\) in the definition of
\(N(\epsilon,e,z')\) follows from \(z'\in e^\perp\).
Therefore this subgroup is the same one for every stated \(z'\).
These observations prove every field of the broad original
\texttt{ref\allowbreak Ng\_\allowbreak is\allowbreak Haar\allowbreak Loc} statement as well as the full
\texttt{ref\allowbreak N\_\allowbreak map\_\allowbreak kap\allowbreak L} formula.
\end{proof}
% END RCV22

% BEGIN RCV23
\paragraph{RCV23. Changing the word: the exact row determinant on all matrices.}
\label{claim:RCV23}
For every \(X\), every pair of words \(u,w\), and every
chosen infinite place \(v\), there is an invertible real linear
map \(L_{u,w,v}\) on the full local coordinate space such that
\[
 \kappa_{w,v}(s)=L_{u,w,v}\kappa_{u,v}(s)
       \quad\hbox{for every ambient unitary element }s,
 \qquad |\det L_{u,w,v}|=(T_w/T_u)^r.
\tag{RCV.42}
\]
The same coordinate identity holds for every ambient matrix on which
the readings are defined by the source's total inverse convention.
In particular this is the full original \texttt{Frame\allowbreak Row}.
\begin{proof}
Both \(S_u=(w_{u,v,1},\ldots,w_{u,v,k},z_u)\) and
\(S_w=(w_{w,v,1},\ldots,w_{w,v,k},z_w)\) are bases of the
same complement \(W_v\).  Write \(S_w=S_u C\).  Their Gram
identities give
\[
       C^*\operatorname{diag}(1_k,-T_u)C
                           =\operatorname{diag}(1_k,-T_w).
\]
Taking determinants proves \(|\det_{\mathbb C}C|^2=T_w/T_u\).
For an arbitrary ambient matrix \(s\), the functional
\(y\mapsto h(f_b,sy)\) is complex linear.  Its row of values
on \(S_w\) is therefore its row of values on \(S_u\), multiplied
on the right by the same matrix \(C\).  This proves the row
transformation without any unipotence assumption on \(s\).

All \(Q\), \(A\), \(M=Q^{-1}A\), and central coordinate
readings use only \(e,f,s\).  They are independent of the word,
even when the total matrix inverse is evaluated off its nonsingular
locus.  Thus \(L_{u,w,v}\) is the identity on those coordinates
and right multiplication by \(C\) on each of the \(r\) row
blocks.  A complex linear map has real determinant
\(|\det_{\mathbb C}C|^2\), as follows for instance by writing
its real block matrix.  The product over the \(r\) rows is
\((T_w/T_u)^r\), proving (RCV.42).
No relation between the separately supplied orthonormal frames is
assumed; their Gram matrices alone determine this determinant.
\end{proof}
% END RCV23

% BEGIN RCV24
\paragraph{RCV24. The complete independent-word reference-box formula.}
\label{claim:RCV24}
For every \(X\) and every independently chosen pair \(u,w\),
let
\[
 \nu_u=\prod_{v\mid\infty}\mathrm{refNX}(X,u,v).
\]
Then \(\nu_u\) is an original \texttt{Is\allowbreak Haar\allowbreak On\allowbreak Set} measure
on \(N_\infty(w)\), and
\[
       \nu_u(\mathrm{Box}_w)
                 =T_w^{-rd}\,2^{d(rk+r^2)}.
\tag{RCV.43}
\]
This proves the full original \texttt{Ref\allowbreak NBox}, including the
reference word different from the lattice word.
\begin{proof}
At each place, RCV22 with the frame of \(u\) gives density
\(2^{r^2+rk}T_u^{-r}\) in the \(u\)-coordinates.  Apply
the linear map in RCV23.  Pushforward of Lebesgue measure divides
the density by its absolute determinant, so in the
\(w\)-coordinates it becomes
\[
 2^{r^2+rk}T_u^{-r}(T_w/T_u)^{-r}
                              =2^{r^2+rk}T_w^{-r}.
\]
This equality keeps the same local reference measure throughout.
The local \(w\)-cube has Lebesgue volume one.  Taking the finite
product over the \(d\) specified places yields (RCV.43), since
the original \(\mathrm{Box}_w\) is exactly the product of these
local boxes.

Each local measure is sigma-finite, nonzero, supported on the local
group, and invariant under its left translations by the final
assertion of RCV22.  Finite products preserve all four properties;
nonzeroness can also be read from the positive finite box mass just
computed.  This proves the full original product Haar assertion
\texttt{ref\allowbreak NX\_\allowbreak haar} for every \(u,w\), not merely for equal
words.  Empty row or pair index sets contribute the factors one.
\end{proof}
% END RCV24

% BEGIN RCV25
\paragraph{RCV25. Full reference covolume and the original discriminant return.}
\label{claim:RCV25}
For every original \(X\), every reference word \(u\), every
lattice word \(w\), and every supplied
\(\operatorname{IsSFund}(\mathrm{NOinf}(X,w),\Omega)\),
\[
                  \nu_u(\Omega)=\Delta^{e_N}
                         \quad\hbox{in }(0,\infty).
\tag{RCV.44}
\]
In particular, for every original bridge \(hb\), every two
supplied domains \(\Omega,\Omega'\), and every \(u,u'\)
satisfying the original \(\mathrm{HypX}(X,hb,\Omega,\Omega',u,u')\),
\[
 (\nu_u(\Omega))^{\mathrm{toReal}}
 (\nu_u(\Omega'))^{\mathrm{toReal}}
 =|\operatorname{disc}(X.B.I.F)|^{e_{DF}},\qquad
 e_{DF}=-\frac{8r^3+6r^2+6kr^2+r}{3}.
\tag{RCV.45}
\]
This proves the full original \texttt{covol\_\allowbreak ref\allowbreak N} and
\texttt{Covol\allowbreak Ref}, with no retained covolume, row-lattice, or
finite-level membership premise.
\begin{proof}
Apply RCV19 to the unchanged measure \(\nu_u\), using its
original Haar fields from RCV24.  Equations (RCV.32), (RCV.35),
and (RCV.43) give
\[
 \begin{split}
 \nu_u(\Omega)
 &=\big(T_w^{rd}2^{-d(rk+r^2)}\Delta^{e_N}\big)
                       \big(T_w^{-rd}2^{d(rk+r^2)}\big)\\
 &=\Delta^{e_N}.
 \end{split}
\]
Every cancelled factor is finite and strictly positive; these
properties follow from \(T_w>0\) and \(\Delta>0\), so there
is no cancellation of zero or infinite extended-real masses.
This proves (RCV.44) for every domain in the broad statement.

The original \(\mathrm{HypX}\) supplies exactly the two required
strict-domain assertions: \(\Omega\) for
\(\mathrm{NOinf}(X,u)\) and \(\Omega'\) for
\(\mathrm{NOinf}(X,u')\).  Apply (RCV.44) twice, with reference
word \(u\) both times, and lattice words \(u\) and \(u'\)
respectively.  Finite positive masses have their ordinary real values
under \(\mathrm{toReal}\).  Their product is
\(\Delta^{2e_N}=\Delta^{e_{DF}}\), giving (RCV.45).
The bridge and the other fields of \(\mathrm{HypX}\) are accepted
at their original scope; this calculation does not need to constrain
or reselect them.

The proof has thus supplied the exact reference volume that the
source's subsequent \texttt{place\allowbreak Split\_\allowbreak of} consumes.  That next
theorem still has its separate local \texttt{Jac\allowbreak Gauge} input.
The present result does not assert its Jacobian or the full
\texttt{Real\allowbreak Bracket}, and it does not reconstruct \texttt{Forms0}
or any missing global spectral producer.
\end{proof}
% END RCV25

\subsection{What has been fixed for the next comparison}

The completed formulas have distinct purposes.  Equation (RCV.35)
works with every supplied Haar normalization and every original
arithmetic domain.  Equation (RCV.43) computes the fixed reference
normalization, with its reference word independent of the word in
the box.  Their cancellation gives the discriminant in (RCV.45).
This distinction is necessary in the next step: local disintegration
must compare the actual measures with these reference measures,
while keeping the same setting, two words, domains, and bridge.

The remaining analytic card starts at the original
\texttt{Jac\allowbreak Gauge} and its all-section local disintegration formula.
It may then use the proved \texttt{Covol\allowbreak Ref} in
\texttt{place\allowbreak Split\_\allowbreak of}, and the independently proved or separately
reconstructed product-disintegration supplier in the exact
\texttt{Real\allowbreak Bracket} consumer.  Neither a change of an infinity
measure nor the already completed actual \texttt{In\allowbreak Volumes}
can substitute for that local Jacobian calculation.


\section{The prescribed cyclotomic prism and the full Wach-existence return}
\label{sec:wach-prism-v056}

\status{This chapter constructs the actual prism required by
\texttt{exists\_\allowbreak wach\_\allowbreak prism\_\allowbreak over\_\allowbreak integers}, including its original
scalar map and both coefficient-coordinate assertions. It uses the
already reconstructed integral period rings and invariant matrix
descent. The final return proves the full original Wach-existence
implication under its original tied-period and crystalline-realization
inputs. Later faithfulness, cohomology and quantitative-comparison
producers remain separate.}

The construction has a particularly concrete underlying ring. Let
$p$ be the original prime, $K\subset\overline{\mathbb Q}_p$ the
original finite extension, $K_0$ its maximal unramified subfield,
and put
\[
 W_0=\mathcal O_{K_0},\qquad O=\mathcal O_{LL},\qquad
 \mathfrak S=W_0[[T]],\qquad A_{\mathrm{inf}}=W(\mathcal O_C^\flat).
 \tag{WPR.1}
\]
For the prism endpoint keep the supplied
$P,AI,\mathrm{Good}(K,s,LL),\Gamma,R,I$ and
$\mathrm{IoW}(I)$. In particular $s\geq1$ and $K\subset K_s$,
where $K_s=K_0(\mu_{p^s})$. No new field or cyclotomic parameter
is substituted for these data. The intermediate statements keep
their own printed scopes: in particular WPR11--WPR15 require
no $\mathrm{Good}$ hypothesis, and WPR15 holds for every $s\geq0$. Write $n=[K_0:\mathbb Q_p]$ and
$e_i:K_0\hookrightarrow LL$, $i\in\mathbb Z/n\mathbb Z$, for
the supplied embeddings. The actual coefficient map is
\[
 a:\mathfrak S\longrightarrow\prod_i O[[T]],\qquad
 a(f)_i=e_i(f),
 \tag{WPR.2}
\]
with each $e_i$ applied coefficientwise. The target is the original
$\mathrm{LetA}$, not an independently chosen product.

WIS1--WIS5 supply the complete integral coefficient rings, actual
Witt coordinates and inclusions. IFD9 supplies the unramified
valuation facts and finite cyclotomic integral calculation. LS4
supplies the actual Fontaine map $\theta:A_{\mathrm{inf}}\to\mathcal O_C$
and its principal kernel. These inputs retain their original maps.
We spell out the additional algebra and every prism field below.

\subsection{The unramified constants and the original power-series action}

% BEGIN WPR1
\paragraph{WPR1: unramified integral coordinates and the coefficient Frobenius.}
\label{claim:WPR1}
The ring $W_0$ is a complete DVR with maximal ideal $pW_0$.
It has a finite $\mathbb Z_p$-basis consisting of roots of unity
of order prime to $p$. Consequently every $w\in W_0$ is a
finite $\mathbb Z_p$-linear combination of such roots. The
original coefficient endomorphism $\sigma:W_0\to W_0$ satisfies
\[
                  \sigma(w)=w^p+p b_w\qquad(b_w\in W_0).
 \tag{WPR.3}
\]

\begin{proof}
The complete-DVR and uniformizer assertions are WIS1 and IFD9,
on this same finite unramified subfield. Let its residue field be
$k_0=\mathbb F_{p^f}$. Choose an $\mathbb F_p$-basis
$\bar\omega_1,\ldots,\bar\omega_f$ of $k_0$; each basis element
is nonzero. IFD9 constructs its unique lift $\omega_i$ satisfying
$\omega_i^{p^f-1}=1$. In particular its order is prime to $p$.

For any $w=w_0\in W_0$, choose digits $d_{i,r}\in\{0,\ldots,p-1\}$
successively so that
\[
 w_r=\sum_i d_{i,r}\omega_i+p w_{r+1}.
\]
The residue-basis property permits this step, and divisibility
by $p$ follows from the actual maximal ideal. Thus
\[
 w=\sum_i\omega_i\sum_{r<N}p^r d_{i,r}+p^Nw_N.
\]
All $w_N$ are integral, so the remainder tends to zero. Completeness
gives $w=\sum_i c_i\omega_i$ with
$c_i=\sum_{r\geq0}p^r d_{i,r}\in\mathbb Z_p$. If such a sum is
zero, reducing modulo $p$ makes every $c_i$ divisible by $p$.
Cancel $p$ and repeat; $p$-adic separation of $\mathbb Z_p$ gives
$c_i=0$. This proves the claimed basis and the original finite
root-of-unity spanning statement.

The map $\sigma$ is the restriction of the supplied
$P.\mathrm{PD}(K).\sigma$, transported to the integral ring by
the original $\mathrm{embW}$. It preserves integral elements.
The defining residue-Frobenius property of that datum says
$|\sigma(w)-w^p|<1$ for integral $w$. Since the value group of
$K_0$ is $p^{\mathbb Z}$, every such difference is in $pW_0$.
This gives (WPR.3), including the actual coefficient $b_w$.
\end{proof}
% END WPR1

% BEGIN WPR2
\paragraph{WPR2: the integral cyclotomic character at its full original scope.}
\label{claim:WPR2}
There is a continuous character
$\chi:G_{\mathbb Q_p}\to\mathbb Z_p^\times$ whose image in
$\mathbb Q_p^\times$ is the supplied $P.\chi_Q$. For every finite
$K$, every $r\geq0$ with $K\subset K_r$, and every
$u\in\mathbb Z_p^\times$ satisfying $p^r\mid u-1$, there is
$g\in G_K$ with $\chi(g)=u$. This is the full original
\texttt{exists\_\allowbreak integral\_\allowbreak cyclotomic\allowbreak Char} conclusion.

\begin{proof}
The supplied cyclotomic predicate has the following concrete
content: if $\zeta^{p^m}=1$, there is a natural $b$ such that
$g\zeta=\zeta^b$ and $|P.\chi_Q(g)-b|\leq p^{-m}$.
For $m=0$ this gives $|P.\chi_Q(g)|\leq1$. Applying it to
$g^{-1}$ and using the character identity gives equality.
Thus $P.\chi_Q(g)$ and its inverse are integral, defining the
unique lift $\chi(g)\in\mathbb Z_p^\times$. Its multiplication
law follows by injectivity into $\mathbb Q_p$.

For a primitive $p^m$-th root, the same predicate says that
$\chi(g)\bmod p^m$ is exactly its exponent under $g$. The
stabilizer of that finite extension is open in the original
Galois topology. Therefore each finite congruence of $\chi$ is
locally constant, proving continuity into the inverse limit
$\mathbb Z_p^\times$. This also proves
\[
 g\text{ fixes }\mu_{p^m}
 \quad\Longleftrightarrow\quad \chi(g)\equiv1\pmod{p^m}.
 \tag{WPR.4}
\]

For surjectivity with the stated restriction, the degree and
noncancellation calculation in IFD9 gives
$[K_0(\mu_{p^m}):K_0]=p^{m-1}(p-1)$ for $m\geq1$.
Hence every unit exponent modulo $p^m$ defines an automorphism
of this field over $K_0$. The exponents supplied by $u$ are
compatible as $m$ varies, so they define an automorphism of
$K_\infty=\bigcup_mK_0(\mu_{p^m})$. Extend it to the algebraic
closure: the usual extension of an embedding along algebraic
elements, or its maximal-extension proof, produces an automorphism
of the algebraic closure because the target is algebraically closed
and the extension is normal. It fixes $K_0$, and (WPR.4) shows
that it fixes $K_r$ when $u\equiv1\pmod{p^r}$. It therefore fixes
the supplied $K\subset K_r$. Its character agrees with $u$
modulo every $p^m$, and separatedness gives $\chi(g)=u$.
For $r=0$ the assertion just uses $K\subset K_0$ and the same
construction. No odd-prime restriction was used.
\end{proof}
% END WPR2

% BEGIN WPR3
\paragraph{WPR3: the full cyclotomic power-series action.}
\label{claim:WPR3}
For every commutative $\mathbb Z_p$-algebra $D$ and $c\in\mathbb Z_p$,
put
\[
 B_c(T)=\sum_{j\geq0}\binom cjT^j,\qquad
 \beta_c(T)=B_c(T)-1,\qquad
 \gamma_c(f)=f(\beta_c(T)),\qquad \phi_0=\gamma_p.
 \tag{WPR.5}
\]
Coefficients are first formed in $\mathbb Z_p$ and then mapped
to $D$. These are the source's actual $\mathrm{cyc}$ and
$\mathrm{act}$. They satisfy all its action laws:
\[
 \begin{gathered}
 \gamma_c(T)=\beta_c(T),\quad
 \gamma_c\gamma_d=\gamma_{cd},\quad \gamma_1=1,\quad
 \beta_k(T)=(1+T)^k-1\quad(k\in\mathbb N),\\
 [T]\beta_c=c,\qquad
 \gamma_c\big((1+T)^k-1\big)
       =\beta_c\big((1+T)^k-1\big).
 \end{gathered}                                                     \tag{WPR.6}
\]
The constant coefficient is preserved. Coefficientwise base change
along every $\mathbb Z_p$-algebra map commutes with the action.
If $c$ is a unit, $\gamma_c$ is an automorphism with inverse
$\gamma_{c^{-1}}$.

\begin{proof}
For fixed $j$, the polynomial $c(c-1)\cdots(c-j+1)/j!$ is
continuous as a $\mathbb Q_p$-valued function. On natural integers
its values are integers. Naturals are dense in $\mathbb Z_p$ and
$\mathbb Z_p$ is closed, so these values lie in $\mathbb Z_p$
for every $c\in\mathbb Z_p$. This justifies (WPR.5).
The series $\beta_c$ has zero constant term, so every coefficient
of substitution is a finite sum. Substitution is therefore an
honest algebra endomorphism over an arbitrary $D$.

In $\mathbb Z_p[[T]]$, each coefficient of
$B_{c+d}=B_cB_d$ and $\beta_c(\beta_d(T))=\beta_{cd}(T)$ is a
continuous function of $c,d$. Both identities hold for natural
integers by the ordinary binomial theorem, hence for all $c,d$
by density. Transport them coefficientwise to $D$. They give
composition, the identity and all the displayed formulas. The
last formula uses the commutation of $\gamma_c$ with $\gamma_k$.
Constant terms are preserved because substitution has zero
constant term. Every coefficient expression is a finite algebraic
expression in the original coefficients and the images of
$\binom cj$, so it commutes with the indicated base change.
Composition with $c^{-1}$ proves the unit assertion.
\end{proof}
% END WPR3

% BEGIN WPR4
\paragraph{WPR4: formal-series homomorphisms and their differences.}
\label{claim:WPR4}
Let $f,g:A[[T]]\to D[[T]]$ be ring maps agreeing on constants
and on $T$, with $f(T)\in TD[[T]]$. Then $f=g$.
In particular an $A$-algebra endomorphism $f$ with $f(T)\in TA[[T]]$
preserves constant coefficients and satisfies
\[
                         f(T)-T\mid f(h)-h
                  \qquad(h\in A[[T]]).
 \tag{WPR.7}
\]
No continuity hypothesis is added to these original algebraic statements.

\begin{proof}
Write $h=h_{<N}+T^Nh_N$, with $h_{<N}$ the degree-$<N$
polynomial truncation. The two maps agree on that polynomial,
so $f(h)-g(h)$ is divisible by $f(T)^N$, hence by $T^N$.
This holds for every $N$, and coefficientwise separation gives
equality. For an $A$-algebra endomorphism, comparison with
substitution at $y=f(T)$ proves $f(h)=h(y)$. The constant term
statement follows at once. Explicitly, a quotient for (WPR.7) is
\[
 \sum_{n\geq1}h_n\sum_{j=0}^{n-1}y^jT^{n-1-j},
              \qquad h=\sum_{n\geq0}h_nT^n.
\]
Its degree-$m$ coefficient has contributions only from $n\leq m+1$,
because $y\in(T)$. It is a well-defined formal series, and the
finite geometric-sum identity proves the claimed product equality
coefficient by coefficient.
\end{proof}
% END WPR4

% BEGIN WPR5
\paragraph{WPR5: uniqueness of the commuting cyclotomic action.}
\label{claim:WPR5}
Let $D$ be a domain with $p\neq0$ and $p$ not a unit.
Suppose $f,g$ are $D$-algebra endomorphisms of $D[[T]]$,
their images of $T$ have zero constant term and the same linear
coefficient, and both commute with the endomorphism
$\phi_0(T)=(1+T)^p-1$. Then $f=g$.
Consequently the full original \texttt{act\_\allowbreak unique} holds:
on its complete DVR, the prescribed Frobenius is $\phi_0$,
and an endomorphism commuting with it, of zero constant term and
linear coefficient $c$, is $\gamma_c$.

\begin{proof}
Put $F=f(T)$ and $G=g(T)$. If $F-G\neq0$, let $k\geq2$ be
the first degree where its coefficient $b$ is nonzero.
Commutation gives
\[
 (F-G)((1+T)^p-1)=(1+F)^p-(1+G)^p.
\]
The degree-$k$ coefficient on the left is $p^k b$; on the right
it is $p b$. Every higher power of $F$ or $G$ contributes to
their difference only in degrees at least $k+1$. Thus
$p(p^{k-1}-1)b=0$. None of these factors is zero: if
$p^{k-1}=1$, then $p$ would be a unit. This contradicts the
domain hypothesis. Therefore $F=G$, and WPR4 gives $f=g$.

In the original uniqueness statement, WPR4 first identifies the
given $\phi$ from its value on $T$. WPR3 shows that $\gamma_c$
has the required linear coefficient and commutes with that $\phi$.
Apply the preceding argument. This proves the full original
complete-DVR specialization and, in fact, the broader uniqueness
statement just printed.
\end{proof}
% END WPR5

% BEGIN WPR6
\paragraph{WPR6: identification of the supplied series operators.}
\label{claim:WPR6}
For the original $R:\mathrm{WSer}(\Gamma)$ and the character
of WPR2,
\[
 R.\varphi_0=\phi_0,
 \qquad R.\gamma_0(\mathrm{res}_\Gamma g)=\gamma_{\chi(g)}
                         \quad(g\in G_K).
 \tag{WPR.8}
\]
These are equalities on the actual coefficient ring $O[[T]]$.

\begin{proof}
The source fields $R.\varphi_0(T)=Tq_0$ and the defining
formula for $q_0$ give $R.\varphi_0(T)=(1+T)^p-1$;
the prime of $P.\mathrm{PD}(K)$ is the original $p$.
The supplied $R.\gamma_0(\mathrm{res}_\Gamma g)$ commutes with
this Frobenius, sends $T$ to a series with zero constant term,
and has linear coefficient $\chi_\Gamma(\mathrm{res}_\Gamma g)$.
The original character compatibility identifies this coefficient
with the image of $P.\chi_Q(g)$, hence of $\chi(g)$ in $O$.
WIS1 gives $p\neq0$ and $p$ nonunit in this complete DVR.
WPR5 proves both equalities. This argument quantifies over every
original $g\in G_K$ and does not replace the supplied $R$.
\end{proof}
% END WPR6

% BEGIN WPR7
\paragraph{WPR7: all cyclotomic parameter and ideal-generator formulas.}
\label{claim:WPR7}
For every commutative $\mathbb Z_p$-algebra $D$ and $s\geq1$, put
\[
 U_s=(1+T)^{p^{s-1}}-1,\quad
 Q_s=\sum_{j=0}^{p-1}(1+U_s)^j,\quad d_s=p^{s-1}(p-1).
 \tag{WPR.9}
\]
Then the full original \texttt{XX\_\allowbreak QQ\_\allowbreak properties} holds:
\[
 \begin{gathered}
 \phi_0(U_s)=U_sQ_s=(1+T)^{p^s}-1,
 \qquad \phi_0^{s-1}(T)=U_s,\quad Q_s(0)=p,\\
 Q_s=T^{d_s}+p b,\qquad Q_s=p+U_s b'
                    \quad\text{for some }b,b'\in D[[T]],\\
 \gamma_c(U_s)=U_s(c+U_s b_c),\quad
 \gamma_c(Q_s)=Q_s(1+U_s b'_c)\quad(c\in\mathbb Z_p^\times),\\
 u=1+p^s w\quad\Longrightarrow\quad
         \gamma_u(f)-f\in U_sQ_sD[[T]]\quad(f\in D[[T]]).
 \end{gathered}                                                     \tag{WPR.10}
\]

\begin{proof}
The geometric-sum identity gives the first product, and iteration
of $\phi_0(T)=(1+T)^p-1$ gives the second. At $T=0$ each of the
$p$ summands is one. The polynomial
$\sum_{j<p}(1+Z)^j$ is congruent to $Z^{p-1}$ modulo $p$:
multiply by $Z$ and use $(1+Z)^p-1\equiv Z^p$ in
$\mathbb F_p[Z]$, then cancel $Z$ in that polynomial ring.
Lifting the integral coefficient identity to any $D$, and using
$U_s\equiv T^{p^{s-1}}\pmod p$, gives
$Q_s=T^{d_s}+pb$. Subtracting the value at $U_s=0$ gives
$Q_s=p+U_sb'$ directly.

The zero and linear coefficients in WPR3 give
$\beta_c(T)=T(c+Tb_c)$. Commuting with $\gamma_{p^{s-1}}$
gives the formula for $U_s$. For $Q_1$ write $b=b_c$ and
$D_c=c+Tb$. It is a unit since $c$ is. Commutation applied to
$\phi_0(T)=TQ_1$ gives, after cancellation of the regular
formal variable $T$,
\[
 D_c\gamma_c(Q_1)=Q_1\{c+TQ_1\phi_0(b)\}.
\]
Thus
$\gamma_c(Q_1)=Q_1[1+TD_c^{-1}(Q_1\phi_0(b)-b)]$.
Applying $\gamma_{p^{s-1}}$ proves the asserted formula for
$Q_s$. Cancellation here is by $T$ in a power-series ring,
which is valid over every coefficient ring $D$.

Finally set $V=U_sQ_s=(1+T)^{p^s}-1$. The additive and
multiplicative binomial identities in WPR3 show
\[
 \gamma_{1+p^s w}(T)-T=(1+T)\beta_w(V).
\]
The right side is divisible by $V$, because $\beta_w$ has
zero constant term. WPR4 then proves the last assertion for
every formal series $f$. All arguments commute with the given
coefficient map, including rings with $p$-torsion.
\end{proof}
% END WPR7

\subsection{Division, completeness and cyclotomic specialization}

% BEGIN WPR8
\paragraph{WPR8: full distinguished division and the original kernel criterion.}
\label{claim:WPR8}
Let $D$ be a complete separated local ring for its maximal ideal,
let $p\in\mathfrak m_D$, and let
$Q=T^d+pb\in D[[T]]$. Every $h\in D[[T]]$ has a unique
expression
\[
                     h=Qq+r,\qquad r\in D[T],\quad\deg r<d.
 \tag{WPR.11}
\]
For $d=0$ this means $r=0$. If a ring map $f:D[[T]]\to C$
satisfies $f(Q)=0$ and is injective on polynomials of degree
less than $d$, then
\[
                              \ker f=(Q).
 \tag{WPR.12}
\]
This proves the full original \texttt{ker\_\allowbreak eq\_\allowbreak span}, with no
Noetherian or continuity assumption added.

\begin{proof}
For $d=0$, $Q=1+pb$ has unit constant coefficient, so it is a
unit in $D[[T]]$ and the statement follows. Suppose $d>0$.
Write $H_d(h)$ for the tail after removing its first $d$
coefficients, so $h=r_0+T^dH_d(h)$ with $\deg r_0<d$.
The required quotient must solve
\[
       q+pL(q)=H_d(h),\qquad L(q)=H_d(bq).
\]
It is given coefficientwise by
\[
                       q=\sum_{j\geq0}(-p)^jL^j(H_d(h)).
\]
For each coefficient, the $j$th term lies in
$\mathfrak m_D^j$, so completeness defines the sum.
Each coefficient of $L(q)$ uses only finitely many coefficients
of $q$ and $b$; hence $L$ commutes with these limits.
Telescoping gives $q+pL(q)=H_d(h)$. Therefore the tail of
$h-Qq$ vanishes, yielding the required polynomial $r$.

If $q+pL(q)=0$, iteration gives
$q=(-p)^NL^N(q)$ for every $N$. Each coefficient belongs to
every $\mathfrak m_D^N$, so separation makes $q=0$.
The quotient and then the remainder are unique. This proof uses
the given maximal-ideal completeness, not a tacit assumption that
every principal ideal is closed in that topology.

Apply $f$ to (WPR.11). If $f(h)=0$, then $f(r)=0$ because
$f(Q)=0$. The stated degree-$<d$ injectivity makes $r=0$,
so $h\in(Q)$. The reverse inclusion is immediate. This also
handles the possible zero polynomial without a natural-degree
convention causing an extra case.
\end{proof}
% END WPR8

% BEGIN WPR9
\paragraph{WPR9: completeness for the actual prism ideal.}
\label{claim:WPR9}
Let $D$ be complete and separated for $(p)$, and suppose
$Q=T^d+pb$ with $d>0$. Then $D[[T]]$ is complete and
separated for $(p,Q)$. This is the original
\texttt{is\allowbreak Adic\allowbreak Complete\_\allowbreak span} at its full coefficient-ring scope.

\begin{proof}
First put $J=(p,T)$. A series belongs to $J^N$ exactly when
its coefficient of $T^i$, for $i<N$, belongs to $p^{N-i}D$.
One implication follows from the generators $p^{N-i}T^i$;
for the other, choose the indicated finitely many quotients and
leave the degree-$\geq N$ tail in $T^ND[[T]]$.

A compatible system modulo $J^N$ therefore determines, at every
fixed degree, a compatible system of coefficients modulo all
powers of $p$. Completeness of $D$ supplies a unique coefficient.
The resulting formal series has every prescribed class by the
same membership criterion. Separation follows coefficientwise.
Thus $D[[T]]$ is $J$-adically complete and separated.

Set $I=(p,Q)$. From $T^d=Q-pb$ and $d>0$ we obtain
$J^d\subset I\subset J$. Consequently
$J^{dN}\subset I^N\subset J^N$ for all $N$. These explicit
containments identify the Cauchy systems and separated limits
for the two topologies. They transfer the completeness just
proved to $I$, without finite-generation assumptions on $D$.
\end{proof}
% END WPR9

% BEGIN WPR10
\paragraph{WPR10: the full integral cyclotomic polynomial statement.}
\label{claim:WPR10}
Let $s\geq1$, let $\zeta$ be any primitive $p^s$-th root in
the original algebraic closure, and let
$t\in\mathcal O_C$ be its specified image $\zeta-1$.
Then
\[
 \deg r<d_s,\quad r\in W_0[T],\quad r(t)=0
                       \Longrightarrow r=0.                  \tag{WPR.13}
\]
If $K\subset K_s$, every $x\in\mathcal O_K$ has
$x=r(t)$ in $\mathcal O_C$ for some $r\in W_0[T]$.
All embeddings are the original $\mathrm{intToOC}$ maps.

\begin{proof}
The root-norm calculation in LS2, used in IFD9, applies to any
primitive root and gives $|\zeta-1|=p^{-1/d_s}$. For
$a_i\in K_0$, the nonzero terms $a_i(\zeta-1)^i$,
$0\leq i<d_s$, have distinct valuations modulo $\mathbb Z$,
because the valuation of each $a_i$ is an integer. Their sum
cannot vanish. This proves (WPR.13).

These $d_s$ powers are a basis of $K_s$ over $K_0$: the preceding
independence gives the lower degree bound, and the cyclotomic
polynomial gives the upper bound. Write an integral $x\in K_s$
in this basis. The norm of the sum is its largest term norm.
Since the exponents $i/d_s$ lie in $[0,1)$ and the coefficient
valuations are integral, every coefficient is integral.
Thus $\mathcal O_{K_s}=W_0[\zeta-1]$. Apply this to the given
$x\in\mathcal O_K\subset\mathcal O_{K_s}$. Polynomial evaluation
commutes with the specified inclusions into the algebraic closure
and then its completion; hence the resulting equality is exactly
the original one in $\mathcal O_C$. The case $p=2,s=1$, where
$d_s=1$, is included.
\end{proof}
% END WPR10

\subsection{Descending the prescribed tensor-valued inclusion}

% BEGIN WPR11
\paragraph{WPR11: a normalized linear retraction.}
\label{claim:WPR11}
Let $Z$ be a local ring with maximal ideal $(p)$, let $O'$ be
a free $Z$-algebra, and suppose the image of $p$ is not a unit
in $O'$. There is a $Z$-linear map $r:O'\to Z$ with $r(1)=1$.
For every $Z$-algebra $A$, the map
\[
 j:A\longrightarrow O'\otimes_Z A,\qquad x\longmapsto1\otimes x
 \tag{WPR.14}
\]
is injective. The induced linear map
$\rho=r\otimes1:O'\otimes_Z A\to A$ has
$\rho j=1$ and $\rho(j(a)y)=a\rho(y)$.

\begin{proof}
Choose a free basis of $O'$ over $Z$. At least one coordinate
$c$ of $1$ is a unit. Otherwise all its finitely many nonzero
coordinates belong to $(p)$, so $1=pb$ in $O'$ for some $b$,
contradicting the nonunit hypothesis. Multiply the corresponding
coordinate functional by $c^{-1}$ to obtain $r$.
The balanced tensor map $\rho$ satisfies
$\rho(o\otimes a)=r(o)a$, giving $\rho(1\otimes a)=a$.
This proves injectivity of $j$. On a pure tensor,
$\rho(j(a)(o\otimes b))=r(o)ab=a\rho(o\otimes b)$;
additivity gives the formula for every tensor. No finiteness
of the free basis was needed in the retraction argument.
\end{proof}
% END WPR11

% BEGIN WPR12
\paragraph{WPR12: full formal-series factorization through a pure tensor.}
\label{claim:WPR12}
Let $Z,O',A$ be commutative rings with the indicated $Z$-algebra
structures. Suppose a $Z$-linear $r:O'\to Z$ satisfies $r(1)=1$.
Let $D$ be any commutative ring and
$F:D[[T]]\to O'\otimes_Z A$ a ring map such that
\[
 F(T)=j(\pi),\qquad F(D)\subset j(A).
\]
Assume only that $O'\otimes_Z A$ is separated for the ideal
$(j(\pi))$. Then there is a unique ring map
$\iota:D[[T]]\to A$ with $F=j\iota$.
This proves the full original \texttt{exists\_\allowbreak factor} conclusion.

\begin{proof}
Use $\rho$ and the injectivity of $j$ from the tensor calculation
in WPR11; that calculation uses only the supplied $r(1)=1$.
Every polynomial has image in $j(A)$, since its constants and
variable do. For any $h$ and any $N$, write
$h=h_{<N}+T^Nh_N$. Then
\[
 F(h)=j(a_N)+j(\pi)^NF(h_N)
\]
for some $a_N\in A$. The identity
$\rho(j(a)y)=a\rho(y)$ gives
\[
 j\rho(F(h))-F(h)
   =j(\pi)^N\{j\rho(F(h_N))-F(h_N)\}.
\]
The left side belongs to every power of $(j(\pi))$ and is
therefore zero. Thus $F(h)\in j(A)$ for every $h$.
Define $\iota(h)$ to be its unique preimage. Injectivity of $j$
and the ring identities for $F$ prove all ring identities for
$\iota$, and also uniqueness. This argument uses formal
truncations and separatedness; the map $F$ need not be supplied
with any topological structure or continuity property.
\end{proof}
% END WPR12

% BEGIN WPR13
\paragraph{WPR13: separation for every permitted Witt parameter.}
\label{claim:WPR13}
Let $e\in\mathcal O_C^\flat$ with $e-1$ a nonunit and put
$\mu=[e]-1\in A_{\mathrm{inf}}$. Then
\[
 \bigcap_{N\geq0}\mu^NA_{\mathrm{inf}}=0,
 \qquad
 \bigcap_{N\geq0}(1\otimes\mu)^N(O\otimes_{\mathbb Z_p}A_{\mathrm{inf}})=0.
 \tag{WPR.15}
\]
In particular this applies to the original $e=\epsilon_s$.

\begin{proof}
If $e=1$, both intersections are zero immediately. Otherwise
$\tau=e-1$ is a nonzero element of norm less than one in the
tilt valuation ring. Hence $\bigcap_N\tau^N\mathcal O_C^\flat=0$.
The image of $\mu$ modulo $p$ is $\tau$.

Every ideal $(\mu^N)$ is $p$-saturated. Indeed, if
$p x=\mu^N y$, reduction modulo $p$ gives $\tau^Ny_0=0$.
The tilt is a domain, so $y_0=0$ and WIS2 gives $y=pz$.
Cancel $p$ in $A_{\mathrm{inf}}$ to obtain $x=\mu^Nz$.
If $x$ belongs to every $(\mu^N)$, its zeroth Witt coordinate
belongs to every $(\tau^N)$ and is zero. Write $x=py$.
Saturation puts $y$ in the same intersection. Repeating gives
$x\in\bigcap_m p^mA_{\mathrm{inf}}=0$.

Choose a finite $\mathbb Z_p$-basis of $O$, available from WIS1.
It identifies $O\otimes A_{\mathrm{inf}}$ with finitely many copies of
$A_{\mathrm{inf}}$ as an $A_{\mathrm{inf}}$-module. Multiplication by $1\otimes\mu$
is coordinatewise multiplication by $\mu$, proving the second
intersection assertion. Finally WIS3 proves that
$\epsilon_s-1$ is a nonunit for every original $s$; alternatively,
its coordinates are nilpotent modulo $p$ by the cyclotomic
power identities. Thus the original parameter is included.
\end{proof}
% END WPR13

% BEGIN WPR14
\paragraph{WPR14: descent of every unramified constant with its theta value.}
\label{claim:WPR14}
For every $w\in W_0$ there is $b_w\in A_{\mathrm{inf}}$ such that
\[
 I.\mathrm{io}(a(Cw))=1\otimes b_w,
             \qquad \theta(b_w)=\mathrm{intToOC}_{K_0}(w).
 \tag{WPR.16}
\]
Here $Cw$ is the constant formal series, and the supplied
$\mathrm{IoW}(I)$ is the only root-compatibility premise used.

\begin{proof}
For $z\in\mathbb Z_p$, every $e_i$ fixes $z$ and
$I.\mathrm{io}$ is an $O$-algebra map. The tensor balancing
identity therefore gives
$I.\mathrm{io}(a(Cz))=1\otimes z$.
The actual theta map agrees on $\mathbb Z_p$ with its usual
embedding into $\mathcal O_C$. This can be checked without a
continuity assumption on a ring map $\psi:\mathbb Z_p\to\mathcal O_C$:
write $z=b_N+p^Ny_N$ with an integer $b_N$ and
$y_N\in\mathbb Z_p$. The images under $\psi$ and the usual map
differ by $p^N$ times an integral element, so their difference
has norm at most $p^{-N}$ for every $N$. They are equal.

Now let $\omega\in W_0$ be a root of unity of order $m$
prime to $p$. Choose the compatible $p$-power roots of $\omega$
inside its finite cyclic group; explicitly their exponents are
$p^{-r}\bmod m$. Their reductions form the original allowed
tilt $\omega^\flat$, with $(\omega^\flat)^\sharp=\omega$.
Apply $\mathrm{IoW}(I)$ to the actual tuple $(e_i\omega)_i$.
It gives
\[
 I.\mathrm{io}(a(C\omega))=1\otimes[\omega^\flat],
       \qquad \theta([\omega^\flat])=\omega.
\]
WPR1 expresses any $w$ as a finite sum
$\sum_i z_i\omega_i$. Both $I.\mathrm{io}\circ a\circ C$ and
$\theta$ are ring maps, and $b\mapsto1\otimes b$ is a ring
map. Taking $b_w=\sum_i z_i[\omega_i^\flat]$ proves both
equalities. The construction uses a finite linear combination;
no additivity of Teichmuller representatives is assumed.
\end{proof}
% END WPR14

% BEGIN WPR15
\paragraph{WPR15: the full original inclusion into one copy of Ainf.}
\label{claim:WPR15}
For every original $P,AI,K,LL$, every $s\geq0$, and every
supplied $\Gamma,R,I,\mathrm{IoW}(I)$, there is a ring map
$\iota:\mathfrak S\to A_{\mathrm{inf}}$ with
\[
 \begin{gathered}
 I.\mathrm{io}(a(f))=1\otimes\iota(f)\quad(f\in\mathfrak S),\\
 \iota(T)=[\epsilon_s]-1,\qquad
 \theta\iota(Cw)=\mathrm{intToOC}_{K_0}(w)\quad(w\in W_0).
 \end{gathered}                                                     \tag{WPR.17}
\]
This supplies the full original
\texttt{exists\_\allowbreak wach\allowbreak Inclusion\_\allowbreak power\allowbreak Series\_\allowbreak into\_\allowbreak one\_\allowbreak tmul\_\allowbreak ainf}
on every permitted supplied $\Gamma,R,I,\mathrm{IoW}(I)$.

\begin{proof}
Use WPR11 with $Z=\mathbb Z_p$ and $O'=O$; WIS1 gives
freeness and noninvertibility of $p$. Let $F=I.\mathrm{io}\circ a$.
Its value on $T$ is $1\otimes([\epsilon_s]-1)$ by the original
$I.\mathrm{io}_\pi$ field, because $a(T)$ is the tuple of
the original variables. WPR14 gives all its constant values in
the image of $1\otimes A_{\mathrm{inf}}$. WPR13 provides the required
separation. WPR12 now supplies $\iota$ and the first equality.

Injectivity of $b\mapsto1\otimes b$ then gives the asserted
value of $\iota(T)$ and identifies $\iota(Cw)$ with the
particular $b_w$ of WPR14. This proves its theta value as well.
No $\mathrm{Good}$, $s\geq1$ or $K\subset K_s$ assumption
is used in WPR11--WPR15: separation in WPR13 holds at every
original root index, and WPR14 uses only the supplied constant
map and $\mathrm{IoW}$. All original witnesses therefore remain
arbitrary at this broader scope. The proof even applies
to the source's broader \texttt{exists\_\allowbreak io} statement with any
permitted $e$ and coefficient-tuple map satisfying its printed
constant and variable hypotheses.
\end{proof}
% END WPR15

% BEGIN WPR16
\paragraph{WPR16: the actual Frobenius and Galois intertwining maps.}
\label{claim:WPR16}
Define the ring endomorphism
\[
 \varphi_{\mathfrak S}(f)
       =\sigma(f)((1+T)^p-1)
\]
by applying $\sigma$ coefficientwise before substitution. Then
\[
 \begin{gathered}
 a\varphi_{\mathfrak S}=\mathrm{let}\varphi_{R.\varphi_0}\,a,
 \qquad a\gamma_{\chi(g)}
       =\mathrm{let}\gamma_{R.\gamma_0(\mathrm{res}_\Gamma g)}\,a,\\
 \iota\varphi_{\mathfrak S}=\varphi_{A_{\mathrm{inf}}}\iota,
 \qquad \iota\gamma_{\chi(g)}=g\iota
                                \quad(g\in G_K).
 \end{gathered}                                                     \tag{WPR.18}
\]
The coefficient map $a$ is injective.

\begin{proof}
The source embeddings satisfy $e_i\sigma=e_{i-1}$, with the
original cyclic index convention. Applying WPR3's base-change
formula and then WPR6 gives the first identity componentwise:
coefficient Frobenius shifts precisely to component $i-1$,
and the variable substitution is the prescribed $R.\varphi_0$.
For the second identity no shift occurs; each $e_i$ commutes
with $\gamma_{\chi(g)}$, and WPR6 identifies its source operator.
The index set is nonempty since $[K_0:\mathbb Q_p]>0$.
Evaluation at one component and injectivity of its field embedding
show that $a$ is injective coefficient by coefficient.

To prove the last two identities, first apply
$b\mapsto1\otimes b$ and then the actual inclusion
$O\otimes A_{\mathrm{inf}}\hookrightarrow O\otimes W(\operatorname{Frac}
\mathcal O_C^\flat)$. This composite is injective by WPR11
and WIS5. Substitute (WPR.17) and the first two identities in
(WPR.18). The defining $I.\mathrm{ioE}$ Frobenius and Galois
compatibilities, together with its extension of $I.\mathrm{io}$,
make the two resulting expressions equal in that tensor ring.
On a pure tensor $1\otimes b$, the tensor Frobenius and Galois
actions give $1\otimes\varphi(b)$ and $1\otimes g(b)$.
These equalities follow directly from the Witt coordinate maps.

The Galois action specified by $AI$ is the same action used here:
its required compatibility with every inverse-limit coordinate
determines the tilt action uniquely, and Witt functoriality then
determines its action on $A_{\mathrm{inf}}$. This is also the coordinate
equivariance checked in IFD8. Injectivity of the composite map
now proves the last two identities in $A_{\mathrm{inf}}$ itself.
\end{proof}
% END WPR16

% BEGIN WPR17
\paragraph{WPR17: the exact specialization kernel and its base structure.}
\label{claim:WPR17}
Put $f=\theta\iota:\mathfrak S\to\mathcal O_C$.
Then
\[
 \ker f=(Q_s),\qquad
 \mathrm{intToOC}_K(\mathcal O_K)\subset f(\mathfrak S).
 \tag{WPR.19}
\]
There is a ring map
$\mathrm{str}:\mathcal O_K\to\mathfrak S/(Q_s)$ such that,
for every $x\in\mathcal O_K$ and $b\in\mathfrak S$,
\[
 [b]=\mathrm{str}(x)
       \quad\Longleftrightarrow\quad f(b)=\mathrm{intToOC}_K(x).
 \tag{WPR.20}
\]

\begin{proof}
By WPR15 and the exact sharp-root values of the original
$\epsilon_s$, $f(T)=\zeta_{p^s}-1$ and
$f(Cw)=\mathrm{intToOC}_{K_0}(w)$. Thus polynomial evaluation
under $f$ is the prescribed evaluation in WPR10. Moreover,
\[
 f(Q_s)=\sum_{j=0}^{p-1}\zeta_p^j=0.
\]
Indeed multiplying this sum by $\zeta_p-1\neq0$ gives zero,
and $\mathcal O_C$ is a domain. WPR7 gives
$Q_s=T^{d_s}+pb$, while WPR10 gives the required injectivity
on polynomials of degree less than $d_s$. WPR8 therefore gives
the kernel equality. The second conclusion in WPR10, using the
actual $K\subset K_s$, gives the image assertion in (WPR.19).

Here is the general algebra behind (WPR.20). For ring maps
$f:S\to C$, $g:D\to C$, with $\ker f=J$ and
$g(D)\subset f(S)$, the induced map $\bar f:S/J\to C$
is injective and its image contains $g(D)$. Define
$\mathrm{str}(x)$ to be the unique preimage of $g(x)$.
Injectivity proves preservation of addition, multiplication and
one, and gives $[b]=\mathrm{str}(x)$ exactly when $f(b)=g(x)$.
Apply this to the present maps. This is the full original
\texttt{exists\_\allowbreak str} argument, including its exact equivalence.
\end{proof}
% END WPR17

\subsection{All fields of the three actual prisms and their maps}

% BEGIN WPR18
\paragraph{WPR18: delta structures from the actual Frobenius lifts.}
\label{claim:WPR18}
Let $A$ be a $p$-torsion-free commutative ring and let
$\varphi:A\to A$ be a ring endomorphism with
$\varphi(x)\equiv x^p\pmod{pA}$. There is a unique function
$\delta$ with $\varphi(x)=x^p+p\delta(x)$, and it satisfies
\[
 \begin{split}
 \delta(1)&=0,\\
 \delta(xy)&=x^p\delta(y)+y^p\delta(x)+p\delta(x)\delta(y),\\
 \delta(x+y)&=\delta(x)+\delta(y)
       -\sum_{j=1}^{p-1}\frac{\binom pj}{p}x^jy^{p-j}.
 \end{split}                                                       \tag{WPR.21}
\]
A ring map between two such rings that commutes with Frobenius
also commutes with $\delta$ whenever the target is $p$-torsion-free.
The maps $\sigma$ on $W_0$, $\varphi_{\mathfrak S}$ on
$\mathfrak S$ and Witt Frobenius on $A_{\mathrm{inf}}$ have these structures.

\begin{proof}
Divisibility gives existence of $\delta(x)$ and cancellation of
$p$ gives uniqueness. Expand the ring identities for $\varphi$
and cancel $p$ to obtain (WPR.21); the middle binomial
coefficients are divisible by the prime $p$. Applying a ring
map to $\varphi(x)=x^p+p\delta(x)$ and comparing with the
same identity in the target proves functoriality by cancellation.

WPR1 gives the Frobenius congruence for $\sigma$. For any
commutative $\mathbb Z_p$-algebra $D$ with an endomorphism
$\sigma$ lifting coefficient Frobenius, its associated series
map is congruent modulo $p$ to
$f\mapsto\sum a_i^pT^{pi}=f^p$. Substitution is coefficientwise
finite, and $(1+T)^p-1\equiv T^p$. Thus every coefficient of
$\varphi_{D[[T]]}(f)-f^p$ is divisible by $p$, and choosing
their quotients gives a formal series quotient. This proves the
full original \texttt{phi\allowbreak S\_\allowbreak lift}, including coefficients with
$p$-torsion. For the actual $W_0$ and $\mathfrak S$,
$p$-torsion-freeness follows from the domain property.
For $A_{\mathrm{inf}}$, WIS2 proves $p$-torsion-freeness and identifies
reduction modulo $p$ with the tilt; Witt Frobenius becomes its
$p$th-power map. The constructions therefore apply to all three
actual rings.
\end{proof}
% END WPR18

% BEGIN WPR19
\paragraph{WPR19: the actual Fontaine prism over the original integers.}
\label{claim:WPR19}
The original $A_{\mathrm{inf}}$, its Witt-Frobenius delta structure,
ideal $\ker\theta$, and structure map induced by
$\mathcal O_K\hookrightarrow\mathcal O_C$ form the source's
bounded prism $\mathrm{ainfObj}(K)$. Its Galois endomorphism
for $g\in G_K$ is the original action on Witt coordinates.

\begin{proof}
Use the cyclotomic generator $\xi$ of LS4. That proof gives
$\ker\theta=(\xi)$ and shows that $\xi$ is a nonzerodivisor.
Multiplication by $\xi$ identifies its ideal with the free
rank-one module $A_{\mathrm{inf}}$, proving the required invertibility
of the ideal as a module. WPR18 supplies the delta structure.
The same cyclotomic formula gives
\[
 \theta\varphi(\xi)=p,
 \qquad p=\{p-\varphi(\xi)\}+\varphi(\xi)
                  \in(\xi)+\varphi((\xi)).                 \tag{WPR.22}
\]
The quotient is $\mathcal O_C$ by the actual surjective theta
map. It is $p$-torsion-free, so its bounded-torsion exponent may
be taken to be zero.

We verify the mixed completeness field. Write $S=\mathcal O_C^\flat$
and $\alpha=\xi\bmod p=(\epsilon_1-1)^{p-1}$.
WIS22 shows that $S$ is complete and separated for
$(\epsilon_1-1)$ and hence for $(\alpha)$, since a positive
power gives cofinal ideal powers. LS4 proves that every
$(\xi^m)$ in $A_{\mathrm{inf}}$ is $p$-saturated. Iterating saturation
gives
\[
             (\xi^m)\cap p^NA_{\mathrm{inf}}=p^N(\xi^m).             \tag{WPR.23}
\]
We claim $A_{\mathrm{inf}}/p^N$ is $\xi$-adically complete and separated
for every $N$. For $N=1$ this is the assertion about $S$.
Induct using the exact sequence with left term
$p^{N-1}A_{\mathrm{inf}}/p^NA_{\mathrm{inf}}\simeq S$ and right term
$A_{\mathrm{inf}}/p^{N-1}$. Equation (WPR.23) makes its reductions
modulo $\xi^m$ exact, with left term $S/\alpha^mS$.
Those left-term transition maps are surjective. Compatible
middle lifts may therefore be chosen successively: the defect
of a lift at stage $m+1$ is a class in $S/\alpha^mS$, which
lifts from $S/\alpha^{m+1}S$ and can be subtracted. The same
exact sequence and separatedness of its two outer terms show
uniqueness. This proves the induction assertion.

Now a compatible rectangular system modulo $(p^N,\xi^m)$
first gives an element of $A_{\mathrm{inf}}/p^N$ for each $N$, and
then an element of $A_{\mathrm{inf}}$ by its $p$-adic completeness.
All maps are the original quotient maps, so this element has
the specified rectangular classes and is unique. Finally
\[
 (p,\xi)^{2N}\subset(p^N,\xi^N)\subset(p,\xi)^N
\]
shows that the diagonal rectangular system is cofinal with
the required mixed-adic system. This proves the prism's
completeness and separation fields.

The structure map is the original inclusion into $\mathcal O_C$
followed by the inverse of the theta quotient isomorphism.
Galois acts on the tilt coordinates and hence on Witt vectors,
commutes with Frobenius and with theta, and preserves its kernel.
It therefore preserves $\delta$ by WPR18. Since $g\in G_K$
fixes the original $\mathcal O_K$, it also preserves the
structure map. These are exactly the fields of
$\mathrm{ainfHom}(K,g)$ on the actual $\mathrm{ainfObj}(K)$.
\end{proof}
% END WPR19

% BEGIN WPR20
\paragraph{WPR20: construction of the original series prism.}
\label{claim:WPR20}
The following data form the required bounded prism over
$\mathcal O_K$:
\[
 P_{\mathfrak S}.A=\mathfrak S,
 \qquad P_{\mathfrak S}.I=(Q_s),
 \qquad P_{\mathfrak S}.D.\varphi=\varphi_{\mathfrak S},
 \qquad P_{\mathfrak S}.\mathrm{str}=\mathrm{str}
                     \text{ of WPR17}.                       \tag{WPR.24}
\]
Its ring is local and a domain. Its quotient has no nonzero
$p$-power torsion.

\begin{proof}
The least nonzero coefficient of a product proves the domain
property of $W_0[[T]]$. A formal series is a unit exactly when
its constant coefficient is a unit in $W_0$; construct the
inverse recursively from that coefficient. Its unique maximal
ideal is therefore $(p,T)$, proving locality. WPR18 gives the
delta structure.

The constant coefficient of $Q_s$ is $p\neq0$, so $Q_s$ is
nonzero and regular. Its ideal is free of rank one via
multiplication by $Q_s$, and hence is an invertible module.
WPR1 gives $p$-adic completeness of $W_0$; WPR7 and WPR9
give completeness for $(p,Q_s)$ in $\mathfrak S$.

Write $Q_s=p+U_s b$ using WPR7. Since $\sigma$ fixes its
integral binomial coefficients,
$\varphi_{\mathfrak S}(U_s)=U_sQ_s$. Consequently
\[
 \varphi_{\mathfrak S}(Q_s)
     =p+U_sQ_s\varphi_{\mathfrak S}(b),\qquad
 p=\varphi_{\mathfrak S}(Q_s)
          -Q_sU_s\varphi_{\mathfrak S}(b).
 \tag{WPR.25}
\]
This proves $p\in(Q_s)+\varphi_{\mathfrak S}((Q_s))$.
By WPR17 the quotient embeds into $\mathcal O_C$ under $f$,
so it has no nonzero $p$-power torsion; exponent zero suffices
for the boundedness field. The map in (WPR.20) supplies its
base structure. These are all the fields of the original
$\mathrm{prismS}$, on the specified ring and parameter.
\end{proof}
% END WPR20

% BEGIN WPR21
\paragraph{WPR21: the constant prism and the original evaluation morphism.}
\label{claim:WPR21}
There is a bounded prism $P_W$ on $W_0$, with ideal $(p)$ and
Frobenius $\sigma$, and a prism morphism
\[
 \mathrm{ev}:P_{\mathfrak S}\longrightarrow P_W,
 \qquad \mathrm{ev}(h)=h(0).
 \tag{WPR.26}
\]
This ring map is surjective and has kernel $(T)$.

\begin{proof}
WPR1 and WPR18 supply the complete ring and delta structure.
The nonzero element $p$ is regular in $W_0$, so its ideal is
free of rank one. Completeness for $(p)+(p)$ is exactly
$p$-adic completeness. The prism containment follows already
from $p\in(p)$. Every element of $W_0/(p)$ is killed by $p$,
so the boundedness exponent is one. Give this quotient the
structure map
\[
 \mathcal O_K\xrightarrow{\mathrm{str}}
 \mathfrak S/(Q_s)\xrightarrow{\overline{\mathrm{ev}}}W_0/(p).
\]
The second arrow exists because $Q_s(0)=p$. These data are
exactly the original $\mathrm{prismW}$ with its specified
structure map.

Constant coefficients give
$\mathrm{ev}\,\varphi_{\mathfrak S}=\sigma\,\mathrm{ev}$.
Since $W_0$ has no $p$-torsion, WPR18 gives compatibility with
delta. Evaluation sends $(Q_s)$ into $(p)$, and the displayed
definition makes it compatible with the base structure. This
verifies every field of a prism morphism. The constant series
$C(w)$ maps to $w$, proving surjectivity. A series has zero
constant coefficient exactly when it is $T$ times the series
obtained by shifting all coefficients down by one. Thus the
kernel is precisely $(T)$, with no topological assumption on
the ring map.
\end{proof}
% END WPR21

% BEGIN WPR22
\paragraph{WPR22: the actual period-ring morphism of prisms.}
\label{claim:WPR22}
The ring map $\mathrm{io}$ of WPR15 defines a morphism
\[
 \mathrm{ioH}:P_{\mathfrak S}\longrightarrow\mathrm{ainfObj}(K)
 \tag{WPR.27}
\]
to the prescribed Fontaine prism, with the original pure-tensor
identity $I.\mathrm{io}(a h)=1\otimes\mathrm{io}(h)$.

\begin{proof}
The ring map is the unique factor of the original
$I.\mathrm{io}\circ a$ constructed in WPR15. WPR16 proves
that it commutes with the original Frobenius maps, and WPR18
therefore proves compatibility with delta in the $p$-torsion-free
target. WPR17 says that $f=\theta\,\mathrm{io}$ has kernel
$(Q_s)$, so $\mathrm{io}$ sends the source ideal into
$\ker\theta$, the target ideal.

For $x\in\mathcal O_K$ and $h\in\mathfrak S$, WPR17 gives
\[
 [h]=\mathrm{str}(x)
 \quad\Longleftrightarrow\quad
 \theta\,\mathrm{io}(h)=\mathrm{intToOC}(x).
\]
The same theta criterion defines the structure map of
$\mathrm{ainfObj}(K)$ in WPR19. Hence $[\mathrm{io}(h)]$
is its structure class of $x$ whenever $[h]$ is the source
structure class. This is exactly the remaining base-compatibility
field of a prism morphism. The underlying target ring is the
actual $A_{\mathrm{inf}}$; its identification $\mathrm{ainfId}$ in the
source merely records that ring, so the pure-tensor identity is
the one required there, including that identification.
\end{proof}
% END WPR22

% BEGIN WPR23
\paragraph{WPR23: all original Galois prism morphisms.}
\label{claim:WPR23}
For every $g\in G_K$ the original action
$\gamma_{\chi(g)}$ defines a prism automorphism
$gH(g):P_{\mathfrak S}\to P_{\mathfrak S}$, and
\[
 \mathrm{ioH}\circ gH(g)=\mathrm{ainfHom}(K,g)\circ\mathrm{ioH},
 \qquad \mathrm{ev}\circ gH(g)=\mathrm{ev}.
 \tag{WPR.28}
\]
These maps respect identity and multiplication in $G_K$.

\begin{proof}
The coefficient map $\sigma$ fixes $\mathbb Z_p$, so it commutes
with $\gamma_c$ for every $c\in\mathbb Z_p$. WPR3 gives
$\gamma_c\gamma_p=\gamma_p\gamma_c$. Thus
$\gamma_c\varphi_{\mathfrak S}=\varphi_{\mathfrak S}\gamma_c$;
WPR18 turns this into delta compatibility. For the unit
$c=\chi(g)$, WPR7 gives
\[
              \gamma_c Q_s=Q_s(1+U_s b_c).
\]
The second factor is a unit, because $U_s(0)=0$. The source
ideal is therefore preserved, in both directions by using
$c^{-1}$.

By WPR16 and the original theta equivariance,
$f\gamma_{\chi(g)}=\sigma_O(g)f$. The automorphism $g\in G_K$
fixes the actual embedded $\mathcal O_K$. The structure-class
criterion in WPR17 consequently proves base compatibility.
This establishes the claimed prism morphism and, using
$\gamma_{c^{-1}}$, its inverse. WPR16 proves the first equality
in (WPR.28) on underlying rings; prism morphisms with equal
underlying ring maps are equal. The second follows from
$\beta_c(0)=0$, so substitution preserves constant coefficients.
Finally WPR3 and the character law of WPR2 give identity and
composition exactly. No chosen lift of an element of $\Gamma$
enters this construction; it is defined for every original
$g\in G_K$.
\end{proof}
% END WPR23

% BEGIN WPR24
\paragraph{WPR24: the original scalar, parameter and distinguished polynomial maps.}
\label{claim:WPR24}
Put $\pi=T$, $Q=Q_s$, and let
$z:\mathbb Z_p[[T]]\to\mathfrak S$ be coefficient extension.
Then the original coefficient map $a$ satisfies
\[
 a\pi=\mathrm{let}\pi,\qquad aQ=R.Q,\qquad
 a\circ z=\mathrm{zlA},
 \tag{WPR.29}
\]
as literal equalities on the original rings. It is injective,
and it intertwines the prescribed Frobenius and Galois operators.

\begin{proof}
Injectivity and both intertwining identities are WPR16.
Each coefficient embedding sends the variable $T$ to $T$,
giving the first equality. Since every $e_i$ is
$\mathbb Q_p$-linear, its integral restriction commutes with
$\mathbb Z_p$ scalars. Coefficientwise this says precisely
$a\circ z=\mathrm{zlA}$; it is not a replacement scalar action.

For the polynomial equality, WPR6 identifies the supplied
$R.\varphi_0$ with substitution by $(1+T)^p-1$. Its defining
equation is $R.\varphi_0(T)=Tq_0$, so regularity of $T$ gives
\[
                     q_0=\sum_{j=0}^{p-1}(1+T)^j.
\]
The source's $R.Q$ is the $(s-1)$-fold full Frobenius iterate
of the constant tuple $q_0$. The tuple shift has no effect on
a constant tuple, and all the coefficient embeddings fix its
integer coefficients. WPR3 identifies its iterate with
\[
 \gamma_p^{s-1}(q_0)
   =\sum_{j=0}^{p-1}(1+T)^{j p^{s-1}}=Q_s.
\]
This proves $aQ_s=R.Q$ on every component. All cancellations
used regular $T$, not a division by $p$ or a change of layer.
\end{proof}
% END WPR24

\subsection{Both coefficient-coordinate assertions}

% BEGIN WPR25
\paragraph{WPR25: the full generic unit-separated coordinate theorem.}
\label{claim:WPR25}
Let $R_0$ be a nonzero commutative ring, $W$ a finite free
commutative $R_0$-algebra of rank $|J|$, $J$ a finite set,
and $O'$ a commutative $R_0$-algebra. Suppose maps
$\phi_k:W\to O'$, $k\in J$, are $R_0$-algebra maps and
for each $k\ne j$ there is $w\in W$ such that
$\phi_k(w)-\phi_j(w)$ is a unit. For every finite
$R_0$-basis $(o_i)_{i<d}$ of $O'$ and every $y\in(O')^J$
there is a unique $(c_i)_{i<d}\in W^d$ with
\[
                    y_k=\sum_{i<d}o_i\phi_k(c_i)
                    \quad(k\in J).                         \tag{WPR.30}
\]
This is the full original \texttt{exists\allowbreak Unique\_\allowbreak of\_\allowbreak separated}.

\begin{proof}
Consider the $O'$-algebra map
\[
 \Phi:O'\otimes_{R_0}W\longrightarrow(O')^J,
 \qquad \Phi(b\otimes w)_k=b\phi_k(w).
\]
Fix $k$. For every $j\ne k$, choose $w_{kj}$ with unit
$u_{kj}=\phi_k(w_{kj})-\phi_j(w_{kj})$. The element
\[
 t_k=\prod_{j\ne k}u_{kj}^{-1}
        \bigl(1\otimes w_{kj}-\phi_j(w_{kj})\otimes1\bigr)
\]
maps to the coordinate vector with value one at $k$ and zero
elsewhere: the $k$-coordinate of every factor is one, while
at coordinate $j$ its own factor is zero. Thus $\Phi$ is
surjective, since these vectors generate $(O')^J$ over $O'$.

Choose an $R_0$-basis of $W$ of size $|J|$. It makes source
and target of $\Phi$ free $O'$-modules of the same finite rank.
Choose preimages of the target basis to obtain a right-inverse
matrix $B$ for its matrix $M$. Then
$\det(M)\det(B)=1$, so the adjugate formula makes $M$ invertible.
This works over an arbitrary commutative $O'$; it does not
require a Noetherian hypothesis. Empty ranks satisfy the same
claim, or may be checked directly.

The given basis $(o_i)$ identifies $W^d$ with
$O'\otimes_{R_0}W$ by $(c_i)\mapsto\sum_i o_i\otimes c_i$.
Its composite with $\Phi$ is exactly the map in (WPR.30).
It is bijective, proving existence and uniqueness for every
supplied basis, with no preferred choice of that basis.
\end{proof}
% END WPR25

% BEGIN WPR26
\paragraph{WPR26: separators and the actual unramified coefficient coordinates.}
\label{claim:WPR26}
Let $K,LL$ be any original finite extensions and let
$e:\mathbb Z/n\mathbb Z\to\operatorname{Emb}_{\mathbb Q_p}(K_0,LL)$
be any injective enumeration, where $n=[K_0:\mathbb Q_p]$.
No $\mathrm{Good}$ or $\Gamma$ datum, or cyclic Frobenius
indexing, is assumed in this claim. These embeddings satisfy
the separation hypothesis of WPR25 on $W_0$. Consequently,
for every $\mathbb Z_p$-basis
$(o_i)_{i<d}$ of the original $O$ and every $y\in O^{\mathbb Z/n\mathbb Z}$,
there is a unique $c\in W_0^d$ with
\[
                    y_k=\sum_{i<d}o_i e_k(c_i).
 \tag{WPR.31}
\]
This proves the full original
\texttt{exists\allowbreak Unique\_\allowbreak unramified\allowbreak Integers\_\allowbreak coords\_\allowbreak basis}.

\begin{proof}
First, in any local domain $D$, distinct $m$-th roots of unity
$a,b$ have unit difference whenever $m>0$ is a unit in $D$.
Indeed
\[
 (a-b)\sum_{j=0}^{m-1}a^j b^{m-1-j}=a^m-b^m=0.
\]
The nonzero first factor and the domain property make the sum
zero. If $a-b$ were a nonunit, their residues would agree.
Reducing the zero sum would give
$0=m\bar b^{m-1}$, a unit in the residue field, a contradiction.
This proves the source's generic unit-difference lemma.

Two distinct $\mathbb Q_p$-embeddings of $K_0$ must differ on
one of the prime-to-$p$ roots spanning $W_0$ in WPR1. Otherwise
they agree on $W_0$ and, after multiplying any field element
by a sufficiently large power of $p$, on all of $K_0$.
For such a root $\omega$ of order $m$ prime to $p$, both
$e_k(\omega)$ and $e_j(\omega)$ are integral roots of unity
in $O$, and are distinct. The number $m$ is a unit in
$\mathbb Z_p$ and hence in $O$. WIS1 gives that $O$ is a
local domain, so the preceding argument makes their difference
a unit. This is separation on the actual coefficient maps.

The supplied enumeration $e_k$ is injective. Its cardinality
is $n=[K_0:\mathbb Q_p]$, which equals the residue degree
and the size of the $\mathbb Z_p$-basis of $W_0$ in WPR1.
Apply WPR25 with $R_0=\mathbb Z_p$, $W=W_0$, $O'=O$.
All hypotheses use the original unramified field and the
supplied injective enumeration; no compatibility with Frobenius
was used. The result holds for every $d$ and every given
basis $(o_i)$, exactly as asserted. In particular it applies
to the embeddings of the original $\Gamma$ in the prism
construction.
\end{proof}
% END WPR26

% BEGIN WPR27
\paragraph{WPR27: the full coefficientwise formal-series lifting theorem.}
\label{claim:WPR27}
Let $A,B$ be commutative rings, $(o_i)_{i<d}$ a finite family
in $B$, and $e_k:A\to B$ ring maps indexed by an arbitrary
set $J$. Suppose every $y\in B^J$ has a unique $c\in A^d$
such that $y_k=\sum_i o_i e_k(c_i)$. Then every
$x\in B[[T]]^J$ has a unique $c\in A[[T]]^d$ such that
\[
                  x_k=\sum_{i<d}C(o_i)e_k(c_i).
 \tag{WPR.32}
\]
If $B$ is an $A$-algebra with finite basis $(b_i)$, this also
gives the full original \texttt{exists\allowbreak Unique\_\allowbreak of\_\allowbreak basis}:
\[
 x=\sum_{i<d}\operatorname{map}(A\to B)(c_i)C(b_i)
 \quad\text{with unique }c_i\in A[[T]].                      \tag{WPR.33}
\]

\begin{proof}
For each degree $r\geq0$ apply the hypothesis to the tuple
$([T^r]x_k)_{k\in J}$. Let its unique coordinate tuple be
$(c_{i,r})_{i<d}$ and define $c_i=\sum_{r\geq0}c_{i,r}T^r$.
The coefficient of degree $r$ on the right of (WPR.32) is
$\sum_i o_i e_k(c_{i,r})$, exactly $[T^r]x_k$. Thus the
series agree. Conversely an equality of series gives this
coefficient equation at every degree; the coefficient-level
uniqueness determines every $c_{i,r}$ and hence every $c_i$.
Only the finite sum over $i$ is involved; $J$ need not be
finite, and no topology on $A$ or $B$ has been introduced.

For (WPR.33), take $J$ to be a singleton, $e$ the algebra map,
and $o_i=b_i$. The coefficient-level hypothesis is exactly
the basis-coordinate theorem in $B$. Commutativity puts its
product in the displayed order. This proves the two original
generic formal-series statements at their full scopes.
\end{proof}
% END WPR27

% BEGIN WPR28
\paragraph{WPR28: both finite coordinate assertions on the actual rings.}
\label{claim:WPR28}
There are $d_0$ and $w_1,\ldots,w_{d_0}\in\mathfrak S$ such
that every $h\in\mathfrak S$ has a unique tuple
$c_i\in\mathbb Z_p[[T]]$ with
\[
                         h=\sum_i z(c_i)w_i.                \tag{WPR.34}
\]
In addition, for \emph{every} $d$, every $\mathbb Z_p$-basis
$(o_i)_{i<d}$ of $O$ and every original $x\in\mathrm{LetA}$,
there is a unique $c\in\mathfrak S^d$ such that
\[
                         x_k=\sum_{i<d}C(o_i)a(c_i)_k
                         \quad\text{for every }k.           \tag{WPR.35}
\]

\begin{proof}
For (WPR.34), use the $\mathbb Z_p$-basis
$\omega_1,\ldots,\omega_n$ of $W_0$ from WPR1 and set
$w_i=C(\omega_i)$. WPR27 gives existence and uniqueness
coefficient by coefficient. Its scalar map is precisely the
$z$ of WPR24, not an abstract finite-free identification.
Thus $d_0=n$ is one valid choice in the original existential
coordinate clause.

For (WPR.35), fix any of the basis data quantified in the
statement. WPR26 gives the coefficient-level unique
decomposition using these same $o_i$ and the original $e_k$.
Apply WPR27 with $A=W_0$ and $B=O$. By definition its series
maps $e_k$ are exactly the $k$-components of $a$. The result
is (WPR.35) for all of $\mathrm{LetA}$. No choice of the
$\mathbb Z_p$-basis of $O$ has been restricted, and no series
coefficient is inverted.
\end{proof}
% END WPR28

\subsection{Return to the two full original endpoints}

% BEGIN WPR29
\paragraph{WPR29: the full original Wach-prism existence theorem.}
\label{claim:WPR29}
For every original prime, $P,AI,K,s,LL$ with
$\mathrm{Good}(K,s,LL)$, every supplied $\Gamma,R,I$ and
$\mathrm{IoW}(I)$, the original
\texttt{exists\_\allowbreak wach\_\allowbreak prism\_\allowbreak over\_\allowbreak integers}
holds, with all sixteen output fields and no extra premise.
Specifically take $P_{\mathfrak S}$, $P_W$, $\mathrm{ev}$,
$\mathrm{ioH}$, $gH$, $a$, $z$, $\pi=T$ and $Q=Q_s$ above.
Their complete output is as follows:
\begin{enumerate}
\item $a$ is injective.
\item $I.\mathrm{io}(a h)=1\otimes\mathrm{ainfId}(\mathrm{ioH}(h))$
for every $h$.
\item $a(\varphi_{P_{\mathfrak S}}h)=\mathrm{let}\varphi_R(a h)$
for every $h$.
\item $a(gH(g)h)=\mathrm{let}\gamma_{R,\Gamma.\mathrm{res}(g)}(a h)$
for every $g,h$.
\item $\mathrm{ioH}\circ gH(g)=\mathrm{ainfHom}(K,g)\circ\mathrm{ioH}$
for every $g$.
\item $\mathrm{ev}\circ gH(g)=\mathrm{ev}$ for every $g$.
\item $a\pi=\mathrm{let}\pi$.
\item $\mathrm{ev}$ is surjective.
\item $\ker\mathrm{ev}=(\pi)$.
\item $P_{\mathfrak S}.I=(Q)$.
\item $aQ=R.Q$.
\item $a\circ z=\mathrm{zlA}$.
\item $P_{\mathfrak S}.A$ is a local ring.
\item $P_{\mathfrak S}.A$ is a domain.
\item The finite unique coordinates in (WPR.34) hold over this $z$.
\item For every original coefficient basis, the unique coordinates
in (WPR.35) hold over this $a$.
\end{enumerate}

\begin{proof}
WPR20 and WPR21 give the two bounded prisms over the same
$\mathcal O_K$, including their delta structures, invertible
ideals, mixed completeness, prism containment, bounded torsion
and structure maps. WPR21--WPR23 give all three kinds of
prism morphisms on these very rings. Thus the displayed
objects have the types required by the original endpoint.

Fields 1, 3 and 4 are WPR16. Field 2 is WPR15 and WPR22.
Fields 5 and 6 are WPR23; fields 7, 11 and 12 are WPR24.
Fields 8 and 9 are WPR21, fields 10, 13 and 14 are WPR20,
and fields 15 and 16 are WPR28. Each uses the same maps
$a,z,\mathrm{io}$ defined above. The original source's
$\mathrm{Good}$ data supplied precisely $s\geq1$ and
$K\subset K_s$ where those facts were needed in WPR17.
All primes, including $p=2$, and all permitted layer and
coefficient-field choices have been retained. This proves
both the source's \texttt{Prism\allowbreak S.\allowbreak exists\_\allowbreak prism} and its
exported statement in their original scope.
\end{proof}
% END WPR29

% BEGIN WPR30
\paragraph{WPR30: the full original conditional all-layer Wach existence.}
\label{claim:WPR30}
Retain the original
\texttt{Bhatt\allowbreak Scholze\allowbreak Crystalline\allowbreak Lattices}
input and the original $\mathrm{LocTied}(p,P,AI)$ premise.
Then the full original $\mathrm{WachEx}(AI)$ holds. Equivalently,
the theorem \texttt{Wach\allowbreak Lattice.\allowbreak main} is now reconstructed
at all of its original hypotheses. There is no remaining
period-ring supplier in this implication.

\begin{proof}
Fix an arbitrary original $\mathrm{Good}(K,s,LL)$ and arbitrary
permitted $\Gamma,R,I$, $\mathrm{IoW}(I)$ and supplied \'etale
module $\mathcal D$. We supply the three inputs isolated by
WLD26 on exactly these data.

First, WPR29 supplies the actual Wach prism, its maps and
both coordinate statements, including $a\circ z=\mathrm{zlA}$.
Second, WIS18 supplies injectivity and reflection of units
for the original $I.\mathrm{ioE}$. Third, IFD25 supplies the
full inhomogeneous fixed-Frobenius matrix theorem; taking its
right-hand side to be zero gives the homogeneous rectangular
descent required by WLD26. Its fixed subgroup is the original
$H=G_K\cap\ker\chi$. In particular this is the actual matrix
input, with every row and column dimension allowed; it is not
merely a scalar substitute.

For clarity, the resulting construction works on each original
crystalline lattice $T$. The retained crystalline-realization
input gives its realizing crystal on the actual Fontaine prism.
Evaluation on $P_{\mathfrak S}$ and $P_W$ supplies a finite
free lattice with Frobenius denominators controlled by $Q_s$.
The original comparison maps put its basis in the tensor period
module. Since $H$ has trivial image in $\Gamma$ (apply the
injective $\chi_\Gamma$ to the defining character identity),
this basis satisfies the fixed-matrix descent hypothesis. The
homogeneous IFD25 case places it in the original \'etale
coefficient module. WLD25 and WLD26 then prove every original
Wach-lattice field: the prescribed spanning equality, the two
Frobenius inclusions after powers of $Q_s$, the Galois
congruence modulo $T$, preservation by every allowed lattice
morphism, and the original trivial-object normalization. The
fullness clause of the retained crystalline-realization input
is the one used for the morphism and trivial-object returns;
the tied-period premise gives its original standard-period
identification.

The result of that construction is $\mathrm{WN}(\mathcal D,a\circ z)$.
WPR24 identifies its scalar map with the original $\mathrm{zlA}$,
so it is literally the required $\mathrm{NS}(\mathcal D)$.
Every layer and every supplied \'etale module were fixed only
after the universal inputs, and were arbitrary. This gives
the full $\mathrm{WachEx}(AI)$ and proves the stated original
implication.
\end{proof}
% END WPR30

\paragraph{Completed return and the next mathematical boundary.}
The prescribed prism and the full original conditional
\texttt{Wach\allowbreak Ex} theorem are now supplied. The latter retains
exactly its original tied-period and crystalline-realization
premises. The independent \texttt{Faith\allowbreak H} and \texttt{Et\allowbreak Basis}
producers, the Herr cohomology class maps, the quantitative
de Rham comparison and the common \texttt{Concl\allowbreak F} family still
lie ahead. Those are needed for the full \texttt{HW/\allowbreak Fon/\allowbreak Cmp}
data and the local row bounds. No whole
\texttt{GLOBAL\_\allowbreak DENOMINATOR} obligation is closed
by this chapter.

\paragraph{Source correspondence.}
All sources are at the accepted commit.
\par\noindent\texttt{7cbf7a0c32fca1badf05931879a37c3f8a2b03b4}.
Exact file, declaration, range and proof-block hashes accompany
each claim. The principal sources under
\texttt{Proofs/\allowbreak Inertia\allowbreak Eigenvalues} are
\texttt{Exists\allowbreak Wach\allowbreak Prism\allowbreak Over\allowbreak Integers},
\texttt{Exists\allowbreak Wach\allowbreak Inclusion\allowbreak Power\allowbreak Series\allowbreak Into\allowbreak One\allowbreak Tmul\allowbreak Ainf},
\texttt{Exists\allowbreak Unique\allowbreak Unramified\allowbreak Integers\allowbreak Coords\allowbreak Basis},
\texttt{Cyclotomic\allowbreak Integers\allowbreak Eq\allowbreak Witt\allowbreak Zeta\allowbreak Sub\allowbreak One},
\texttt{CDvd\allowbreak Of\allowbreak Map\allowbreak Eq\allowbreak Zero},
\texttt{Geom\allowbreak Sum\allowbreak One\allowbreak Add\allowbreak Prime},
and \texttt{Wach\allowbreak Lattices\allowbreak Exist\allowbreak All\allowbreak Layers\allowbreak Of\allowbreak Bhatt\allowbreak Scholze}.
The bounded-prism definitions and actual Fontaine object are
read on their original rings; the completed source theorem is
used as a contract for reconstruction, not as a substitute for
the human arguments above.

\endgroup
\end{document}
